
%
%
\section{Introduction}

The notion of a \textbf{pseudorandom generator (PRG)}~\cite{BlumMi84,Yao82,Nisan91,NisanWi94} is ubiquitous in theoretical computer science, with vast applicability in cryptography and derandomization.  (See the texts \cite{Goldreich10-prgprimer,vadhan2012pseudorandomness} for more background on pseudorandomness.)
A recent %
paper by %
Braverman, Cohen, and Garg~\cite{BravermanCoGa18} introduced the following intriguing generalization of a PRG, in which we attach real coefficients to the outputs of the generator. %

\begin{definition}\label{def:booleanWPRG}
	Let $\mathcal{B}$ be a class of boolean functions $B \colon \zo^n \to \{0, 1\}$. An \textbf{$\eps$-weighted pseudorandom generator (WPRG)} for $\mathcal{B}$ is a function $(G,\rho) \colon \zo^s \to \zo^n\times \R$ such that for every $B \in \mathcal{B}$,
	$$
		\left |\E_{x \la U_n}[B(x)] - \E_{x \la U_{s}}[\rho(x)\cdot B(G(x))]\right | \leq \eps.
	$$
	The value $s$ is the \textbf{seed length} of the WPRG, and $n$ is the \textbf{output length} of the WPRG. 
	We say that the WPRG is \textbf{(mildly)\footnote{We consider this definition to correspond to \emph{mild} explicitness because requiring that
	the generator be computable in space linear in its seed length only implies that it is computable in time exponential in its seed length (\ie, time polynomial in the size of its truth table), which is
	mildly explicit according to the terminology in \cite{vadhan2012pseudorandomness}.  \emph{Strong} explicitness, in contrast,
	would require that each bit of the truth table %
be %
computable in 
	time polynomial in $s$.} explicit} if given $x$, $G(x)$ and $\rho(x)$ are computable in space $O(s)$, and $\rho(x)$ 
	has absolute value at most $2^{O(s)}$.\footnote{Naively the coordinates could be as large as $2^{2^{O(s)}}$, but this restriction is essentially WLOG for most models (\expref{Section}{app:intro}).}  %
\end{definition}
Above and throughout, we use the standard definition of space-bounded complexity, which counts the working, read-write memory of the algorithm, and does not include the length of the read-only input or write-only output, which can be exponentially longer than the space bound.

In the original %
article by %
Braverman, Cohen, and Garg~\cite{BravermanCoGa18} and previous versions of this paper~\cite{PyneVa20}, generators as above were called \textbf{pseudorandom pseudodistributions (PRPDs)}.  The terminology of weighted pseudorandom generators (WPRGs) was introduced
by Cohen et al.~\cite{CohenDoReSbTa21}, and we find it more intuitive
(and it avoids the double use of the `pseudo-' prefix).

With \expref{Definition}{def:booleanWPRG}, a PRG is a special case of a WPRG with $\rho(x)=1$ for all $x$.  The power of WPRGs comes from allowing the coefficients
to be negative, which yields cancellations.  Indeed, an explicit $\eps$-WPRG with seed length $s$ in which all of the coefficients are nonnegative can be converted into an explicit $O(\eps)$-PRG with seed length $s+O(\log(1/\eps))$ (see \expref{Section}{app:intro}). A WPRG for a class of functions that includes testing if the input equals an arbitrary fixed string must have coefficients bounded by $2^{O(s)}$ (see \expref{Section}{app:intro}).

A general WPRG can be converted into a linear combination of two unweighted generators. That is, for every explicit WPRG 
$(G,\rho) \colon \zo^s \to \zo^n\times \R$, there are 
explicit generators $G_+ : \zo^{s'}\rightarrow \zo^n$ and $G_- : \zo^{s'}\rightarrow \zo^n$ with seed length
$s'=O(s+\log(1/\eps))$ and (explicitly computable) coefficients $\rho_+,\rho_-\in \R^{\geq 0}$ such that for every function $B : \zo^n\rightarrow \zo$, we have %
$$\E_x[\rho(x)\cdot B(G(x))] = 
\rho_+\cdot \E_x\left[B(G_+(x))\right] - \rho_-\cdot \E_x\left[B(G_-(x))\right]\pm \eps.$$
(See \expref{Section}{app:intro}).

The motivation for WPRGs is that they can be used to derandomize algorithms in the same way as a PRG:
we can estimate the acceptance probability of any function $B\in \cB$ by enumerating over the seeds $x$ of the WPRG $(G,\rho)$ and calculating the average of the values $\rho(x)\cdot B(G(x))$.
Furthermore, ~\cite{BravermanCoGa18} observe that if $(G,\rho)$ is an $\eps$-WPRG then $G$ is an $\eps$-\textbf{hitting set generator (HSG)} for the same class.  That is, if $B$ is any function in $\cB$ with $\Pr[B(U_n)=1]>\eps$, then there exists an $x\in \zo^s$ such that $B(G(x))=1$.

Given this motivation, it is natural to ask whether WPRGs are more powerful than PRGs.  That is, can WPRGs achieve a shorter seed length than PRGs for a natural computational model $\cB$?  (There are simple constructions of artificial examples, one of which we give in \expref{Section}{app:intro}.) 
As discussed below, Braverman, Cohen, and Garg~\cite{BravermanCoGa18} gave an explicit construction of WPRGs achieving a shorter seed length than the \emph{best known} construction of PRGs for ordered branching programs, but not beating the best possible seed length for that model (given by a non-explicit application of the Probabilistic Method). In this %
paper,  %
we give an explicit construction of WPRGs for a natural computational model (ordered permutation branching programs of unbounded width) with a
seed length that beats all possible PRGs for that model. 


\subsection{Ordered branching programs}

    The work of Braverman, Cohen, and Garg~\cite{BravermanCoGa18}, as well as our paper, focuses
    on WPRGs for classes $\cB$ of functions computable by \emph{ordered branching programs}, a nonuniform model that captures how a space-bounded randomized algorithm accesses its random bits.
	
	\begin{definition} \label{def:obp}
	An \textbf{(ordered) branching program} $B$ of length $n$ and width $w$ computes a function $B:\zo^n \ra \{0,1\}$. On an input $\sigma\in\{0,1\}^n$, the branching program computes as follows. It starts at a fixed start state $v_0\in [w]$. Then for $r=1,\ldots,n$, it reads the next symbol $\sigma_r$ and updates its state according to a transition function $B_r:[w]\times \zo\ra[w]$ by taking $v_{t}=B_r(v_{t-1},\sigma_t)$. Note that the transition function $B_r$ can differ at each time step. 
	
	The branching program \textbf{accepts} $\sigma$, denoted $B(\sigma)=1$, if $v_{n}\in \acc$, where $\acc \subseteq [w]$ is the set of accept states, and otherwise it \textbf{rejects}, denoted $B(\sigma)=0$.
	Thus an ordered branching program is specified by the transition functions $B_1,\ldots,B_n$, the start state $v_0$ and the set $\acc$ of accept states.
	\end{definition}
	An ordered branching program is necessarily read-once, where the read order is fixed in advance.
	An ordered branching program of length $n$ and width $O(w)$ can compute the output of an algorithm that uses $\log w$ bits of memory and $n$ random bits, by taking the state at each layer as the contents of memory and configuration of the Turing machine at that time.  We note that we can convert any ordered branching program into one with a single accept state by collapsing all of $\acc$ to a single state.
	
	Using the Probabilistic Method, it can be shown that there \emph{exists} an $\eps$-PRG for ordered branching programs of length $n$ and width $w$ with seed length
	$s=O(\log(nw/\eps))$.  The classic construction of Nisan~\cite{Nisan92} gives an explicit
	PRG with seed length $s = O(\log n \cdot \log(nw/\eps))$, and this bound has not been improved except for extreme ranges of $w$, namely when
	$w$ is at least quasipolynomially larger than $n/\eps$~\cite{NisanZu96,Armoni98,KaneNeWo08} or when $w\leq 3$ \cite{BogdanovDvVeYe09,SimaZa11,GopalanMeReTrVa12,MekaReTa18}.
	Braverman, Cohen, and Garg~\cite{BravermanCoGa18} gave an explicit construction of a WPRG that achieves improved dependence on the error parameter $\eps$, with seed length  $$s=\tO\left(\log n\cdot \log(nw) + \log(1/\eps)\right).$$
	In particular, for error $\eps=n^{-\log n}$ and width $w=\poly(n)$, their seed length improves Nisan's from $O(\log^3 n)$ to $\tO(\log^2 n)$.  
    %
	Chattopadhyay %
    and Liao~\cite{ChattopadhyayLi20} gave a simpler construction of WPRGs with a slightly shorter seed length than \cite{BravermanCoGa18}, with an additive dependence on $O(\log(1/\eps))$ rather than $\tO(\log(1/\eps))$.

	
\subsection{Permutation branching programs}\label{subsec:PBP_overview}

Due to the lack of progress in constructing improved PRGs for general ordered branching programs as well as some applications~\cite{Indyk06,KaplanNaRe05}, attention has turned to more restricted classes of ordered branching programs.  In this %
article, %
our focus is on \emph{permutation} branching programs. %
	
	\begin{definition}\label{def:prbp}An \textbf{(ordered) permutation branching program} is an ordered branching program $B$ where for all $t \in [n]$ and $\sigma \in \zo$, $B_t(\cdot,\sigma)$ is a permutation on $[w]$. 
	\end{definition}
	This can be thought of as the computation being time-reversible on any fixed input $\sigma$. We note that we cannot assume without loss of generality that a permutation branching program has a single accept state, as merging a set of accept states will destroy the permutation property. Nevertheless, ordered permutation branching programs with a single accept state can compute interesting functions, such as testing
	whether a $\sum_{i\in S} x_i \equiv 0 \pmod m$, for any $m\leq w$ and any $S\subseteq [n]$.  An ordered permutation branching program with a single accept state can also compute permutation test functions, which test whether $x|_T = \pi(x|_S)$ for any permutation $\pi : \zo^\ell\rightarrow \zo^\ell$ and any two subsets $S,T\subseteq [n]$ of size $\ell$ such that all elements of $T$ are larger than all elements of $S$, provided that $w\geq 2^{\ell}$~\cite{HozaPyVa21}.
	
	Previous %
work %
on various types of PRGs for permutation branching programs~\cite{RozenmanVa05,ReingoldTrVa06,BravermanRaRaYe10,KouckyNiPu11,De11,Steinke12,HozaPyVa21} %
has  %
achieved seed lengths that are logarithmic or nearly logarithmic in the length $n$ of the branching program, improving the $\log^2 n$ bound in Nisan's generator. In particular, Braverman, Rao,
Raz, and Yehudayoff~\cite{BravermanRaRaYe10} gave a PRG for the more general model of \emph{regular} branching programs (with an arbitrary number of accept states) with seed length
	$$s=O\left(\log n\cdot (\log w + \log(1/\eps)+\log\log n)\right).$$
For getting an %
    HSG, they also showed how to %
    eliminate
	the $\log\log n$ and $\log(1/\eps)$ terms at the price of a worse dependence on $w$,\footnote{The lack of dependence on $\eps$ can be explained by the observation of Braverman et al. that any regular branching program that 
	has nonzero acceptance probability has acceptance probability at least
	$1/2^{w-1}$, so WLOG $\eps> 1/2^w$, \ie, $w> \log(1/\eps)$.} 
	achieving a seed length of
	$$s\leq \log(n+1) \cdot w.$$

	For the specific case of permutation branching programs, Kouck{\'y}, Nimbhorkar, and Pudl{\'a}k~\cite{KouckyNiPu11}, De~\cite{De11}, and Steinke~\cite{Steinke12}
	showed how to remove the $\log\log n$ term in the Braverman et al. PRG at the price of a worse dependence on $w$, achieving seed length
	$$s=O(\log n\cdot(\poly(w)+\log(1/\eps))).$$
	Most recently, Hoza, Pyne, and Vadhan~\cite{HozaPyVa21} showed that the dependence on the width $w$ could be entirely eliminated if we restrict to permutation branching programs with a \emph{single accept state}, constructing a PRG with seed length
	$$s=O(\log n\cdot (\log\log n+\log(1/\eps)).$$
	In particular, they show that this	seed length is provably better than what is achieved by the Probabilistic Method; that is, a random function with seed length $o(n)$ fails to be a PRG for unbounded-width permutation branching programs with high probability. To do so, they use the family of permutation test programs discussed above.
	Like the prior PRGs for bounded-width permutation branching programs, the seed length has a term of $O(\log n\cdot \log(1/\eps))$.  However, in contrast to the bounded-width case, this cannot be improved to $O(\log(n/\eps))$ by a non-explicit construction.  Hoza et al. prove that seed length $\Omega(\log n \cdot \log(1/\eps))$ is \emph{necessary} for any $\eps$-PRG against unbounded-width permutation branching programs.  For hitting-set generators (HSGs), they show that seed length $O(\log(n/\eps))$ is possible via the Probabilistic Method, thus leaving an explicit construction as an open problem. 
	
\subsection{Our results}
In this paper, we construct an explicit WPRG for permutation branching programs of unbounded width and a single accept state that beats the aforementioned lower bounds for PRGs. %

\begin{theorem}
\label{thm:cycle}
For all $n\in \N$ and $\eps \in (0,1/2)$, there is an explicit $\eps$-WPRG (and hence $\eps$-HSG) for ordered permutation branching 
programs of length $n$, arbitrary width, and a single accept state, with seed length
$$s=O\left(\log(n)\sqrt{\log(n/\eps)}\sqrt{\log\log(n/\eps)}+\log(1/\eps)\log\log(n/\eps)\right).$$
\end{theorem}
In particular, when $\eps=1/\poly(n)$, we achieve seed length $\tO(\log^{3/2} n)$, while
a PRG requires seed length $\Omega(\log^2 n)$~\cite{HozaPyVa21}.

As noted in \cite{HozaPyVa21}, an $\eps$-WPRG for branching programs with a single accept state is also an $(a\cdot \eps)$-WPRG for branching programs with at most $a$ accept states.  For bounded-width permutation branching programs, we can take $a=w$ and obtain the following. %
\begin{corollary} \label{cor:permutation-polywidth}
For all $n,w\in \N$ and $\eps \in (0,1/2)$, there is an explicit $\eps$-WPRG (and hence $\eps$-HSG) for ordered permutation branching 
programs of length $n$ and width $w$ (and any number of accept states), with seed length
$$s=O\left(\log(n)\sqrt{\log(nw/\eps)}\sqrt{\log\log(nw/\eps)}+\log(w/\eps)\log\log(nw/\eps)\right).$$
\end{corollary}
In particular for $w=\poly(n)$ and $\eps=1/\poly(n)$, we achieve seed length $\tO(\log^{3/2} n)$.  Note that the previous explicit PRGs (or even HSGs) for permutation branching programs (as mentioned in \expref{Subsection}{subsec:PBP_overview}) achieved seed length $o(\log^2 n)$ only when both $w=n^{o(1)}$ and $\eps=1/n^{o(1)}$.  With seed length $o(\log^2 n)$, 
\expref{Corollary}{cor:permutation-polywidth} can handle width as large as $w=n^{\tilde{\Omega}(\log n)}$ and error as small as $\eps=1/n^{-\tilde{\Omega}(\log(n))}$. 
We summarize these results in a table.\\

%
%
%
\hspace{-0.6cm}
\begin{tabular}{c|c|c|c}    %
     Citation & Type & Model & Seed Length \\\hline
     Non-explicit (folklore) & PRG & General & $\Theta(\log(nw/\eps)$\\
     \cite{Nisan92,ImpagliazzoNiWi94} & PRG & General & $O(\log n \cdot \log(nw/\eps))$\\
     \cite{BravermanRaRaYe10} & PRG & Regular & $\tO(\log n \cdot \log(w/\eps))$\\
     \cite{BravermanRaRaYe10} & HSG & Regular & $\log(n+1)\cdot w$\\
     \cite{KouckyNiPu11,De11,Steinke12} & PRG & Permutation & $O(\log n\cdot (\poly(w)+\log(1/\eps))$\\
     \cite{BravermanCoGa18,ChattopadhyayLi20}, \expref{Thm.}{thm:path} & WPRG & General & $\tO(\log n\cdot \log nw+\log(1/\eps))$\\
%
     \cite{HozaPyVa21} & PRG & Permutation ($1$ accept) & $\widetilde{\Theta}(\log n\cdot \log(1/\eps))$\\
     Non-explicit~\cite{HozaPyVa21} & HSG & Permutation ($1$ accept) & $O(\log (n/\eps))$\\
     \expref{Theorem}{thm:cycle} & WPRG & Permutation ($1$ accept) & $\tO(\log n\sqrt{\log(n/\eps)}+\log(1/\eps))$\\
%
     \expref{Corollary}{cor:permutation-polywidth} & WPRG & Permutation & $\tO(\log n\sqrt{\log(nw/\eps)}+\log(w/\eps))$\\
\end{tabular}

\subsection{Subsequent work}
Subsequent to the dissemination of the preprint version of this %
paper %
%
%
%
%
%
%
as \href{https://eccc.weizmann.ac.il/report/2021/019/}{ECCC TR21-019},
there have been several further developments. Independently and concurrently, Cohen, Doron, Renard, Sberlo, and Ta-Shma~\cite{CohenDoReSbTa21} obtained \expref{Theorem}{thm:path}, and both papers appeared at the 2021 Computational Complexity Conference. 
%
Subsequent to both preprints appearing, Hoza~\cite{Hoza21} combined \expref{Theorem}{thm:path} with the Armoni PRG~\cite{Armoni98} and the framework of Saks and Zhou~\cite{SaksZh99} to slightly improve the derandomization of $\mathbf{BPL}$ to %
$\mathbf{BPSPACE}(S)\subset \mathbf{DSPACE}(S^{3/2}/\sqrt{\log S})$. %
He also constructed a weighted PRG for ordered branching programs with seed length $O(\log(n)\log(nw)+\log(1/\eps))$, obtaining a weighted PRG with matching dependence on $n$ and $w$ to~\cite{Nisan92} with optimal dependence on $\eps$. 
%
%
Subsequently, 
%
Bogdanov, Hoza, Prakriya, and Pyne~\cite{BHPP22} obtained a hitting set for regular branching programs with seed length $O(\log^{3/2}n+\log(n)\log(w)+\log(n)\sqrt{\log(1/\eps)})$, improving on prior work~\cite{BravermanRaRaYe10} in the case that $\eps=1/\poly(n)$ and $w=n^{o(1)}$. %
%
Furthermore, Chen, Hoza, Lyu, Tal, and Wu~\cite{CHLTW23} gave a new proof of \expref{Theorem}{thm:cycle}. Their new proof could be extended to regular branching programs, allowing them to construct a weighted PRG for regular branching programs with seed length matching that of the hitting set of~\cite{BHPP22}. Finally, Chattopadhyay and Liao~\cite{CL24} substantially simplified the proof of~\cite{CHLTW23} in the case of permutation programs. All of these works use the core idea of analyzing an error reduction procedure as a weighted PRG.

\subsection{Versions of this paper}
The original ECCC preprint version of this paper contained the main result of \expref{Theorem}{thm:cycle}. The CCC publication was an extended abstract that did not include detailed proofs of the main claims. The present version contains proofs of the main claims, and a new section (\expref{Section}{app:intro}) with some ancillary observations about weighted PRGs (this section also appeared in an updated version of the paper posted to ECCC).
%

\section{Overview of proofs}
The starting point for our %
%
work was %
the recent
family of %
%
%
space-efficient algorithms for estimating
random-walk probabilities in directed graphs by
%
Ahmadinejad, %
Kelner, Murtagh,
Peebles, Sidford, and Vadhan~\cite{AhmadinejadKeMuPeSiVa19}, %
%
based on spectral graph
theory and space-efficient Laplacian solvers. 
    We interpret these algorithms as giving WPRGs with large seed length, which we then derandomize to obtain our results.  
    
    The specific problem considered by %
Ahmadinejad  %
et al. is the following: given a directed graph $\cG=(V,E)$, two vertices $s,t\in V$, a walk-length $k\in \N$, and an error parameter $\eps>0$, estimate the probability that a random walk of length $k$ started at $s$ ends at $t$ to within $\pm \eps$.  Such an algorithm can be applied to the following graph in order to estimate the acceptance probability of an ordered branching program:
    \begin{definition} \label{def:OBP_graph}
%
Given a length-$n$, width-$w$    %
branching program $B$ with transition functions $(B_1,\ldots,B_n)$ with start vertex $v_0\in [w]$, and a single accept vertex $\va$, the \textbf{(layered) graph associated with $B$} is the graph $\cB$ with vertex set $\{0,1,\ldots,n\}\times[w]$ and directed edges from $(i-1,v)$ to
        $(i,B_i(v,0))$ and $(i,B_i(v,1))$ for every $i=1,\ldots,n$ and $v\in [w]$.
    \end{definition} 
    Applying the algorithms of %
Ahmadinejad  %
et al. to the graph $\cG$ with $s=(0,v_0)$, $t=(n,\va)$, and $k=n$, we obtain an estimate of the acceptance probability of $B$ to within $\pm \eps$, just like an $\eps$-WPRG for $B$ would allow us to obtain. But a WPRG $(G,\rho)$ is much more constrained than an arbitrary space-efficient algorithm, which can directly inspect the graph. Instead, a WPRG is limited to generating $S=2^s$ walks of length $n$ in the layered graph, described by sequences $G(x_1),\ldots, G(x_S)\in \zo^n$ of edge labels, and then combining the indicators $B(G(x_1)),\ldots,B(G(x_S))$ of whether the walks ended at $t$ via a linear combination with fixed coefficients $\rho(x_1),\ldots,\rho(x_S)\in \R$.
    
    Note that if $B$ is a permutation branching program, then the graph $\cG$ above is $2$-regular (in that every vertex has $2$ in- and out- edges, except for layer $0$ which has no incoming edges and layer $n$ which has no outgoing edges).  Thus, the basis for \expref{Theorem}{thm:cycle} is the (main) result 
    of %
Ahmadinejad  %
et al., which applies to regular (or more generally, Eulerian) directed graphs $G$.
    However, they also give a new algorithm for estimating random-walk probabilities in arbitrary directed graphs.  This
    algorithm is not as space-efficient as the ones for regular graphs, but is significantly simpler, so we begin by describing how to
    obtain a WPRG based on that algorithm.  The resulting WPRG matches the parameters of the WPRG of Braverman, Cohen, and Garg~\cite{BravermanCoGa18}, but has a significantly simpler proof (and is also simpler than the construction of %
     Chattopadhyay %
     and Liao~\cite{ChattopadhyayLi20}). A similar construction was independently discovered by Cohen, Doron, Renard, Sberlo, and Ta-Shma~\cite{CohenDoReSbTa21}.
    
    \subsection{WPRG for arbitrary ordered branching programs}\label{subsec:OBP_overview}
    Let $B$ be an arbitrary width-$w$, length-$n$ ordered branching program, with associated layered graph $\cG$ as in \expref{Definition}{def:OBP_graph}.  The algorithm of
%
Ahmadinejad  %
et al. starts with the $(n+1)w\times (n+1)w$
%
random-walk transition matrix $\mW$ of $\cG$, which has the following block structure %

    $$\mW=
    \begin{bmatrix}
    0 & \mB_1 & 0 & \cdots & 0 \\
    0 & 0 & \mB_2 & \cdots & 0 \\
    \vdots & & \ddots & & \vdots\\
    0 & 0 & 0 & \ddots & \mB_n\\
    0 & 0 & 0 & \cdots &  0 \\
    \end{bmatrix}.$$ %
    Here entry $((i,u),(j,v))$ is the probability that taking one random step in $\cG$ from vertex $(i,u)$ ends at $(j,v)$.  Thus $\bp_i$ is the $w\times w$ transition matrix for the random walk from layer $i-1$ to $i$ in the branching program.
   (Note that the matrix $\mW$ is not quite stochastic due to layer $n$ having no
   outgoing edges.)
   
%
Ahmadinejad  %
et al. consider the Laplacian $\mL = \I_{(n+1)w} - \mW$.  Its
   inverse $\mL^{-1} = (\I_{(n+1)w}-\mW)^{-1} = \I_{(n+1)w}+\mW+\mW^2+\mW^3+\cdots$ sums up random-walks of all lengths in $G$, and thus has the following form. %
   $$\mL^{-1} =
   \begin{bmatrix}
    \mB_{0\ldots0} & \mB_{0\ldots 1} & \mB_{0\ldots 2} & \cdots & \mB_{0\ldots n} \\
    0 & \mB_{1\ldots 1} & \mB_{1\ldots 2} & \cdots & \mB_{1\ldots n} \\
    \vdots & & \ddots & & \vdots\\
    0 & 0 & 0 & \ddots & \mB_{n-1\ldots n}\\
    0 & 0 & 0 & \cdots &  \mB_{n\ldots n} \\
    \end{bmatrix},$$
   where $\mB_{i,i}=\I$ and 
   $$\mB_{i .. j} = \mB_{i+1}\mB_{i+2}\cdots \mB_j.$$
   
   In particular, the $(0,n)$'th block of $\mL^{-1}$ gives the random-walk
   probabilities from layer $0$ to layer $n$, and thus the acceptance probability
   of $G$ is exactly the $(v_0,\vac)$'th entry of the $(0,n)$'th block of $\mL^{-1}$.  Therefore, the task reduces to producing a sufficiently good estimate of $\mL^{-1}$.
   
%
Ahmadinejad  %
et al. estimate $\mL^{-1}$  in two steps.  First, they observe that the Saks--Zhou derandomization of logspace~\cite{SaksZh99} can be used to produce,
   in deterministic space $O(\log(nw)\sqrt{\log(n)})$, %
   %
   approximations %
   $\widetilde{\mB_{i .. j}}$ of the blocks $\mB_{i .. j}$ to within entrywise error $1/\poly(nw)$, resulting in an approximate pseudoinverse
   \begin{equation}
       \Lt =
   \begin{bmatrix}
    \widetilde{\mB_{0\ldots0}} & \widetilde{\mB_{0\ldots 1}} & \widetilde{\mB_{0\ldots 2}} & \cdots & \widetilde{\mB_{0\ldots n}} \\
    0 & \widetilde{\mB_{1\ldots 1}} & \widetilde{\mB_{1\ldots 2}} & \cdots & \widetilde{\mB_{1\ldots n}} \\
    \vdots & & \ddots & & \vdots\\
    0 & 0 & 0 & \ddots & \widetilde{\mB_{n-1\ldots n}}\\
    0 & 0 & 0 & \cdots &  \widetilde{\mB_{n\ldots n}} \\
    \end{bmatrix},\label{eqn:path-inverse}
   \end{equation}
   with the property that
   $$\lbar \I_{(n+1)w}-\Lt \mL\rbar_{1}\leq 1/nw,$$
   where $\|\cdot\|_1$ denotes the $\ell_1$ operator norm on row vectors, \ie, $\|\mathbf{M}\|_1 =\sup_{x\neq 0}\|x\mathbf{M}\|_{1}/\|x\|_{1}$.

   Next, %
Ahmadinejad  %
et al. reduce the approximation error to an arbitrary $\eps<1/(nw)^{O(1)}$ by using preconditioned Richardson iterations, as captured by the following lemma. %
   
   
    \begin{restatable}[preconditioned Richardson iteration, \cite{AhmadinejadKeMuPeSiVa19} Lemma 6.2]{lemma}{Richardson}
    \label{lem:richardson}
    Let $\|\cdot\|$ be a submultiplicative norm on $N\times N$ real matrices. Given matrices $\mA,\mP_0\in \R^{N\times N}$ such that $\|\I_N-\mP_0\mA\|\leq \alpha$ for some $\alpha >0$, let $\Pm=\sum_{i=0}^m(\I_N-\mP_0\mA)^i\mP_0$. Then $\|\I_N-\Pm \mA\|\leq \alpha^{m+1}$.
    \end{restatable}
    
    Setting $N=(n+1)w$, $\mA=\mL$, $\mP_0=\Lt$, and $\alpha=1/n$, and
    $m=O(\log_{n}(1/\eps))$, we obtain $\Lte=\Pm$ such that
    $\|\I_N-\Lte \mL\|_1\leq \eps/(nw)^{O(1)}$, which implies that $\Lte$ and $\mL^{-1}$ are entrywise equal up to $\pm \eps$, for
    \begin{equation}
        \Lte = \sum_{i=0}^m(\I_N-\Lt\mL)^i\Lt \label{eqn:iterated_L}
    \end{equation}
    In particular, the $(v_0,\vac)$'th entry of the $(0,n)$'th block of $\Lte$ is an estimate of the acceptance probability of the branching program to within $\pm \eps$.  Computing $\Lte$ from $\mL$ and $\Lt$ 
    can be done in space $O((\log nw)\cdot \log m),$ yielding 
%
Ahmadinejad  %
et al.'s space bound of 
    $$O(\log(nw)\sqrt{\log(n)} + (\log nw)\cdot \log \log_{n}(1/\eps)).$$
    
Now we show how, with an appropriate modification, we can interpret this algorithm of
%
Ahmadinejad  %
et al. as a WPRG (albeit with large seed length).
    We replace the use of the Saks--Zhou algorithm (which requires looking at the branching program) with Nisan's pseudorandom generator.  Specifically, we take 
    $\widetilde{\mB_{i .. j}}$ to be the matrix whose $(u,v)$'th entry
    is the probability that, if we start at state $u$ in %
    the $i$'th layer and
    use a random output of Nisan's pseudorandom generator to take $j-i$ steps in the branching program, we end at state $v$ in the $j$'th layer.  For 
    $\widetilde{\mB_{i .. j}}$ to approximate $\mB_{i .. j}$ to within error $\pm 1/\poly(nw)$ as above, Nisan's pseudorandom generator requires seed length 
    
    $$\sNisan = O(\log(j-i)\cdot \log nw) = O(\log n\cdot \log nw).$$
    Observe that for every $i$, $\widetilde{\mB_{i\ldots i}} = \I_w = \mB_{i\ldots i}$. Without loss of generality, we may also assume that
    $\widetilde{\mB_{(i-1)\ldots i}}=\mB_{(i-1)\ldots i}$, since taking one step only requires one random bit.  
    
    Next, we observe from \expeqref{Equation}{eqn:iterated_L} that the matrix $\Lte$ is a polynomial of degree $2m+1$ in the matrices $\mL$ and $\Lt$.  In particular the $(0,n)$'th block of $\Lte$ is a polynomial of degree at most $2m+1$ in the matrices $\widetilde{\mB_{i .. j}}$.  Specifically, using the upper-triangular structure of the matrices $\mL$ and $\Lt$ and noting that the product of $d$ $(n+1)\times (n+1)$ block matrices expands into a sum of $(n+1)^{d-1}$ terms, each of which is a product of $d$ individual blocks, we show the following. %
    
    \begin{observation} \label{obs:path}
    The $(0,n)$'th block of $\Lte$ equals the sum of at most $(n+1)^{O(m)}$ terms, each of which is of the form 
    \begin{equation} \label{eqn:terms}
    \pm \widetilde{\mB_{i_0\cdots i_1}}
    \widetilde{\mB_{i_1\cdots i_2}}
    \cdots
    \widetilde{\mB_{i_{r-1}\cdots i_r}},
    \end{equation}
    where $0=i_0<i_1<i_2<\cdots<i_r=n$ and $r\leq 2m+1$.
    \end{observation}
    
    Notice that, up to the sign, each term as expressed in \expeqref{Equation}{eqn:terms} is the transition matrix for a pseudorandom walk from layer 0 to layer $n$ of the branching program, where we use $r\leq m+1$ independent draws from Nisan's generator, with the $j$'th draw being used to walk from layer $i_{j-1}$ to layer $i_j$.  In particular,
    the $(v_0,\vac)$ entry of \expeqref{Equation}{eqn:terms} equals the acceptance probability of the branching program on such a pseudorandom walk (up to the $\pm$ sign).
    Thus the algorithm now has the form required of a WPRG. 
    
    The seed length
    for the WPRG is the sum of the seed length $\ssum$ needed to select a random term in the sum (using the coefficients of the WPRG to rescale the sum into an expectation) and the seed length $\sterm$ to generate a walk for the individual term.
    To select a random term in the sum requires a seed of length
    $$\ssum = \log n^{O(m)}=O(m\cdot \log(n))=O(\log_{n}(1/\eps)\cdot\log(n))=O(\log(1/\eps)).$$
    The seed length needed for an individual term is at most $$\sterm = O(m)\cdot \sNisan=O(\log_{n}(1/\eps)\cdot \log(n)\cdot\log nw) = O(\log(1/\eps)\cdot \log(n)).$$ 
    The latter offers no improvement over Nisan's PRG. (Recall that $\eps<1/nw$.) To obtain a shorter seed length, we just need to derandomize the product in \expeqref{Equation}{eqn:terms}.  Instead of using $r$ independent seeds, we use dependent seeds generated using the Impagliazzo--Nisan--Wigderson pseudorandom generator~\cite{ImpagliazzoNiWi94}.  Specifically, we can produce a pseudorandom walk that approximates the product to within entrywise error $\pm \gamma$ using a seed of length 
    $$\sterm' = \sNisan+O((\log r)\cdot \log(rw/\gamma)).$$
    The entrywise error of $\gamma$ in each term may accumulate over the $n^{O(m)}$ terms, so to achieve a WPRG error of $O(\eps)$, we should set $\gamma=\eps/n^{O(m)}=1/\eps^{O(1)}$.  Recalling that $r\leq 2m+1 = O(\log_{n}(1/\eps))$, we attain a seed length of
    \begin{eqnarray*}
    \ssum+\sterm' 
    &=& O(\log(1/\eps)) + O(\log n\cdot \log nw) + O(\log\log_{n}(1/\eps)\cdot \log(1/\eps))\\
    &=& O(\log n\cdot \log nw + \log(1/\eps)\cdot \log\log_{n}(1/\eps)),
    \end{eqnarray*}
    which slightly improves over the bound of Braverman, Cohen, and Garg~\cite{BravermanCoGa18}, and is incomparable to that of Chattopadhyay and Liao~\cite{ChattopadhyayLi20}. Specifically, our first term of $O(\log n\cdot \log nw)$ is better than~\cite{ChattopadhyayLi20} by a factor of $\log\log(nw)$, but our second term of $O(\log(1/\eps)\cdot \log\log_{n}(1/\eps))$ is worse by a factor of $\log\log_{n}(1/\eps)$.
    
    \subsection{WPRG for permutation branching programs}\label{subsec:PBP_proof_overview}
    Now we give an overview of our WPRG for permutation branching programs, as stated in \expref{Theorem}{thm:cycle}. This is based on the
    the algorithm of 
%
Ahmadinejad  %
et al. that estimates random-walk probabilities in \emph{regular} (or even Eulerian) digraphs with better space complexity than the algorithm described in \expref{Subsection}{subsec:OBP_overview}.  As before, we will review their algorithm as applied to the $((n+1)\cdot w)$-vertex graph $\cG$ associated with an ordered branching program $B$ of length $n$ and width $w$.  Since we assume that the branching program $B$ is a permutation program, the graph $\cG$ will be 2-regular at all layers other than $0$ and $n$.  For the spectral graph-theoretic machinery used by 
%
Ahmadinejad  %
et al., it is helpful to work with random-walk matrices that correspond to strongly connected digraphs, so we also add a complete bipartite graph of edges from layer $n$ back to layer $0$, resulting in the following modified version of the matrix $\mW$:
    
    \begin{equation}
        \mW_0 =
    \begin{bmatrix}
    0 & \mB_1 & 0 & \cdots & 0 \\
    0 & 0 & \mB_2 & \cdots & 0 \\
    \vdots & & \ddots & & \vdots\\
    0 & 0 & 0 & \ddots & \mB_n\\
    \Jw & 0 & 0 & \cdots &  0 \\
    \end{bmatrix}, \label{eqn:bWhat}
    \end{equation}
    where the $\Jw$ in the lower-left corner is the $w\times w$ matrix in which every entry is $1/w$ (corresponding to the complete bipartite graph we added).  Notice that the matrix $\Jw$ is identically zero when applied to any vector that is orthogonal to the uniform distribution, so it is not very different %
    from having $0$ in the lower-left block as we had before. Indeed, the powers of $\mW_0$ have the following form. %
    \begin{equation} \label{eqn:bWhatpowers} 
    \mW_0^2 =
    \begin{bmatrix}
    0 & 0 & \mB_{0..2} & 0 & 0 \\
    \vdots & 0 & 0 & \ddots & 0 \\
    0 & & \vdots & & \mB_{n-2..n}\\
    \Jw & 0 & 0 & \cdots & 0\\
    0& \Jw & 0 & \cdots &  0 \\
    \end{bmatrix}, \ldots,
    \mW_0^n =
    \begin{bmatrix}
    0 & 0 & \cdots & 0 & \mB_{0..n} \\
    \Jw & 0 &  &  & 0 \\
    0 & \ddots &  & & 0\\
    \vdots & 0 & \Jw & 0 & 0\\
    0& 0 & 0 & \Jw &  0 \\
    \end{bmatrix}
    \end{equation}
    where as before 
    $$\mB_{i..j} = \mB_{i+1}\mB_{i+2}\cdots \mB_j.$$
    Notice in particular that $\mW_0^{n+1}$ will be a block-diagonal
    matrix with %
    $\Jw$ %
on the diagonal (\ie, $\mW_0^{n+1}=\I_{n+1}\otimes \Jw$), and thus has no dependence on the branching program $B$.
    
    Now the Laplacian $\I_{(n+1)w}-\mW_0$ is no longer invertible (the uniform distribution is in the kernel).  In \cite{AhmadinejadKeMuPeSiVa19}, they instead estimate the Moore--Penrose pseudoinverse of $\I_{(n+1)w}-\mW_0$. We instead scale $\mW_0$ by a factor $c=1-1/(n+1)$, and consider the Laplacian $\mL_0=\I_{(n+1)w}-c\mW_0$. Looking ahead, this scaling factor ensures that the condition number of $\mL_0$ depends only on $n$, allowing us to obtain a seed length independent of $w$. Then, by the expressions above for the powers of $\mW_0$, it can be shown that from
    $$\mL^{-1}_0 = \I_{(n+1)w}+c\mW_0+c^2\mW_0^2+c^3\mW_0^3+\dots$$
    we can compute $\mB_{0..n}$, which appears in $\mW_0^n$ with a scaling factor $c^n\geq 1/4$.
   
   So again to estimate the acceptance probability of $B$, it suffices to compute a sufficiently good approximation to $\mL_0^{-1}$.
   As before, it suffices to compute a matrix
   $\Ltz$ such that $\|\I_N-\Ltz \mL_0\| \leq \alpha$ for some constant $\alpha<1$ (letting $N:=(n+1)w$) and a submultiplicative matrix norm $\|\cdot\|$, because then we can use preconditioned Richardson iterations (\expref{Lemma}{lem:richardson}) to estimate $\mL_0$ to within arbitrary entrywise accuracy.
   
   Unfortunately %
   %
   %
   %
%
%
%
%
   we don't know how to directly obtain such an initial approximation $\Ltz$ efficiently
enough for our result. Instead, following 
%
Ahmadinejad  %
et al., we tensor $\mW_0$ with a sufficiently long directed cycle.  Specifically, we let $\mC_i$ be the directed cycle on $2^i$ vertices, and consider $\mC_q$ for $q=\log(n+1)$ (which we assume is an integer WLOG). We consider the \emph{cycle lift}, whose transition matrix is
    $$\mC_q\otimes\mW_0 =
    \begin{bmatrix}
    0 & \mW_0 & 0 & \cdots & 0 \\
    0 & 0 & \mW_0 & \cdots & 0 \\
    \vdots & & \ddots & & \vdots\\
    0 & 0 & 0 & \ddots & \mW_0\\
    \mW_0 & 0 & 0 & \cdots &  0 \\
    \end{bmatrix}.$$  %
    Then, we seek to invert the Laplacian $\mL = \I_{2^qN} - c \mC_q \otimes \mW_0$.  Similarly
to the above, we have %

    \begin{eqnarray*}
    \mL^{-1} \!\!\! &=&\!\!\!\! \left(\I_{2^qN} - c \mC_q\otimes \mW_0\right)^{-1}\\
    &=& \!\!\! \left(\I_{2^qN} - c^{n+1} \mC_q^{n+1} \otimes \mW_0^{n+1}\right)^{-1}\cdot
    \left(\I_{2^qN} + c \mC_q\otimes \mW_0 +c^2 \mC^2_q\otimes \mW_0^2 + \cdots +
    c^n \mC^n_q\otimes \mW_0^n\right)\\
    &=& \!\!\! \left(\I_{2^qN} - c^{n+1} \mC_q^{n+1}\otimes(\I_{n+1} \otimes \Jw)\right)^{-1} \cdot 
        \left(\I_{2^qN} + c \mC_q\otimes \mW_0 +c^2 \mC^2_q\otimes \mW_0^2 + \cdots +
    c^n \mC^n_q\otimes \mW_0^n\right).  %
        %
    \end{eqnarray*}
%
    Thus, letting $$\mM = \I_{2^qN} - c^{n+1} \mC_q^{n+1}\otimes(\I_{n+1} \otimes \Jw)= \I_{2^qN} - c^{n+1} \I_{2^q}\otimes(\I_{n+1} \otimes \Jw),$$
    which has no dependence on the branching program,
    we %
    %
    obtain %
    \begin{eqnarray*}
    \mM\cdot \mL^{-1}
    &=& \I_{2^qN} + c \mC_q\otimes \mW_0 +c^2 \mC^2_q\otimes \mW_0^2 + \cdots + 
    c^n \mC^n_q\otimes \mW_0^n\\
    &=&
    \begin{bmatrix}
    \I_N & c\mW_0 & c^2\mW_0^2 & \cdots & c^n\mW_0^n \\
    c^n\mW_0^n & \I_N & c\mW_0 & \cdots & c^{n-1}\mW_0^{n-1} \\
    \vdots & & \ddots & & \vdots\\
    c^2\mW_0^2 & c^3\mW_0^3 & c^4\mW_0^4 & \ddots & c\mW_0\\
    c\mW_0^1 & c^2\mW_0^2 & c^3\mW_0^3 & \cdots &  \I_N \\
    \end{bmatrix}.
    \end{eqnarray*}

    Thus, if we can accurately estimate $\mL^{-1}$, we can
    obtain an accurate estimate of  $\mW_0^n$, whose upper-right block equals $\mB_{0..n}$ and thus contains the acceptance probability of the 
    branching program.
    
    To compute an approximate inverse of $\mL=\I_{2^qN} - c \mC_q \otimes \mW_0$, 
%
Ahmadinejad  %
et al. provide a recursive formula
    expressing $(\I_{2^qN} - c \mC_q \otimes \mW_0)^{-1}$ in terms of
    $(\I_{2^{q-1}N} - c^2 \mC_{q-1} \otimes \mW^2_0)^{-1}$ and some applications of the matrix $\mW_0$.
    That is, computing the inverse of the Laplacian of the cycle lift of $\mW_0$ reduces
    to computing the inverse of the Laplacian of a cycle lift of $\mW_0^2$ with a cycle of half the length.  At the bottom of the recursion (after $q$ levels of recursion), we need to compute the inverse of $$\I_N-c^{2^q}\mW_0^{2^q} = \I_N - c^{n+1}\mW_0^{n+1} = \I_N - c^{n+1} \I_{n+1}\otimes \Jw,$$
    which is easy (and does not depend on the branching program).
    The resulting formula for $(\I_{2^qN} - c \mC_q \otimes \mW_0)^{-1}$
    is a polynomial in $\mW_0,\mW_0^2,\mW_0^4,\ldots,\mW_0^{2^{q-1}}$.
    However, computing these high powers of $\mW_0$ exactly is too expensive in space usage.  
    
    Thus, instead %
Ahmadinejad  %
et al. use the \emph{derandomized square}~\cite{RozenmanVa05} which allows for computing a sequence
    $\mW_0,\mW_1,\ldots,\mW_q$ where $\mW_i$ is a sparsification of $\mW^2_{i-1}$ with two desirable properties. First, $\mW_q$ can be constructed in
    deterministic space $$O(\log nw+q\cdot \log(1/\delta))$$ for an error parameter $\delta$, rather than the space $O(q\cdot \log nw)$ of exact repeated squaring.  Second, they introduce a new notion of spectral approximation called \emph{unit-circle approximation}, and show that
    the derandomized square
    $\mW_i$ is a unit-circle approximation of $\mW_{i-1}^2$ to within 
    error $\delta$.
    Using these repeated derandomized squares in the recursion, 
%
Ahmadinejad  %
et al. obtain an approximate inverse $\Lt$ with the
%
following properties.  %
    \begin{enumerate}
        \item The $N\times N$ blocks of $\mM\cdot \Lt$ are 
        each of the form $\mW_{i_1}\mW_{i_2}\cdots \mW_{i_r}$
        where $r=O(q)$ \label{itm:degree}
        \item There is a submultiplicative matrix norm $\|\cdot \|_\mF$ such that
        $\|\I_{2^qN}-\Lt \mL\|_\mF = O(q^2\delta)$. Moreover, achieving an $\eps/\poly(n)$ approximation of $\mM\cdot \mL^{-1}$ in $\mF$-norm implies an $\eps$ approximation of $\mM\cdot \mL^{-1}$ in $\max$-norm.   %
Ahmadinejad  %
et al. actually lose a factor of $\poly(nw)$ in moving from $\mF$-norm to approximation in $\max$-norm, but we improve this bound to $\poly(n)$ by our choice of scaling factor $c=1-1/(n+1)$.   
        \label{itm:Fnorm} 
    \end{enumerate}

    
    \expref{Item}{itm:degree} allows for constructing $\mM\cdot \Lt$ from
    $\mW_0,\mW_1,\ldots,\mW_q$
    in space
    $$O(\log q \cdot \log nw).$$
    By \expref{Item}{itm:Fnorm}, if we take $\delta<1/\Omega(q^2)$, we 
    can apply 
    preconditioned Richardson iterations (\expref{Lemma}{lem:richardson})
    with degree $m=O(\log(n/\eps)/\log(1/q\delta))$ to obtain $\Lte=\Pm$ 
    such that $\mM\cdot \Lte$ approximates $\mM\mL^{-1}$ to within entrywise error $\eps.$
    The preconditioned Richardson iterations have an additive
    space cost of %
    $$O(\log m\cdot \log nw).$$
    Taking $\delta=1/\Omega(q^2)$ and recalling that $q=\log(n+1)$, the final space complexity is
    $$O(\log(nw)+q\log q) + O(\log q\cdot \log nw)+O(\log\log (n/\eps)\cdot \log nw) = O(\log nw \cdot \log\log(n/\eps)).$$
    
    To view this algorithm as a WPRG for permutation branching programs, we use
    the equivalence between the Impagliazzo--Nisan--Wigderson (INW) generator on permutation branching programs and the derandomized square of the corresponding graph, as established in \cite{RozenmanVa05,HozaPyVa21}.
    Using this correspondence, the matrix $\mW_i$ has the same structure as $\mW^{2^i}$ (see \expeqref{Equation}{eqn:bWhatpowers}), except that each block of the form $\mB_{j..j+2^i}$ is replaced with a matrix $\widetilde{\mB_{j..j+2^i}}$ that is the transition matrix of a pseudorandom
    walk from layer $j$ of the branching program to layer $j+2^i$ using the INW generator.  
    The seed length to generate this pseudorandom walk is
    $$\sINW = O(q\log(q/\delta)),$$
    which, as highlighted in \cite{HozaPyVa21}, is independent of the width $w$ of the branching program.
    This is the place where we use the fact that $B$ is a permutation branching program rather than a regular branching program.  Even though the algorithm of Ahmadinejad et al. works for regular directed graphs (and hence regular branching programs), the derandomized square operations used in that case can no longer be viewed as being obtained by using a pseudorandom generator to derandomize walks in the graph.
    
    Then, again assuming without loss of generality that 
    $\widetilde{\mB_{(j-1)\ldots j}}=\mB_{(j-1)\ldots j}$ for $j=1,\ldots,n$, we have
    the following analogue of \expref{Observation}{obs:path}. %
    \begin{observation} \label{obs:cycle}
    The upper-right $w\times w$ block of $\mM \cdot \Lte$ equals the sum of at most $n^{O(m)}$ terms, each of which is of the form 
    \begin{equation} \label{eqn:permutation-terms}
    \pm \widetilde{\mB_{i_0\cdots i_1}}
    \widetilde{\mB_{i_1\cdots i_2}}
    \cdots
    \widetilde{\mB_{i_{r-1}\cdots i_r}},
    \end{equation}
    where $0=i_0<i_1<i_2<\cdots<i_r=n$ and $r=O(qm)$.
    \end{observation}
    
    As in \expref{Subsection}{subsec:OBP_overview}, the algorithm now has the form required of a WPRG and our only remaining challenge is to keep the seed length small. 
    The seed length
    for the WPRG is the sum of the seed length needed to select a random term in the sum (using the coefficients of the WPRG to rescale the sum into an expectation) and the seed length to generate a walk for the individual term.
    To select a random term in the sum requires a seed of length
    $$\ssum = \log (n^{O(m)}).$$
    The seed length needed for an individual term is at most $$\sterm = O(qm)\cdot \sINW,$$
    which again would be too expensive for us.
    To derandomize the product in \expeqref{Equation}{eqn:permutation-terms}, 
    we again use the INW generator, but
    rely on the analysis in \cite{HozaPyVa21} for permutation branching
    programs to maintain a seed length that is independent of the width.  Specifically, we can produce a pseudorandom walk that approximates the product to within entrywise error $\pm \gamma$ using a seed of length 
    $$\sterm' = \sINW+O((\log r)\cdot \log(\log(r)/\gamma))
    = \sINW + O(\log qm\cdot \log(\log(qm)/\gamma)).$$
    The entrywise error of $\gamma$ in each term may accumulate over the $n^{O(m)}$ terms, so to achieve a WPRG error of $O(\eps)$, we should set $\gamma=\eps/n^{O(m)}$, which means that $\sterm' \geq \ssum$.  
    
    All in all, we attain a seed length of
    \begin{eqnarray*}
    \ssum+\sterm' 
    &=& O(m\log n)+\sINW+O((\log qm)\cdot \log(\log(qm)/\gamma))\\
    &=& O(q\log(q/\delta)) + \tO(m\log n) + O(\log qm\cdot \log(n/\eps))\\
    &=& \tO\left(\log n\cdot \log(1/\delta) + \frac{\log(n/\eps)}{\log(1/(\delta\log n))}\cdot \log n + \log\log(n/\eps)\cdot \log(n/\eps)\right)\,.
     %
    \end{eqnarray*}
    Optimizing the choice of $\delta$ as $\delta = \exp(-\widetilde{\Theta}(\sqrt{\log(n/\eps)}))$, we get a seed length of
    $$\tO(\log n\sqrt{\log(n/\eps)} + \log(1/\eps)).$$
    Note that the choice of $\delta$ here is much smaller than in 
    the %
Ahmadinejad  %
et al. algorithm, which used $\delta=1/\polylog(n)$.
    The reason we need the smaller choice of $\delta$ is to reduce the effect of the $\log(n^{O(m)})$ price we pay in $\ssum$ and $\sterm'$, which does not have an analogue in the algorithm of
%
Ahmadinejad  %
et al.
    
\subsection{Perspective}



Some intuition for the ability of WPRGs to beat the parameters of PRGs can come from
the study of \emph{samplers}~\cite{Goldreich97}.  A \emph{sampler} for a class $\cF$ of
functions $f : \zo^m\rightarrow \bbR$ is %
a 
%
randomized algorithm $\Samp$ that is given 
oracle access to a function $f\in \cF$ and, with probability at least $1-\delta$, outputs an estimate of $\E[f(U_n)]$ to within additive error $\pm \eps$. Most often, the class $\cF$ is taken to be the class of all bounded functions $f : \zo^m \rightarrow [0,1]$, but some
%
authors  %
have considered the general definition and other classes, such as the class $\cF$ of unbounded functions $f$ such that the random variable $f(U_n)$ has subgaussian tails~\cite{Blasiok18,Agrawal19}.  Two key complexity parameters of a sampler are its \emph{randomness complexity} (the number of coin tosses it uses, typically as a function of $m$, $\delta$, and $\eps$)  and its 
\emph{sample complexity} (the number of queries it makes to the oracle $f$).
An \emph{averaging sampler} is one that has a restricted form, where it uses its coin tosses to generate (possibly correlated) samples $x_1,\ldots,x_S$, and then outputs the average of $f$
on the samples, \ie, $(f(x_1)+\cdots+f(x_S))/S$.

As noted by Cheng and Hoza~\cite{ChengHo20}, PRGs and WPRGs can be viewed as deterministic averaging samplers (\ie, with randomness complexity and failure probability zero).  Specifically,
a PRG $G : \zo^s\rightarrow \zo^m$ for a class $\cF$ is a deterministic averaging sampler for the class $\cF$ with sample complexity $S=2^s$.  Indeed, the sampler simply outputs the set of all $S=2^s$ outputs of $G$. A WPRG can be viewed as a more general form of a nonadaptive deterministic sampler for the class $\cF$, one that is restricted to output a linear combination of the function values.

So comparing the power of PRGs vs. WPRGs is a special case of the more general problem of comparing the power of averaging samplers vs. more general nonadaptive samplers.  In this more general framing, there are some natural examples of classes $\cF$ where nonadaptive samplers can have smaller sample complexity than any averaging sampler.  Specifically, if we consider the class $\cF$ of \emph{unbounded} functions $f : \zo^m\rightarrow \bbR$ with bounded variance, \ie, $\mathrm{Var}[f(U_n)]\leq 1$, then the best sample complexity for an averaging sampler is
$\min\{\widetilde{\Theta}(1/\eps^2\delta),2^{\Theta(m)}\}$. (Essentially, Chebychev's Inequality is tight for such functions.)  However, there is a nonadaptive sampler with sample complexity $O(\log(1/\delta)/\eps^2)$, namely the \emph{median-of-averages sampler}, which outputs the median of $O(\log(1/\delta))$ averages, with each average being on $O(1/\eps^2)$ samples (see \expref{Section}{app:intro}).

This example suggests two areas of investigation.  First, can we gain further benefits in seed length by considering further generalizations of PRGs that are allowed to estimate acceptance probability with more general functions than linear combinations (or possibly even with adaptive queries)?  Some examples are the
line of work on converting hitting-set generators for circuits~\cite{AndreevClRo98,AndreevClRoTr99,BuhrmanFo99,GoldreichVaWi11} or ordered branching programs~\cite{ChengHo20} into deterministic samplers.
Second, is there a benefit in the study of samplers in restricting attention to ones that output linear combinations like WPRGs?  Perhaps these still %
retain  %
some of the useful composition properties and connections to other pseudorandom objects that are enjoyed by averaging samplers %
(see  %
\cite{Zuckerman97,vadhan2012pseudorandomness,Agrawal19}), while allowing for gains in sample and/or randomness complexity.

\subsection{Organization of the remaining sections}
In \expref{Section}{sec:prelim} we introduce arithmetic over WPRGs and the view of branching programs as %
matrix-valued  %
functions. In \expref{Section}{sec:path} we prove \expref{Theorem}{thm:path}, using a simple analysis that introduces preconditioning methods. In \expref{Section}{sec:cycle} we prove \expref{Theorem}{thm:cycle}, using more sophisticated preconditioning tools from \cite{AhmadinejadKeMuPeSiVa19} and \cite{CKPPRSV17} and the analysis of the INW PRG from~\cite{HozaPyVa21}.

\section{Preliminaries}\label{sec:prelim}
Following~\cite{ReingoldStVa13}, we will view branching programs not as boolean functions, but as matrix-valued functions $\bp:\zo^n\ra\R^{w\times w}$ where $\bp[s]_{i,j}=1$ if the branching program started at state $i$ ends at state $j$ upon reading input $s$. In all cases, we will index the rows and columns of matrices starting from \emph{zero}.

\begin{definition}
Let $B$ be a width-$w$, length-$n$ branching program with transition functions $B_1,\ldots,B_n$. For $t\in [n]$ let $\bp_t:\zo\ra \R^{w\times w}$ be defined as
\[\bp_t[s]_{a,b} = \begin{cases}   
1 & \text{if } B_t(a,s)=b\\
0 & \text{otherwise}
\end{cases}\]
For $0\leq i<j\leq n$ let $\bp_{i..j}$ be defined via matrix multiplication as 
\[\bp_{i..j}[s_{i+1}\dots s_{j}]=\bp_{i+1}[s_{i+1}]\cdots\bp_j[s_j]\]
and let $\bp=\bp_{0..n}$. Observe that $\bp_{i..j}[s_{i+1}\dots s_j]_{u,v}=1$ if and only if $\bp$ reaches state $v$ in layer $j$ when started in state $u$ in layer $i$ and reading $(s_{i+1}\dots s_j)$. Define the constant function on zero bits $\bp_{i,i}[]=\I_w$ for all $i$. This is purely for cleanliness of later derivations.
\end{definition}

We now formally %
define  %
weighted distributions.
\begin{definition}
A \textbf{weighted distribution} %
(a.\,k.\,a.  %
\textbf{pseudodistribution}) on $\zo^n$ is a jointly distributed random variable $(X,Y)$ taking values in $\zo^n\times \R$. For a (possibly %
matrix-valued) %
function $f:\zo^n\ra\R^d$, $f[(X,Y)]$ denotes the random variable $Y\cdot f(X)$ and $\overline{f}[(X,Y)]$ its expectation. For a weighted generator $(G,\rho):\zo^s\ra \zo^n\times \R$ we write $\overline{f}[(G,\rho)]$ as shorthand for $\E_{x\la U_s}[\rho(x)\cdot f(G(x))]$. We say a weighted generator $(G,\rho)$ is \textbf{$r$-bounded} if $|\rho|\leq r$.
\end{definition}
We can now define %
matrix-valued %
functions evaluated on distributions or weighted distributions. We write $U_n=U_{\zo^n}$ for convenience.
\begin{definition}
    Let $(X,Y)$ be a weighted distribution on $\zo^n$ and $\|\cdot\|$ a norm on $w\times w$ matrices and let $\mathcal{B}$ be a class of %
    length-$n$, width-$w$ %
    ordered branching programs. We say that $(X,Y)$ is \textbf{$\eps$-pseudorandom} for $\mathcal{B}$  with respect to $\|\cdot\|$  if for all $\bp \in \mathcal{B}$ we have
    $$\left\| \bbp[(X,Y)] - \bbp[U_n]]\right\|\leq \eps.$$
\end{definition}
A weighted PRG generates a weighted distribution, exactly analogous to a PRG generating a distribution.
\begin{definition}
A weighted generator $(G,\rho):\zo^s\ra\zo^n\times\R$ is \textbf{$\eps$-pseudorandom} for $\mathcal{B}$ with respect to $\|\cdot\|$ if the weighted distribution $(G(U_s),\rho(U_s))$ is $\eps$-pseudorandom for $\mathcal{B}$ with respect to $\|\cdot\|$. That is, for all $\bp\in \mathcal{B}$, $\|\bbp[(G,\rho)]-\bbp[U_n]\|\leq \eps$. We also say $(G,\rho)$ is an $\eps$-\textbf{weighted pseudorandom generator} ($\eps$-\textbf{WPRG}) for $\mathcal{B}$ with respect to $\|\cdot\|$.
\end{definition}

To use this definition, we need to select a matrix norm. We will work with several different norms on matrices $\mA \in \R^{w\times w}$. Some examples include %
\begin{itemize}
    \item $\|\mA\|_{\max} = \max_{i,j}|\mA_{i,j}|$
    \item $\|\mA\|_{1} = \max_{x\in \R^w\setminus\{0\}}\|x\mA\|_1/\|x\|_1 = \max_i \|\mA_{i,\cdot}\|$ where $\mA_{i,\cdot}$ is the $i$th row of $\mA$. We note that we define the norm in terms of left multiplication.
    \item $\|\mA\|_2 = \max_{x\in \R^w\setminus \{0\}} \|x\mA \|_2/\|x\|_2=\sigma_{\max}(\mA)$ where $\sigma_{\max}(\mA)$ is the maximum singular value of $\mA$.
\end{itemize}
Above and throughout the paper, all vectors are \emph{row} vectors.
\begin{lemma}
Suppose $(X,Y)$ is $\eps$-pseudorandom for a class $\mathcal{B}$ of branching programs with respect to $\|\cdot\|_{\max}$. Then
\begin{enumerate}
    \item $(X,Y)$ is $\eps$-pseudorandom for the class of boolean functions (according to \expref{Definition}{def:booleanWPRG}) obtained by selecting $B\in \mathcal{B}$, choosing any start vertex $v_0\in[w]$ and a single accept vertex $\vac\in [w]$.
    \item $(X,Y)$ is $w\cdot \eps$-pseudorandom for the class of boolean functions (again according to \expref{Definition}{def:booleanWPRG}) obtained by selecting $B\in \mathcal{B}$, choosing any start vertex $v_0\in[w]$ and a set $\acc\subseteq[w]$ of accept vertices.
\end{enumerate}
\end{lemma}
\begin{proof}
Unwinding the definitions we have
$\| \E[Y\cdot \bp[X]] - \E[\bp[U_n]]\|_{\max}\leq \eps,$
which implies 
$$\left| \E[Y\cdot \bp[X]_{v_0,\vac}] - \E[\bp[U_n]_{v_0,\vac}]\right|\leq \eps$$
for all states $v_0,\vac$. Since $\bp[s]_{v_0,\vac}:\zo^n\ra \zo$ is precisely the boolean function obtained by choosing start vertex $v_0$ and accept vertex $\vac$ for $B$, this shows the first bound. For the second case, enumerating over all states $v\in \acc$ and applying a union bound completes the proof. 
\end{proof}
Thus, our end goal is to construct WPRGs with respect to $\|\cdot\|_{\max}$. But at intermediate stages in our analysis we work with other norms, such as the $\ell_2$ and $\ell_1$ norms, because they are submultiplicative (\ie, $\|AB\|\leq \|A\|\cdot\|B\|$) and $\max$-norm is not.

We can define rules for ``arithmetic'' on WPRGs, which naturally translate into operations on matrix forms. Below and throughout the paper, all logs are base 2.

\begin{definition}[Sum Rule for WPRGs]\label{def:addRule}
Given WPRGs $F_a=(G_a,\rho_a),F_b=(G_b,\rho_b)$ each with seed length $s$, let $F_a+F_b$ be the WPRG with seed length $s+1$, where for $(x,y)\in \zo^s\times\zo$ we define
\[
(F_a+F_b)((x,y)) = \begin{cases}
(G_a(x),2\rho_a(x)) & y=1\\
(G_b(x),2\rho_b(x)) & y=0\\
\end{cases}\]
\end{definition}
\begin{lemma}
For WPRGs $F_a,F_b$ as defined above, for every branching program $\bp$
$$\bbp[F_a]+\bbp[F_b]=\bbp[F_a+F_b].$$
Furthermore, if $F_a$ and $F_b$ are explicit then $F_a+F_b$ is, and if $F_a$ and $F_b$ are $r$-bounded then $F_a+F_b$ is $2r$-bounded.
\end{lemma}
\begin{proof}
The explicitness and boundedness properties are immediate from the definition.
Then %
\begin{align*}
    \bbp[F_a]+\bbp[F_b] &= \E_{x\la U_s}\rho_a(x)\bp[G_a(x)]+\E_{x\la U_s}\rho_b(x)\bp[G_b(x)]\\
    &= \frac{1}{2^{s+1}}\sum_{x\in \zo^s}\left(2\rho_a(x)\bp[G_a(x)]+2\rho_b(x)\bp[G_b(x)]\right)\\
    &= \bbp[F_a+F_b]\,.  \qedhere
      %
\end{align*}
\end{proof}
We frequently take sums over large collections of WPRGs, so we state a recursive definition of the addition rule.
\begin{proposition}\label{prop:multiAdd}
Let $\{F_i:\zo^s\ra\zo^n:i\in [V]\}$ be a set of (explicit) WPRGs where given $i$, $F_i$ can be evaluated in space $O(s)$. Then
$\sum_{i=1}^VF_i$
is an explicit $2V-$bounded WPRG with seed length $s+\lceil\log(V)\rceil$.
\end{proposition}
\begin{proof}
We can recursively construct the WPRGs $F_a=\sum_{i=1}^{\lfloor V/2\rfloor}F_i$ and $F_b=\sum_{i=\lfloor V/2\rfloor+1}^{V}F_i$, which by induction have seed length at most $s+\lceil \log(V/2)\rceil$, and apply \expref{Definition}{def:addRule} (padding the seed length of $F_b$ to make it equal to the seed length of $F_a$ if necessary) to obtain seed length $s+\lceil \log(V/2)\rceil +1=s+\lceil \log(V)\rceil$.
\end{proof}

\begin{definition}[Product Rule for WPRGs]\label{def:prodrule}
Given WPRGs $F_a=(G_a,\rho_a),F_b=(G_b,\rho_b)$ each with seed length $s$ and output lengths $n,n'$, respectively, %
let $F_aF_b$ be the WPRG with seed length $2s$ and output length $n+n'$, where for $(x,y)\in \zo^s\times\zo^s$ we define
\[(F_a F_b)((x,y)) = (G_a(x)\|G_b(y),\rho_a(x)\rho_b(y))\]
where $\|$ denotes concatenation. We define the product of a WPRG $F_a$ and scalar $\lambda \in \R$ as $(\lambda F_a)(x)=(G_a(x),\lambda\cdot\rho_a(x))$.
\end{definition}
\begin{lemma}
For WPRGs $F_a,F_b$ as defined above, for every pair of branching programs $\bp,\bp'$ of lengths $n,n'$ and equal width $w$,
$$(\bbp\bbp')[F_aF_b]=\bbp[F_a]\bbp'[F_b].$$
Furthermore, if $F_a$ and $F_b$ are explicit then $F_aF_b$ is, and if $F_a$ and $F_b$ are $r$-bounded then $F_aF_b$ is $r^2$-bounded.
\end{lemma}
\begin{proof}
The explicitness and boundedness properties are immediate from the definition. 
Then %
\begin{align*}
\bbp[F_a]\bbp'[F_b] &= \E_{x\la U_s}\rho_a(x)\bp[G_a(x)]\E_{y\la U_s}\rho_b(y)\bp'[G_b(y)]\\
&= \E_{x,y\la U_s}\rho_a(x)\rho_b(y)\bp[G_a(x)]\bp'[G_b(y)]\\
&= (\bbp\bbp')[F_aF_b] \qedhere
\end{align*}
\end{proof}
We implicitly define the sum and product of WPRGs with different seed lengths, by first padding the shorter seed to equal that of the longer.
    
\section{Pseudodistributions for general branching programs}\label{sec:path}
    We first develop the WPRG for general ordered branching programs, as it outlines some of the ideas and constructions used in our main result (\expref{Theorem}{thm:cycle}) but with simpler analysis.
    \begin{restatable}{theorem}{path}\label{thm:path}
        For all $n,w \in \N$ and $\eps \in (0,1/2)$, there exists an explicit $\eps$-WPRG for the class of ordered branching programs of length $n$ and width $w$ with respect to $\|\cdot\|_{\max}$ with seed length
        \[s=O(\log(n)\log(nw)+\log(1/\eps)\log\log_{n}(1/\eps)).\]
        Moreover, the generator is $\poly(1/\eps)$ bounded.
    \end{restatable}

\subsection{A WPRG with large seed length}
In this subsection, we construct an explicit WPRG for ordered branching programs with large seed length.
\begin{restatable}{lemma}{pathRichardson}\label{lem:pathRichardson}
Given $n,w \in \N$ and $\eps \in (0,1/2)$, define $\ell=\lceil \log_{n}(1/\eps)\rceil+1$. Then there exists an explicit weighted generator $\gen_0$ such that $\gen_0$ is $\eps$-pseudorandom for the class of ordered branching programs of length $n$ and width $w$ with respect to $\|\cdot\|_{\max}$ and
\[\gen_0 = \sum_{i\in [V]}\tau_i\cdot P_{i,1}P_{i,2}\cdots P_{i,k}\]
such that %
\begin{enumerate}
    \item $V=n^{O(\ell)}=\poly(1/\eps)$
    \item $k=O(\ell)$
    \item For all $i$, $\tau_i\in \{-1,1\}$.
    \item For all $i,j$, $P_{i,j}$ is an (unweighted) PRG with seed length $s=O(\log n \cdot\log(nw))$.
    \item Given $i\in [V]$ and $j\in [k]$, $\tau_i$ and $P_{i,j}$ are evaluatable in space $O(s+\log V)$.
\end{enumerate}
\end{restatable}
Applying the sum and product rules to the output of the lemma, we obtain an $\eps$-WPRG for ordered branching programs with seed length $O(\ell\cdot \log n\cdot\log(nw)+\log V)=O(\log(nw/\eps)\log n)$, no better than the Nisan PRG. However, since $k=O(\log_{n}(1/\eps))$ we will later apply standard derandomization results to shrink the seed length of each summand.

The pseudorandom generators $P_{i,j}$ in \expref{Lemma}{lem:pathRichardson} are instantiations
of Nisan's classic generator. %
\begin{theorem}[\cite{Nisan92}]\label{thm:nisan}
    For all $n,w$ and $\eps \in (0,1/2)$, there exists an explicit $\eps$-PRG for the class of ordered branching programs of length $n$ and width $w$ with respect to $\|\cdot\|_{1}$ with seed length $s=O(\log n\cdot \log(nw/\eps))$.
\end{theorem}
We will use Nisan's generator to approximate the random walk matrix of an arbitrary branching program $B$.
\begin{definition}\label{def:pathLap}
Let $R$ be the trivial PRG on one bit. Given $n,w\in \N$, for every length-$n$, width-$w$ branching program $\bp$, define the \textbf{random walk matrix} $\mW$ of $\bp$ as the $(n+1)\times (n+1)$ block matrix
\[\mW=\begin{bmatrix}
0 & \bbp_{0..1}[R] & 0 & 0\\
0 & 0 & \ddots & 0\\
0 & 0 & 0 & \bbp_{n-1..n}[R]\\
0 & 0 & 0 & 0\\
\end{bmatrix}.\]
And the \textbf{Laplacian} $\mL$ of $\bp$ as
\[\mL=\I_{(n+1)w}-\mW=\begin{bmatrix}
\I_w & -\bbp_{0..1}[R] & 0 & 0\\
0 & \I_w & \ddots & 0\\
0 & 0 & \I_w & -\bbp_{n-1..n}[R]\\
0 & 0 & 0 & \I_w\\
\end{bmatrix}.\]
\end{definition}
The inverse of the Laplacian of $B$ holds information about random walk probabilities, which we will exploit.
\begin{remark}\label{rmk:pathLinv}
Let $\mL$ be the Laplacian of a length-$n$, width-$w$ branching program $\bp$. Then $\mL^{-1}$ has the form %
\[\mL^{-1}_{i,j}=\begin{cases}
\bbp_{i..j}[U_{j-i}] & i<j\\
\I_w & i=j\\
0 & i>j
\end{cases}.\] 
\end{remark}
Note in particular that the $(0,n)$ block of $\mL^{-1}$ equals $\bbp[U_n]$, the distribution of truly random input over the branching program. To obtain a high-quality %
estimate of $\mL^{-1}$, we first obtain a weak estimate by substituting truly random input for the output of Nisan's PRG.
\begin{proposition}\label{prop:pathErrSmall}
Let $\mL$ be the Laplacian of a length-$n$, width-$w$ branching program $\bp$. Then for every $k\in[n]$, let $\nis[k]$ be the PRG obtained from \expref{Theorem}{thm:nisan} with width $w$, length $r$ and error $\delta=1/4n^2$. Define
\[\Lt_{i,j}=\begin{cases}
\bbp_{i..j}[\nis[j-i]] & i<j\\
\I_w & i=j\\
0 & i>j
\end{cases}.\] 
Let $\Err=\I_{(n+1)w}-\Lt \mL$. Then $\lbar\Err\rbar_1\leq 1/2n$. %
\end{proposition}
%
Note that the set $\nis[r]$ of PRGs has %
no dependence on the branching program $\bp$. To prove
that  %
the error matrix has small norm, we describe its structure.
\begin{lemma}\label{lem:pathErrStruct}
Given $n,w\in \N$, let $\mL$ be the Laplacian of a length-$n$, width-$w$ branching program $\bp$, and let $\Lt$ and $\Err$ be as defined in \expref{Proposition}{prop:pathErrSmall}. Then viewing $\Err\in \R^{(n+1)w\times(n+1)w}$ as an $(n+1)\times(n+1)$ block matrix, 
we have %
\[\Err_{i,j} = \begin{cases}
\bbp_{i..j}[\nis[j-i-1]R-\nis[j-i]] & i<j\\
0 & i\geq j
\end{cases}.\]
\end{lemma}
\begin{proof}
We first consider when $i<j$. %
\begin{align*}
    \Err_{i,j} &= -(\Lt \mL)_{i,j}\\
    &= -\sum_{k=0}^{n}\Lt_{i,k}\mL_{k,j}\\
    &= -\left(\Lt_{i,j}\cdot \mL_{j,j}+\Lt_{i,j-1}\cdot \mL_{j-1,j}\right)\\
    &= -\bbp_{i..j}[\nis[j-i]]\cdot \I_w+\bbp_{i..j-1}[\nis[j-i-1]]\bbp_{j-1..j}[R]\\
    &= \bbp_{i..j}[\nis[j-i-1]R-\nis[j-i]]\,.
      %
\end{align*}
%
And for $i=j$ we have  %
\begin{align*}
    \Err_{i,i} &= \I_{(n+1)w}-(\Lt \mL)_{i,i}\\
    &= \I_{(n+1)w} - \sum_{k=0}^{n}\Lt_{i,k}\mL_{k,i}\\
    &= \I_w - \Lt_{i,i}\mL_{i,i}\\
    &= 0\,.  %
\end{align*}
The $i>j$ case is immediate, so $\Err$ has the desired form.
\end{proof}

We can then prove \expref{Proposition}{prop:pathErrSmall}.
\begin{proof}
We bound the $1$ norm of each block of $\Err$ and then take a union bound over the at most $n$ nonzero blocks in each row. Diagonal and lower triangular blocks of $\Err$ are identically zero from \expref{Lemma}{lem:pathErrStruct}. For every $i<j$ we have 
\begin{align*}
    \|\Err_{i,j}\|_1 &=\|\bbp_{i..j}[\nis[j-i-1]R-\nis[j-i]]\|_1
    %
    \tag*{\text{(\expref{Lemma}{lem:pathErrStruct})}}\\
    &\leq \|\bbp_{i..j}[\nis[j-i-1]R] -\bbp_{i..j}[U_{j-i}]\|_1+\|\bbp_{i..j}[U_{j-i}]- \bbp_{i..j}[\nis[j-i]]\|_1  \\
    &\leq  \|\bbp_{i..j-1}[\nis[j-i-1]] -\bbp_{i..j-1}[U_{j-i-1}]\|_1\cdot\|\bbp_{j-1..j}[U_1]\|_1\\ 
%
%
%
%
%
    &\quad\quad\quad + \|\bbp_{i..j}[U_{j-i}]- \bbp_{i..j}[\nis[j-i]]\|_1 %
    \tag*{\text{(Submultiplicativity)}} \\
    &\leq \frac{1}{4n^2}+\frac{1}{4n^2} = \frac{1}{2n^2}\,. %
    \tag*{
%
       \text{(\expref{Theorem}{thm:nisan})}}%
\end{align*}
%
where %
the third inequality uses that $\bbp_{j-1..j}[U_1]$ is row-stochastic hence $\|\bbp_{j-1..j}[U_1]\|\leq 1$.
Thus $\|\Err\|_1 \leq n\cdot(1/2n^2)= 1/2n$ as desired.
\end{proof}

Therefore, by replacing truly random input with a PRG of the correct length we obtain a weak approximation of $\mL^{-1}$. Following~\cite{AhmadinejadKeMuPeSiVa19} we use preconditioned Richardson iteration to boost this to a %
high-quality  %
approximation, and by describing this output in terms of a WPRG prove \expref{Lemma}{lem:pathRichardson}.

\Richardson*
\begin{proof}
We have $\I_N-\Pm\mA= (\I_N-\mP_0\mA)^{m+1}$, and then the proof follows by the submultiplicativity of $\|\cdot\|$.
\end{proof}

We now apply this to boost our weak estimate of $\mL^{-1}$ to a strong estimate.
\begin{lemma}
\label{lem:pathStrongErr}
For every $n,w\in \N$ and $\eps \in (0,1/2)$, set $\ell=\lceil \log_{n}(1/\eps)\rceil+1$. Then for every length-$n$, width-$w$ ordered branching program $\bp$ with Laplacian $\mL$, let $\Lt$ and $\Err$ be defined as in \expref{Proposition}{prop:pathErrSmall}. Then
$$\lbar\sum_{i=0}^{\ell}\Err^i\cdot \Lt-\mL^{-1}\rbar_{\max}\leq \eps/2.$$
\end{lemma}
\begin{proof}
We apply \expref{Lemma}{lem:richardson} with $\mA=\mL$, $\mP_0=\Lt$, $\|\cdot\|=\|\cdot\|_1$ and $\alpha\leq 1/2n$ (which follows from \expref{Proposition}{prop:pathErrSmall}) and obtain $\Pm=\sum_{i=0}^{\ell}\Err^i\cdot \Lt$ satisfying $\|\I-\Pm\mL\|_1\leq \eps/2n$. Finally,
\begin{align*}
    \|\Pm-\mL^{-1}\|_{\max} &\leq \|\Pm-\mL^{-1}\|_1\\ 
    &= \|(\I-\Pm \mL)\mL^{-1}\|_1\\
    &\leq \|\I-\Pm\mL\|_1\cdot \|\mL^{-1}\|_1\\
    &\leq \frac{\eps}{2n}(n+1)\,. \qedhere
      %
\end{align*}
\end{proof}

Given this error guarantee, it remains to interpret the ``output'' of Richardson iteration in an oblivious manner. Intuitively, taking powers of the $\Err$ matrix corresponds to concatenating WPRGs to create more complex WPRGs on layers. We first define an index set for products of combinations of WPRGs. The index set is equivalent to all possible divisions of the layers $\{0,\ldots,t\}$ for all $t\leq n$ into at most $\ell$ sections.
\begin{definition}\label{def:pathIndex}
Given $n,\ell \in \N$, define the index set $\vnl$ as
\[\vnl = \{(d_1,\dots, d_r) : d_i\in \Z^+, \quad 0\leq r\leq \ell,\quad \sum_{i=1}^rd_i\leq n\}.\]
For $\sigma=(d_1,\dots,d_r) \in \vnl$ we write $|\sigma|=r$. Note that this includes the empty tuple where $r=0$.
\end{definition}
Then the nonzero summands in the output of preconditioned %
Richardson  %
iteration correspond to WPRGs indexed by $\vnl$.
\begin{lemma}\label{lem:richardSummandsPath}
For all $n,w,\ell \in \N$, let $\vnl$ be defined as in \expref{Definition}{def:pathIndex} with the same $n$ and $\ell$ and for all $k\in[n]$ let $\nis[k]$ be defined as in \expref{Proposition}{prop:pathErrSmall} with the same $n,w$. For all $\sigma=(d_1,\dots,d_r) \in \vnl$, let $t=\sum_{i=1}^rd_i$ and define the WPRG (using the sum and product rules of 
%
\expref{Definitions}{def:addRule} \expref{and}{def:prodrule})  %
\[M_\sigma = \prod_{i=1}^r\left(\nis[d_i-1]R-\nis[d_i]\right)\nis[n-t].\]
Then for every length-$n$, width-$w$ branching program $\bp$ with Laplacian $\mL$, let $\Err$ and $\Lt$ be defined as in \expref{Proposition}{prop:pathErrSmall}. Then we have
\[\left(\sum_{r=0}^\ell\Err^r\cdot \Lt\right)_{0,n} = \sum_{\sigma \in \vnl}\bbp[M_\sigma].\]
\end{lemma}
\begin{proof}
Fix any $1\leq r\leq \ell$. %
Then we have  %
\begin{align*}
    \MoveEqLeft{(\Err^r\cdot\Lt)_{0,n}}\\
    &= \sum_{t_i\in \{0..n\}^{r}}\Err_{0,t_1}\left(\prod_{i=1}^{r-1}\Err_{t_i,t_{i+1}}\right)\Lt_{t_r,n}\\
    &= \sum_{0=t_0<\ldots<t_r\leq n}\left(\prod_{i=0}^{r-1}\Err_{t_i,t_{i+1}}\right)\Lt_{t_r,n} && \text{(\expref{Lemma}{lem:pathErrStruct})}\\
    &= \sum_{0=t_0<\ldots<t_r\leq n}\left(\prod_{i=0}^{r-1}\bbp_{t_i..t_{i+1}}[\nis[t_{i+1}-t_i-1]R-\nis[t_{i+1}-t_i]]\right)\bbp_{t_r..n}[\nis[n-t_r]] && \text{(\expref{Lemma}{lem:pathErrStruct})}\\
    &= \sum_{0=t_0<\ldots<t_r\leq n}\bbp\left[\left(\prod_{i=0}^{r-1}\nis[t_{i+1}-t_i-1]R-\nis[t_{i+1}-t_i]\right)\nis[n-t_r]\right] && \text{(\expref{Definition}{def:prodrule})}\\
    &= \sum_{(d_i)_{i\in [r]}:t=\sum_{i=1}^rd_i\leq n}\bbp\left[\left(\prod_{i=1}^{r}\nis[d_i-1]R-\nis[d_i]\right)\nis[n-t]\right]\\
    &= \sum_{\sigma \in \vnl,|\sigma|=r}\bbp[M_\sigma]\,.
      %
\end{align*} 
For $r=0$ we have $$(\Err^0\cdot\Lt)_{0,n}=\Lt_{0,n}=\bbp[\nis[n]]=\bbp[M_{()}].$$
Thus,
\[\left(\sum_{r=0}^\ell \Err^r\cdot\Lt\right)_{0,n}=\sum_{r=0}^\ell \sum_{\sigma \in \vnl,|\sigma|=r}\bbp[M_\sigma]= \sum_{\sigma \in \vnl}\bbp[M_\sigma].\qedhere\]
\end{proof}
Furthermore, this family of WPRGs is of the form required for \expref{Lemma}{lem:pathRichardson}.

\begin{corollary}\label{cor:MstructPath}
Given $n,w,\ell \in \N$, let $\vnl$ be defined as in \expref{Definition}{def:pathIndex} with the same $n$ and $\ell$ and let $\{M_\sigma:\sigma\in\vnl\}$ be defined as in \expref{Lemma}{lem:richardSummandsPath} with the same $n,w,\ell$. For all $\sigma=(d_1,\dots,d_r)\in \vnl$ we have
\[M_\sigma=\sum_{x\in \zo^r}\tau_{\sigma,x}\cdot P_{\sigma,x,1}\cdots P_{\sigma,x,r+1},\]
where for all $\sigma,x,i$, $\tau_{\sigma,x}\in \{-1,1\}$ and $P_{\sigma,x,i}$ is an explicit PRG with seed length $s=O(\log n\cdot\log (nw))$.
Furthermore given $\sigma=(d_1,\dots,d_r)\in \vnl,x\in \zo^r$ and $i\in [r+1]$, $\tau_{\sigma,x}$ can be computed and $P_{\sigma,x,i}$ can be evaluated in space $O(s+\log(|\vnl|\cdot r))$.
\end{corollary}
\begin{proof}
For all $\sigma\in \vnl$ and $x\in \zo^{r}$, let $\tau_{\sigma,x}=(-1)^{\sum_{i=1}^{r}x_i}$. For all $i\in[r]$, define
\[P_{\sigma,x,i} = \begin{cases}
\nis[d_{i}-1]R & x_i=0\\
\nis[d_{i}] & x_i=1
\end{cases}\]
and letting $l=\sum_{i=1}^rd_i$, define $P_{\sigma,x,r+1}=\nis[n-l]$. Then by construction
\[M_\sigma=\sum_{x\in \zo^r}\tau_{\sigma,x}\cdot P_{\sigma,x,1}\cdots P_{\sigma,x,r+1}.\] 
Given any $\sigma,x,i$ we have from \expref{Theorem}{thm:nisan} and \expref{Definition}{def:prodrule} that $P_{\sigma,x,i}$ is an explicit PRG with the desired seed length.
\end{proof}

We can now prove the main lemma of this subsection.
\pathRichardson*
\begin{proof}
Note that we can assume $\eps=1/n^{\Omega(1)}$ since otherwise the statement is satisfied by a single Nisan PRG. Let $\{M_\sigma:\sigma \in \vnl\}$ be defined as in \expref{Lemma}{lem:richardSummandsPath} with the same $n,\ell$, and let 
$$\{\tau_{\sigma,x}\cdot P_{\sigma,x,1}\cdots P_{\sigma,x,r+1}:\sigma\in \vnl,x\in \zo^{|\sigma|}\}$$ 
be the family obtained from \expref{Corollary}{cor:MstructPath} ranging over $\sigma$. Then let $[V]$ be the set of terms $(\sigma,x)$, let $k=r+1=O(\ell)$, and define
\[\gen_0=\sum_{i\in[V]}\tau_{i}\cdot P_{i,1}\cdots P_{i,k}\,.\]
     %
All explicitness and seed length conditions are satisfied from \expref{Corollary}{cor:MstructPath}, and we have $V=n^{O(\ell)} =\poly(1/\eps)$ and $\gen$ is $2V=\poly(1/\eps)$ bounded as desired.
Now fix an arbitrary length-$n$, width-$w$ branching program $\bp$ with Laplacian $\mL$ and let $\Lt$ and $\Err$ be as defined as in \expref{Proposition}{prop:pathErrSmall}. Then
\begin{align*}
    \eps/2 &\geq \lbar \left(\sum_{i=0}^\ell \Err^i\cdot \Lt\right)_{0,n}-\left(\mL^{-1}\right)_{0,n}\rbar_{\max} && \text{(\expref{Lemma}{lem:pathStrongErr})}\\ 
    &= \lbar\sum_{\sigma \in \vnl}\bbp[M_\sigma]-\left(\mL^{-1}\right)_{0,n}\rbar_{\max} && \text{(\expref{Lemma}{lem:richardSummandsPath})}\\
    &= \lbar\bbp\left[\sum_{\sigma \in \vnl}M_\sigma\right]-\bbp[U_n]\rbar_{\max} && \text{(\expref{Remark}{rmk:pathLinv})}\\
    &= \lbar\bbp\left[\sum_{i \in [V]}\tau_{i}\cdot P_{i,1}\cdots P_{i,k}\right]-\bbp[U_n]\rbar_{\max} && \text{(\expref{Corollary}{cor:MstructPath})}\\
    &= \lbar \bbp[\gen_0] - \bbp[U_n]\rbar_{\max} && \,. \qedhere
       %
\end{align*}
\end{proof}

\subsection{Shorter seed length via derandomized PRG products}
In this section, we apply standard derandomization results (namely the INW generator) to reduce the seed length of the WPRG in \expref{Lemma}{lem:pathRichardson}. Specifically, we derandomize the products of PRGs $P_{i,1}\cdots P_{i,k}$ as follows. %

\begin{lemma}
\label{lem:pathInnerTool}
Given $w\in \N$ and $\delta \in (0,1/2)$ and a tuple of PRGs $M_1,\ldots,M_{k}$ where $M_i:\zo^{s}\ra \zo^{l_i}$ and given $i$, $M_i$ is computable in space $O(s)$, there exists an explicit PRG $\widetilde{M}:\zo^{\ts}\ra\zo^l$ for $l=\sum_{i=1}^k l_i$ such that for every length-$l$, width-$w$ ordered branching program $\bp$, we have 
\[\|\bbp[\widetilde{M}]-\bbp[M_1\cdots M_k]\|_{\max}\leq \delta\] 
and $\widetilde{M}$ has seed length
\[\ts=s+O(\log k\cdot\log(kw/\delta)).\]
\end{lemma}

This result is obtained from recursive application of the derandomization lemma below. We view the Nisan PRG with seed length $s$ composed with a branching program of width $w$ as a branching program of width $w$, degree $2^s$, and length $1$. Thus, a product of these compositions can be derandomized using a PRG for branching programs of high degree. We use the formulation as stated in Lemma 11 of~\cite{ChattopadhyayLi20}, which is an application of the 
INW generator. %
\begin{lemma}[\cite{ImpagliazzoNiWi94}]
Let $G_1:\zo^s \ra \zo^{l_1}$ and $G_2:\zo^s\ra \zo^{l_2}$ be explicit PRGs. Then for every $\delta \in (0,1/2)$ there is an explicit PRG $G:\zo^{s'}\ra \zo^{l_1+l_2}$ where $s'=s+O(\log(w/\delta))$ such that for every pair of ordered branching programs $\bp$, $\bp'$ of width $w$ and lengths $l_1,l_2$, respectively, %
we have 
$$\lbar(\bbp\bbp')[G] -\bbp[G_1]\bbp'[G_2]\rbar_{\max}\leq \delta.$$
\end{lemma}

Given \expref{Lemma}{lem:pathInnerTool}, we can reduce the seed length of all products of PRGs appearing in the WPRG of \expref{Lemma}{lem:pathRichardson}.
\begin{corollary}\label{cor:pathInner}
Fix $w \in \N$ and $\gamma \in (0,1/2)$ and a family of length $n$ WPRGs $\{\tau_i\cdot P_{i,1}\cdots P_{i,k}:i\in[V]\}$ where for all $i,j$, $\tau_i\in \{-1,1\}$ and $P_{i,j}$ is a PRG with seed length $s$, and given $i$ and $j$, the coefficient $\tau_i$ can be computed and the generator $P_{i,j}$ can be evaluated in space $O(s+\log V)$. Then there is an explicit $2V$-bounded WPRG $\gen$ with seed length $s+O(\log k\cdot \log(wkV/\gamma))$ such that for every length-$n$, width-$w$ branching program $\bp,$
$$\lbar\bbp[\gen] - \bbp\left[\sum_{i\in [V]}\tau_i\cdot P_{i,1}\cdots P_{i,k}\right]\rbar_{\max}\leq \gamma.$$
\end{corollary}
\begin{proof}
For all $i\in [V]$, let $\gen_i$ be the PRG obtained from applying \expref{Lemma}{lem:pathInnerTool} to $P_{i,1}\cdots P_{i,k}$ with $\delta=\gamma/V$. Then $\gen_i$ is explicit and has seed length $s+O(\log k\cdot\log(wkV/\gamma))$.
Finally, we apply \expref{Proposition}{prop:multiAdd} and define
$$\gen=\sum_{i\in [V]} \tau_i\cdot \gen_i.$$ 
Then for every length-$n$, width-$w$ ordered branching program $\bp$,
\begin{align*}
    \lbar\bbp[\gen] - \bbp\left[\sum_{i\in [V]}\tau_i\cdot P_{i_1}\cdots P_{i_k}\right]\rbar_{\max} &\leq  \sum_{i\in [V]}\lbar \bbp[\gen_i] - \bbp\left[P_{i,1}\cdots P_{i,k}\right] \rbar_{\max}\\
    &\leq \frac{\gamma}{V}\cdot V
\end{align*}
and by \expref{Proposition}{prop:multiAdd} $\gen$ is explicit and $2V$-bounded and has seed length
%
\[\ s+O(\log(V)+\log(k)\log(wkV/\gamma))\,.   \qedhere  \]
%
\end{proof}

\subsection{Putting it all together}
We are now prepared to prove \expref{Theorem}{thm:path}.
\path*
\begin{proof}
We assume $\eps=1/n^{\Omega(1)}$ since otherwise the statement is satisfied by the Nisan PRG. Applying \expref{Lemma}{lem:pathRichardson} with the same $n,w$ and $\eps$, we obtain a generator \
\[\gen_0 = \sum_{i\in [V]}\tau_i\cdot P_{i,1}P_{i,2}\cdots P_{i,k}\] satisfying for every branching program $\bp$ of length $n$ and width $w$,
\[\lbar\bbp[\gen_0] - \bbp[U_n]\rbar_{\max}\leq \eps/2.\]

Furthermore, the family $\{\tau_i\cdot P_{i,1}P_{i,2}\cdots P_{i,k}:i\in [V]\}$ satisfies the requirements of \expref{Corollary}{cor:pathInner} with $V=\poly(1/\eps)$ and $k=O(\log_{n}(1/\eps))$ and $s=O(\log n\cdot\log nw)$. Therefore, let $\gen$ be the explicit WPRG obtained from applying \expref{Corollary}{cor:pathInner} to this family with error $\gamma=\eps/2$. Thus $\gen$ is explicit, $2V=\poly(1/\eps)$ bounded, and has seed length 
\[s=O(\log(nw)\log(n)+\log(w/\eps)\log\log_{n}(1/\eps)).\]
As the seed length is greater than $n$ otherwise, we can assume $\log(n)=\Omega(\log\log_{n}(1/\eps))$ and ignore the $\log(w)\log\log_{n}(1/\eps)$ term.

Finally, for every branching program $\bp$ of length $n$ and width $w$, we have
\begin{align*}
\lbar\bbp[\gen] - \bbp[U_n]\rbar_{\max}
    &\leq \lbar \bbp[\gen] -  \bbp[\gen_0]\rbar_{\max} + \lbar \bbp[\gen_0] - \bbp[U_n]\rbar_{\max}\\ 
    &\leq \frac{\eps}{2}+\frac{\eps}{2}
\end{align*}
where the final line comes from our choice of error in \expref{Corollary}{cor:pathInner} and \expref{Lemma}{lem:pathRichardson}.
\end{proof}

\section{Pseudodistributions for permutation branching programs}\label{sec:cycle}
In this section we prove \expref{Theorem}{thm:cycle}. To do so, we restate it in
the language of %
matrix-valued %
functions.  
\begin{restatable}{theorem}{cycleMatrix}\label{thm:cycleMatrix}
For all $n\in \N$ and $\eps \in (0,1/2)$, there exists an explicit $\eps$-WPRG for the class of permutation branching programs of length $n$ with respect to $\|\cdot\|_{\max}$ with seed length
\[O(\log(n)\sqrt{\log(n/\eps)}\sqrt{\log\log(n/\eps)}+\log(1/\eps)\log\log(n/\eps)).\]
\end{restatable}
We prove \expref{Theorem}{thm:cycleMatrix} in a similar way to \expref{Theorem}{thm:path}, with two major modifications. First, we use machinery from Ahmadinejad et al.~\cite{AhmadinejadKeMuPeSiVa19} for a more sophisticated estimate of the norm of the error matrix $\Err$. Second, we use tools from Hoza, Pyne and Vadhan~\cite{HozaPyVa21} to derandomize concatenations of WPRGs in a way that avoids dependence on the width of the branching programs being fooled.

\subsection{A WPRG with large seed length}
For the duration of the section we will assume that $n+1$ is a power of two. This is without loss of generality, as any prefix of an $\eps$-WPRG for permutation branching programs must also $\eps$-fool permutation branching programs, as the final layers could be the identity.\\ 

In the next few subsections we prove the following analogue of 
\expref{Lemma}{lem:pathRichardson}. %
\begin{restatable}{theorem}{cycleRichardson}
\label{thm:cycleRichardson}
Given $n \in \N$ and $\eps,\delta \in (0,1/2)$, let $\ell=O( \log_{1/\delta}(n/\eps))$. 
Then there exists an explicit weighted generator $\gen_0$ such that $\gen_0$ is $\eps$-pseudorandom for the class of permutation branching programs of length $n$ and arbitrary width with respect to $\|\cdot\|_{\max}$ and 
\[\gen_0 = \sum_{i\in [V]}\tau_i\cdot P_{i,1}P_{i,2}\cdots P_{i,k}\]
such that %
\begin{enumerate}
    \item $V=n^{O(\ell)}$
    \item $k=O(\ell\cdot \log n)$
    \item For all $i$, $\tau_i\in \{-1,1\}$.
    \item For all $i,j$, $P_{i,j}$ is an (unweighted) PRG with seed length $s=O(\log n \cdot\log(\log(n)/\delta))$.
    \item Given $i\in [V]$ and $j\in [k]$, $\tau_i$ and $P_{i,j}$ are evaluatable in space $O(s+\log V)$.
\end{enumerate}
\end{restatable}

Applying the sum and product rules to the output of the lemma, we obtain an $\eps$-WPRG for permutation branching programs with one accept vertex with seed length $\omega(\ell\cdot \log^2 n)$, worse than the PRG of~\cite{HozaPyVa21}. However, since $\ell\cdot \log n$ is only polylogarithmic in $n$ for our choice of $\delta$, we can apply derandomization results of~\cite{HozaPyVa21} to shrink the seed length of each summand. 

\subsection{The lift transition matrix}
Given a branching program $\bp$, we define a %
matrix-valued %
function that holds information about transitions between all pairs of layers.

\begin{definition}\label{def:hbp}
Given $n,w\in \N$ and a length-$n$, width-$w$ permutation branching program $\bp$, define the \textbf{lift transition matrix} $\hbp:\zo\ra \R^{(n+1)w\times(n+1)w}$ by
\[\hbp[s] = \begin{bmatrix}
0 & \bp_{0..1}[s] & 0 & \dots & 0\\
0 & 0 & \bp_{1..2}[s] &  & \vdots\\
\vdots & &  & \ddots & 0\\
0 & ~ &   & 0 & \bp_{n-1..n}[s]\\
\Jw & 0 & \dots & & 0\\
\end{bmatrix}\]
where we view the output as an $(n+1)\times(n+1)$ block matrix. For all $i\in\N$, define $\hbp^{(i)}:\zo^{i}\ra \R^{(n+1)w\times(n+1)w}$ as $\hbp^{(i)}[s_1\dots s_i]=\hbp[s_1]\hbp[s_2]\cdots\hbp[s_{i}]$.
\end{definition}

We start by analyzing the structure of these matrices.
\begin{proposition}\label{prop:HBPstruct}
Let $\hbp$ be the lift transition matrix for a length-$n$, width-$w$ permutation branching program. %
For  %
all $x \in \N$, $s \in \zo^x$ and $j,k\in \{0,\ldots,n\}$, let $m=j+x\mod n+1$. 
%
Then we have  %
\[(\hbp^{(x)}[s])_{j,k} = \begin{cases}
0 & k\neq m\\
\Jw & m \leq j \text{ or } x>n\\
\bp_{j..k}[s] & j < m \text{ and } x\leq n
\end{cases}\]
\end{proposition}
\begin{proof}
The only nonzero blocks of $\hbp$ have index $(i,i+1 \mod n+1)$. Thus, 
\begin{align*}
    \left(\hbp^{(x)}[s]\right)_{j,k} &= \sum_{j_1..j_{x-1}}\hbp[s_1]_{j,j_1}\hbp[s_2]_{j_1,j_2}\cdots \hbp[s_x]_{j_{x-1},k}\\
    &= \hbp[s_1]_{j,j+1}\hbp[s_2]_{j+1,j+2}\cdots \hbp[s_x]_{j+x-1,k} %
\end{align*}
    %
%
where %
all block indices are written mod $n+1$.

Clearly if $k\neq j+x\mod n+1$ this is zero. Then if $\hbp[s_i]_{n,0}=\Jw$ appears in this product we have $(\hbp^{(x)}[s])_{j,k}=\Jw$, as $\mA\cdot\Jw=\Jw\cdot\mA=\Jw$ for every doubly stochastic matrix $\mA$, and $\hbp[b]_{j,j+1}$ is doubly stochastic for all $j,b$. Otherwise the product is of the form $\bp_{j..j+1}[s_1]\bp_{j+1..j+2}[s_2]\cdots \bp_{j+x-1,j+x}[s_x]=\bp_{j..j+x}[s]$ as desired.
\end{proof}

Note that when $x>n$, $\hbp^{(x)}[s]$ has no dependence on the branching program $\bp$ or the input $s$; all nonzero blocks equal $\Jw$ and the location of those blocks depends only on $x$ and $n$. In particular, we %
have the following corollary.  %
\begin{corollary}\label{cor:longPRGtriv}
Let $\hbp$ be the lift transition matrix for a length-$n$, width-$w$ permutation branching program and let $F:\zo^s\ra\zo^x,G:\zo^{s'}\ra\zo^x$ be arbitrary PRGs. If $x>n$, then $\hbbp{x}{F}=\hbbp{x}{G}$.
\end{corollary}
\begin{proof}
By \expref{Proposition}{prop:HBPstruct} the nonzero blocks of $\hbbp{x}{F}$ and $\hbbp{x}{G}$ are located only at indices $i,i+x \mod n+1$ and are all equal to $\Jw$.
\end{proof}
If $x<n$, we can eliminate the dependence on the generator by multiplying by $(\I_{n+1}\otimes \Jw)$.
\begin{corollary}\label{cor:hbpCommutes}
Let $\hbp$ be the lift transition matrix for a length-$n$, width-$w$ permutation branching program and let $F:\zo^s\ra\zo^x,G:\zo^{s'}\ra\zo^x$ be arbitrary PRGs. Then
\[(\I_{n+1}\otimes \Jw)\hbbp{x}{F}=(\I_{n+1}\otimes \Jw)\hbbp{x}{G}.\]
\end{corollary}
\begin{proof}
By \expref{Proposition}{prop:HBPstruct} the nonzero blocks of $\hbbp{x}{F}$ and $\hbbp{x}{G}$ are located only at indices $i,i+x \mod n+1$ and these blocks are convex combinations of doubly stochastic matrices and are thus doubly stochastic, so the result follows from the fact that $\Jw\cdot\mA=\Jw$ for every doubly stochastic matrix $\mA$.
\end{proof}
These corollaries will enable long outputs to exactly cancel in the error-reduction procedure we give later.

\subsection{Approximating powers}
To analyze the distribution of PRGs over these transition matrices, we introduce the idea of a cyclic branching program, and recall a consequence of~\cite{HozaPyVa21}.
    
\begin{definition}
    A length-$n$, width-$w$ permutation branching program $B$ is \textbf{cyclic} if it has transition functions $B_1,\ldots,B_n,B_0$, where $B_0$ is a transition function from layer $n$ to layer $0$. Given a cyclic branching program $B$, define the \textbf{cyclic transition matrix} as the function $\hbp:\zo\ra \R^{(n+1)w\times(n+1)w}$ where
    \[\hbp[s] = \begin{bmatrix}
    0 & \bp_{0..1}[s] & 0 & \dots & 0\\
    0 & 0 & \bp_{1..2}[s] &  & \vdots\\
    \vdots & &  & \ddots & 0\\
    0 & ~ &   & 0 & \bp_{n-1..n}[s]\\
    \bp_{n..0}[s] & 0 & \dots & & 0\\
    \end{bmatrix}.\]
    Furthermore, for all $i$ define $\hbp^{(i)}[s_1\dots s_i]=\hbp[s_1]\dots\hbp[s_i]$.
\end{definition}
    
We then state a convenient form of the main theorem of~\cite{HozaPyVa21}. To do so, we review the notion of approximation introduced by Ahmadinejad et al.~\cite{AhmadinejadKeMuPeSiVa19}, which plays a central role in their analysis.
For a complex number $z \in \C$ we write $z^*$ to denote the complex conjugate of $z$ and $|z|$ to denote the magnitude of $z$. Note that our use of row vectors means the right vector has the conjugate transpose applied, where in~\cite{AhmadinejadKeMuPeSiVa19} this is reversed.

\begin{definition}[Unit-Circle Approximation~\cite{AhmadinejadKeMuPeSiVa19}]
    For $\mA,\mtA \in \C^{N\times N}$ and $\eps\geq 0$, we say $\mA$ is an $\eps$-\textbf{unit-circle approximation} of $\mtA$, denoted $\mA\uc_\eps \mtA$, if
    \[ \forall x,y \in \C^N, \quad \left|x(\mA-\mtA)y^*\right| \leq \frac{\eps}{2}\left(\|x\|^2+\|y\|^2-\left|x\mtA x^*+y\mtA y^*\right|\right).\]
\end{definition}
    
One nice feature of unit circle approximation is that it is preserved under convex combinations. %
\begin{proposition}\label{prop:UCavg}
    Let $\mA,\mtA ,\mX,\mtX\in \C^{N\times N}$ where $\mA\uc_\eps \mtA$ and $\mX \uc_\eps \mtX$. Then for every $\lambda\in [0,1]$, $\lambda \mA+(1-\lambda)\mX  \uc_\eps \lambda \mtA+(1-\lambda)\mtX$. %
\end{proposition}
\begin{proof}
    Fix arbitrary $\forall x,y \in \C^N$, then
    \begin{align*}
        \MoveEqLeft{\left| x(\lambda \mA+(1-\lambda)\mX -(\lambda \mtA+(1-\lambda)\mtX))y^* \right|}\\
        &\leq \lambda\left|x(\mA-\mtA)y^*\right|+(1-\lambda)\left|x(\mX -\mtX)y^*\right|\\
        &\leq \frac{\eps}{2}\left(\|x\|_2^2+\|y\|_2^2 - \lambda \left|x\mtA x^*+y\mtA y^*\right|-(1-\lambda)\left|x\mtX x^*+y\mtX y^*\right|\right)\\
        &\leq \frac{\eps}{2}\left(\|x\|_2^2+\|y\|_2^2 - \left|x(\lambda \mtA +(1-\lambda)\mtX )x^*+y(\lambda \mtA +(1-\lambda)\mtX )y^*\right|\right). \qedhere
    \end{align*}
\end{proof}
As a consequence, multiplying by subunit scalars preserves unit circle approximation.
\begin{corollary}\label{cor:UCscale}
    For $\mA,\mtA \in \C^{N\times N}$, if $\mA\uc_\eps \mtA$, for all $c \in \R$ with $c\in[0,1]$ we have $c\mA\uc_\eps c\mtA$.
\end{corollary}
\begin{proof}
    This follows from \expref{Proposition}{prop:UCavg}, taking $\lambda=c$ and $\mX=\mtX=0$.
\end{proof}
    
    We can then recall the consequence of~\cite{HozaPyVa21}. The result uses that cyclic transition matrices correspond to transition matrices of consistently-labeled graphs, and the correspondence between the INW generator and the derandomized square of Rozenman and Vadhan~\cite{RozenmanVa05}.
\begin{theorem}[Consequence of~\cite{HozaPyVa21} Theorem 1.4]\label{thm:HozaPyVa21gen}
    For all $q \in \N$ and $\eps \in (0,1/2)$, there is a family of explicit PRGs $\INW_0,\ldots,\INW_q$ such that $\INW_0$ is the trivial PRG on one bit, and for all $i\in [q]$, $\INW_i$ has seed length $s=O(q\cdot\log(1/\eps))$ and produces $2^i$ bits of output. Furthermore, for all $i \in \{0,\ldots,q-1\}$, for every cyclic permutation branching program $B$ with cyclic transition matrix $\hbp$,
    \[\hbbp{2^{i+1}}{\INW_{i+1}}\uc_{\eps} \hbbp{2^{i+1}}{\INW_i\cdot\INW_i}.\]
\end{theorem}
We can then write the lift transition matrix of a permutation branching program as the convex combination of cyclic transition matrices, and by doing so obtain a bound of the same form as \expref{Theorem}{thm:HozaPyVa21gen}.

\begin{definition}\label{def:basePRGfamily}
For all $n\in \N$ and $\delta\in (0,1/2)$, let $\INW_{0},\ldots,\INW_{q}$ be the family of PRGs obtained from applying \expref{Theorem}{thm:HozaPyVa21gen} with $q=\log(n+1)$ and error $\delta=\delta/40q^2$. These generators have seed length $s=O(q\log(q/\delta))$. Then for every length-$n$ permutation branching program $\bp$ with lift transition matrix $\hbp$, for all $i\in\{0,\dots,q\}$ define
$$\mW_i = \hbbp{2^i}{\INW_i}.$$ 
\end{definition}
We remark that since $\INW_0$ is the trivial PRG on one bit, $\mW_0$ is defined identically to what is stated in \expref{Subsection}{subsec:PBP_proof_overview}. We now show that these matrices successively approximate each other with respect to unit circle approximation.

\begin{lemma}\label{lem:WsUC}
Given $n\in\N$ and $\delta\in(0,1/2)$ and a length-$n$ permutation branching program $\bp$, for all $i\in\{0,\dots,q\}$ let $\mW_i$ be as defined as in \expref{Definition}{def:basePRGfamily} with the same $n,\delta$. Then for all $i\in [q]$,
$$\mW_{i} \uc_{\frac{\delta}{40q^2}} \mW_{i-2}^2.$$
\end{lemma}

\begin{proof}
For all $y\in [w]$ let $\hmA_y$ be the cyclic transition matrix of the length-$n$ cyclic permutation branching program $A_y$, which has transition functions $B_1,\ldots,B_n,A^y_0$ where $A^y_{0}(v,b)=(v-1+y \mod w)+1$. Observe that with %
%
our
%
choice of $A^y_0$, all walks passing from layer $n$ to layer $0$ are distributed uniformly over the relevant set of end states, where the distribution is over $y$. We first claim $\frac{1}{w}\sum_{y=1}^w\hmA_y[s]=\hbp[s]$. Fixing arbitrary $s$, $(\hmA_y[s])_{i,j}=\hbp[s]_{i,j}$ for all $y\in [w]$ and all blocks $(i,j)\neq (n,0)$. For the $(n,0)$ block, for all $u,v\in \{0,\dots,w-1\}$,
$$\left(\left(\frac{1}{w}\sum_{y=1}^w\mA_y[s]\right)_{n,0}\right)_{u,v} = \Pr_{y\in[w]}[u+y \mod w =v] = \frac{1}{w}=(\Jw)_{u,v}.$$

Furthermore by \expref{Theorem}{thm:HozaPyVa21gen}, for all $y \in [w]$ and $i\in [q]$ we have 
\begin{align*}
     \hbba{2^{i}}{\INW_{i}}
    &\uc_{\delta/40q^2}\hbba{2^{i}}{\INW_{i-1}\cdot\INW_{i-1}}%
\end{align*}
%
so 
we obtain %
\begin{align*}
    \mW_i &= \hbbp{2^i}{\INW_i}\\ 
    &= \frac{1}{w}\sum_{y=1}^w\hbba{2^{i}}{\INW_{i}}\\
    &\uc_{\delta/40q^2}\frac{1}{w}\sum_{y=1}^w \hbba{2^{i}}{\INW_{i-1}\cdot\INW_{i-1}} && \text{(\expref{Proposition}{prop:UCavg})}\\
    &= \hbbp{2^i}{\INW_{i-1}\cdot\INW_{i-1}}\\ 
    &= \mW_{i-1}^2
    &&\qedhere
\end{align*}

\end{proof}

\subsection{The cycle-lift Laplacian}
We wish to use this sequence of approximations $\mW_0,\ldots,\mW_q$ to construct a good approximation to $\mW_0^n$. To do so, we slightly modify the construction described in Section 6 of~\cite{AhmadinejadKeMuPeSiVa19}. They tensor the matrix to be approximated with a cycle of sufficient length (in our case $n+1$). To do so, they first define a series of cycles. Let $\mC_i$ be the directed cycle on $2^i$ vertices. Without loss of generality, we use the following ordering of rows and columns of $\mC$.

\begin{definition}
Let $\mC_0=[1]$ and given $\mC_i$ let $\mC_{i+1}=\begin{bmatrix}0&\I_{2^i}\\ \mC_{i}&0\end{bmatrix}$.
\end{definition}
Now let $\pi_r:\{0,\ldots,2^r-1\}\ra\{0,\ldots,2^r-1\}$ be the bijection from the usual ordering of the $2^r$-cycle to the indexing in $\mC_r$. Specifically, writing $u\in \{0,\ldots,2^{r}-1\}$ in binary as $u=u_{r-1}\ldots u_0$, we have $\pi_r(u)=u_0\ldots u_{r-1}=u_0\cdot 2^{r-1}+\pi_{r-1}(u_{r-1}\ldots u_1)$. From this, we can relate the indexing of block submatrices to that of a larger 
matrix. %
\begin{claim}\label{clm:indexStruct}
Let $\mM=\begin{bmatrix}\mA_{0,0} & \mA_{0,1}\\ \mA_{1,0} & \mA_{1,1}\end{bmatrix}$ be a $2^r\times 2^r$ matrix. Then for every $u,v\in \{0,\ldots,2^{r-1}-1\}$ and $b,c\in \{0,1\}$, we have $\mM_{\pi_r(2u+b),\pi_r(2v+c)}=(\mA_{bc})_{\pi_{r-1}(u),\pi_{r-1}(v)}$.
\end{claim}
\begin{proof}
Using the definition of the bijection $\pi_r$, we have that $\pi_r(2u+b)=b\cdot 2^{r-1}+\pi_{r-1}(u)$, which is precisely equivalent to selecting a block via $b$ and the row index of the $2^{r-1}\times 2^{r-1}$ submatrix via $\pi_{r-1}(u)$, and the same holds for the column $\pi_r(2v+c)$, so the claim follows.
\end{proof}


We now take a Laplacian of $\mC_q\otimes \mW_0$ where $q=\log(n+1)$. 
\begin{definition}\label{def:cycleDecomp}
Given $n\in \N$, $\delta\in(0,1/2)$ and a length-$n$ permutation branching program $\bp$ with lift transition matrix $\hbp$, for all $i\in \{0,\dots,q\}$ let $\mW_i$ be defined as in \expref{Definition}{def:basePRGfamily} with the same $n,\delta$. For convenience we define $N=(n+1)w$ for the remainder of the section. Then fix $c=1-1/(n+1)$ and define the \textbf{cycle-lift Laplacian} of $\bp$ as
$$\mL=\mL^{(0)}=\I_{2^qN}-\mC_q\otimes c\mW_0.$$
\end{definition}
For convenience, we write $c_i=c^{2^i}$. This definition is identical to that in~\cite{AhmadinejadKeMuPeSiVa19} except we multiply $\mW_0$ by a factor strictly less than $1$. This makes $\mL$ invertible, making the analysis cleaner and providing a bound on the singular values of $\mL$ that is independent of the width of the branching program, which enables us to obtain a seed length independent of width. We now detail the decomposition of this Laplacian using repeated Schur complements, identical to that of Section 6 of~\cite{AhmadinejadKeMuPeSiVa19}.\\

Let $H_q$ be the set of indices $\{2^{q-1}N,2^{q-1}N+1,\dots ,2^q N-1\}$, so that the cycle $\mC_q$ alternates between the indices in $H_q$ and those in $H^c_q$. The Schur complement of $\mL$ onto the indices $H_q$ is an $|H_q|\times|H_q|$ matrix which is shown in~\cite{AhmadinejadKeMuPeSiVa19} to have the following nice form. %

\[\Sc(\mL,H_q)=\I_{2^{q-1}N}-\mC_{q-1}\otimes c^2\mW_0^2\]
This is exactly the Laplacian of the cycle lift of $c^2\mW_0^2$ with $\mC_{q-1}$. We then replace $\mW_0^2$ with $\mW_1$ and again take the Schur complement with respect to $H_{q-1}$. We repeat this procedure $q$ times to obtain a decomposition of the original Laplacian $\mL$ in terms of repeated Schur complements. Formally, define $\mL^{(0)}=\mL$ and for all $i\in [q]$ define


\begin{equation}
    \mL^{(i)}=\mX_1\dots \mX_i\begin{bmatrix}\I_{(2^q-2^{q-i})N} & 0\\ 0 & \I_{2^{q-i}N}-\mC_{q-i}\otimes c_i\mW_i\end{bmatrix}\mY_i\dots \mY_1 \label{eqn:Lq}%
\end{equation}
%
where %
\[\mX_j=\begin{bmatrix} \I_{(2^q-2^{q-j+1})N} & 0 & 0\\
0 & \I_{2^{q-j}N} & 0\\
0 & -\mC_{q-j}\otimes c_{j-1}\mW_{j-1} & \I_{2^{q-j}N}
\end{bmatrix},\\
\mY_j=\begin{bmatrix} \I_{(2^q-2^{q-j+1})N} & 0 & 0\\
0 & \I_{2^{q-j}N} & -\I_{2^{q-j}}\otimes c_{j-1}\mW_{j-1}\\
0 & 0 & \I_{2^{q-j}N}\,. %
\end{bmatrix}\]
Later we will argue that all these matrices $\mL^{(i)}$%
are good approximations of $\mL$ (in an appropriate sense). Moreover, 
a computation shows that the inverse of $\Lq$ is
%
\begin{equation}\label{eqn:Lqinv}
    (\Lq)^{-1} = \mY_1^{-1}\dots \mY_q^{-1}\begin{bmatrix}\I_{(2^q-1)N} & 0\\ 0 & (\I_{N}-\mC_{0}\otimes c_q \mW_q)^{-1}\end{bmatrix}\mX_q^{-1}\dots \mX_1^{-1}%
\end{equation}
%
where %
\[\mX_j^{-1}=\begin{bmatrix} \I_{(2^q-2^{q-j+1})N} & 0 & 0\\
0 & \I_{2^{q-j}N} & 0\\
0 & \mC_{q-j}\otimes c_{j-1}\mW_{j-1} & \I_{2^{q-j}N}
\end{bmatrix}, 
\mY_j^{-1}=\begin{bmatrix} \I_{(2^q-2^{q-j+1})N} & 0 & 0\\
0 & \I_{2^{q-j}N} & \I_{2^{q-j}}\otimes c_{j-1}\mW_{j-1}\\
0 & 0 & \I_{2^{q-j}N}
\end{bmatrix}.\]
Note that it is easy to invert the final block diagonal matrix at the bottom of the recursion, as
$$(\I_{N}-\mC_{0}\otimes c_q \mW_q)^{-1}=(\I_N-c_q\I_{n+1}\otimes \Jw)^{-1}=\I_N+\left(\sum_{i=1}^\infty c_q^i\right)(\I_{n+1}\otimes \Jw)$$ 
where the first equality holds because $\mW_q = \hbbp{2^q}{\INW_q}=\I_{n+1}\otimes\Jw$ (by \expref{Corollary}{cor:hbpCommutes}) and the second because $c_q<1$. Note that this matrix has no dependence on the branching program $B$.

We now describe the form of the blocks of $(\Lq)^{-1}$ in terms of the base matrices $\mW_i$. To do so, we define $\mD_r$ for $r\in \{0,\ldots,q\}$. The rows and columns of $\mD_r$ correspond to vertices of the original cycle $\mC_{q}$. 
\begin{definition}\label{def:D}
Let $\mD_{0}=(\I_N-c_q\I_{n+1}\otimes \Jw)^{-1}\in \R^{N \times N}$. Given $\mD_{r} \in \R^{2^{r}N\times 2^{r}N}$, define $\mD_{r+1}\in\R^{2^{r+1}N\times 2^{r+1}N}$ as
\begin{align*}
    \mD_{r+1} &= \begin{bmatrix}\I_{2^{r}N}+(\I_{2^{r}}\otimes c_{q-r-1}\mW_{q-r+1})\mD_{r}(\mC_{r} \otimes c_{q-r-1}\mW_{q-r-1}) & (\I_{2^r}\otimes c_{q-r-1}\mW_{q-r+1})\mD_{r}\\ \mD_{r}(\mC_{r} \otimes c_{q-r-1}\mW_{q-r-1}) & \mD_{r}\end{bmatrix}.
\end{align*}

\end{definition}
%
\begin{lemma}\label{lem:Dzeroform}
We have $\mD_q=(\Lq)^{-1}.$
\end{lemma}
\begin{proof}
We show by induction that from $r=0$ up to $q$ that
\[(\Lq)^{-1}=\mY_1^{-1}\dots \mY_{q-r}^{-1}\begin{bmatrix}
 \I_{(2^q-2^{r})N} & 0\\ 0 & \mD_r\end{bmatrix}\mX_{q-r}^{-1}\dots\mX_1^{-1}.
\]
This is immediate for $\mD_0$. Then assuming it holds for $\mD_r$, 
we have %
\begin{align*}
    (\Lq)^{-1} &=\mY_1^{-1}\dots \mY_{q-r}^{-1}
    \begin{bmatrix}\I_{(2^q-2^{r})N} & 0\\ 0 & \mD_r\end{bmatrix}
    \mX_{q-r}^{-1}\dots\mX_1^{-1}\\
    &= \mY_1^{-1}\dots \mY_{q-r}^{-1}\begin{bmatrix}
 \I_{(2^q-2^{r+1})N} & 0 & 0\\ 0 & \I_{2^{r}N} & 0\\ 0 & 0 & \mD_r\end{bmatrix}\mX_{q-r}^{-1}\dots\mX_1^{-1}\\
 &= \mY_1^{-1}\dots \mY_{q-r}^{-1}\begin{bmatrix}
 \I_{(2^q-2^{r+1})N} & 0 & 0\\ 0 & \I_{2^{r}N} & 0\\ 0 & \mD_{r}(\mC_{r}\otimes c_{q-r-1}\mW_{q-r-1}) & \mD_r\end{bmatrix}\mX_{q-r-1}^{-1}\dots\mX_1^{-1}\\
 &= \mY_1^{-1}\dots \mY_{q-r-1}^{-1}\begin{bmatrix}
 \I_{(2^q-2^{r+1})N} & 0 \\ 0 & \mD_{r+1} \end{bmatrix}\mX_{q-r-1}^{-1}\dots\mX_1^{-1}\qedhere
\end{align*}
\end{proof}

We now show that the blocks of $\mD_q=(\Lq)^{-1}$ consist essentially of products of
%
some of the $\mW_i$.  %
\begin{lemma}\label{lem:labelstruct}
For all $r \in \{0,\ldots,q\}$ and $u,v\in \{0,\ldots,2^{r}-1\}$, let $m=2^{q-r}(v-u \mod 2^{r})$. Then
$$(\mD_r)_{\pi_{r}(u),\pi_{r}(v)}=(\I_N-c_q\I_{n+1}\otimes\Jw)^{-1}c^{m}\mW_{i_1}\dots \mW_{i_k}$$
for some indices $i_1,\dots,i_k$ such that $\sum_{j=1}^k 2^{i_j}=m$ and $k\leq 2r$. Furthermore, given $u,v$ the indices $i_1,\ldots,i_k$ are computable in space $O(\log n)$.
\end{lemma}
\begin{proof}
We prove this using induction from $r=0$ up to $q$. For the base case, there is only one block $u=v=0$ and $(\mD_{0})_{\pi_0(0),\pi_0(0)}=(\I_N-c_q\I_{n+1}\otimes\Jw)^{-1}$ as claimed.
Now assume that $\mD_{r}$ has the claimed form.

Let $u',v'\in \{0,\ldots,2^{r+1}-1\}$ be arbitrary indices in the index set of $\mD_{r+1}$. We verify $(\mD_{r+1})_{\pi_{r+1}(u'),\pi_{r+1}(v')}$ has the claimed form via casework. In all cases we take $u,v\in \{0,\ldots,2^r-1\}$.

\begin{enumerate}
    \item If $u'=2u+1$ and $v'\in 2v+1$, we have
    \begin{align*}
        (\mD_{r+1})_{\pi_{r+1}(u'),\pi_{r+1}(v')} &= (\mD_{r})_{\pi_r(u),\pi_r(v)}\\
        &= (\I_N-c_q\I_{n+1}\otimes\Jw)^{-1}c^{l}\widetilde{\mW}
    \end{align*}
    where the first line follows from \expref{Definition}{def:D} and \expref{Claim}{clm:indexStruct}. From the inductive assumption, $\widetilde{\mW}$ is a product of matrices $\mW_{i_j}$ with total length
    \begin{align*}
        l &=2^{q-r}(v-u \mod 2^{r})&\\
        &=2^{q-r-1}(2v-1-(2u-1) \mod 2^{r+1})\\
        &=2^{q-r-1}(v'-u' \mod 2^{r+1})
    \end{align*} 
    as desired.
    
    \item If $u'=2u+1$ and $v'=2v$, we have
    \begin{align*}
        (\mD_{r+1})_{\pi_{r+1}(u'),\pi_{r+1}(v')} &=  (\mD_{r}(\mC_{r} \otimes c_{q-r-1}\mW_{q-r-1}))_{\pi_r(u),\pi_r(v)}\\
        &= (\mD_{r})_{\pi_r(u),\pi_r(v-1)}c_{q-r-1}\mW_{q-r-1}\\
        &= (\I_N-c_q\I_{n+1}\otimes\Jw)^{-1}c^{l}\widetilde{\mW}c_{q-r-1}\mW_{q-r-1}
    \end{align*}
    where the second line follows from the fact that $(\mD_{r}(\mC_r\otimes \I_N))_{\pi(a),\pi(b)}=(\mD_{r})_{\pi_r(a),\pi_r(b-1)}$ for all $a,b$. From the inductive assumption, we have that the product of matrices has length 
    \begin{align*}
        l+2^{q-r-1} &=2^{q-r}(v-1-u \mod 2^{r})+2^{q-r-1}\\
        &=2^{q-r-1}(2v-(2u+1) \mod 2^{r+1})\\
        &= 2^{q-r-1}(v'-u' \mod 2^{r+1})
    \end{align*} 
    as desired.
    
    \item If $u'=2u$ and $v'=2v+1$, we have
    \begin{align*}
        (\mD_{r+1})_{\pi_{r+1}(u'),\pi_{r+1}(v')} &= ((\I_{2^r}\otimes c_{q-r-1}\mW_{q-r+1})\mD_{r})_{\pi_r(u),\pi_r(v)}\\ 
        &= c_{q-r-1}\mW_{q-r-1}(\mD_{r})_{\pi_r(u),\pi_r(v)}\\
        &=c_{q-r-1}\mW_{q-r-1}(\I_N-c_q\I_{n+1}\otimes\Jw)^{-1}c^{l}\widetilde{\mW}\\
        &= c_{q-r-1}(\I_N-c_q\I_{n+1}\otimes\Jw)^{-1}c^{l}\mW_{q-r-1}\widetilde{\mW} && (\text{\expref{Corollary}{cor:hbpCommutes}})
    \end{align*}
    From the inductive assumption, we have that the product of matrices has length
    \begin{align*}
        l+2^{q-r-1} &= 2^{q-r}(v-u \mod 2^{r})+2^{q-r-1}\\
        &=2^{q-r-1}(2v+1-2u \mod 2^{r+1})\\
        &=2^{q-r-1}(v'-u' \mod 2^{r+1})
    \end{align*} 
    as desired.
    
    \item If $u'=2u$ and $v'=2v$, we do casework based on if $u'=v'$, due to the presence of $\I_{2^rN}$ in the relevant quadrant of $\mD_{r+1}$. If $u'\neq v'$, we have
    \begin{align*}
        (\mD_{r+1})_{\pi_{r+1}(u'),\pi_{r+1}(v')}
        &= ((\I_{2^r}\otimes c_{q-r-1}\mW_{q-r+1})\mD_{r}(\mC_{r} \otimes c_{q-r-1}\mW_{q-r-1}))_{\pi_r(u),\pi_r(v)}\\
        &=   c_{q-r-1}\mW_{q-r-1}(\mD_r)_{\pi_r(u),\pi_r(v-1)}c_{q-r-1}\mW_{q-r-1}\\
        &= (\I_N-c_q\I_{n+1}\otimes\Jw)^{-1}c^{l}c_{q-r}\mW_{q-r-1}\cdot \widetilde{\mW}\cdot \mW_{q-r-1}%
    \end{align*}
    %
    where %
    the second line follows from identical reasoning to Case 2. From the inductive assumption, we have that the product of matrices has length 
    \begin{align*}
        l+2^{q-r}&=2^{q-r}(v-1-u \mod 2^{r})+2^{q-r}\\
        &=2^{q-r-1}(2v-2u \mod 2^{r+1})\\
        &=2^{q-r-1}(v'-u' \mod 2^{r+1})
    \end{align*}
    as desired.\\
    
    Finally, if $u'=v'$ we have 
     \begin{align*}
        (\mD_{r+1})_{\pi_{r+1}(u'),\pi_{r+1}(v')} &= (\I_{2^{r}N}+(\I_{2^r}\otimes c_{q-r-1}\mW_{q-r+1})\mD_{r}(\mC_{r} \otimes c_{q-r-1}\mW_{q-r-1}))_{\pi_r(u),\pi_r(v)}\\
        &= \I_N+c_{q-r-1}\mW_{q-r-1}(\mD_r)_{\pi_r(v),\pi_r(v-1)}c_{q-r-1}\mW_{q-r-1}\\
        &= \I_N+(\I_N-c_q\I_{n+1}\otimes\Jw)^{-1}c^{l}c_{q-r}\mW_{q-r-1}\cdot \widetilde{\mW}\cdot \mW_{q-r-1}
    \end{align*}
    Here we have that the product of matrices has length 
    \begin{align*}
        l+2\cdot 2^{q-r-1} &=2^{q-r}(v-1-v \mod 2^r)+ 2^{q-r}\\
        &=2^{q-r}(2^r-1)+2^{q-r}=2^q
    \end{align*} 
    so we obtain 
    \begin{align*}
        (\mD_{r+1})_{\pi_{r+1}(u'),\pi_{r+1}(v')} &= \I_N+(\I_N-c_q\I_{n+1}\otimes\Jw)^{-1}c^{l}c_{q-r}\mW_{q-r-1}\cdot \widetilde{\mW}\cdot \mW_{q-r-1}\\
        &= \I_N+(\I_N-c_q\I_{n+1}\otimes\Jw)^{-1}c_q(\I_{n+1}\otimes\Jw)\\
        &=(\I_N-c_q\I_{n+1}\otimes\Jw)^{-1}
    \end{align*}
    as desired, so $\mD_{r+1}$ has the claimed form.
\end{enumerate}

Finally, given $u',v' \in \{0,\ldots,2^{r+1}-1\}$ in the index set of $\mD_{r+1}$, the algorithm can determine which indices $u,v$ corresponding to rows/columns of $\mD_{r}$ were used to produce $(\mD_{r+1})_{u',v'}$. Given this, the machine can determine if $\mW_{q-r-1}$ were left and/or right concatenated in $(\mD_{r+1})_{u',v'}$ and recurse on $u,v$ (and storing if each concatenation occurred requires only two bits per level). Since this requires storing a constant number of bits per level and the current level $r$, we require space $O(\log n)$ to output the index set. Furthermore, note that this algorithm does not depend on the branching program $\bp$, only on $n$ and the definition of $\mC_q$.
\end{proof}
For the remainder of the section we index all relevant block matrices with respect to the conventional ordering of the cycle $\mC_q$, dropping the $\pi_q$ notation.  
\begin{corollary}\label{cor:trueCycleInverse}
We have
$$(\mL^{-1})_{0,n}=(\mL^{(0)})^{-1}_{0,n} = (\I_N-c_q\I_{n+1}\otimes\Jw)^{-1}c^n\mW_0^n.$$
\end{corollary}
\begin{proof}
  If we replace $\mW_i$ with $\mW_0^{2^i}$ in the construction above, by Equation 10 of~\cite{AhmadinejadKeMuPeSiVa19} we have that $\mL^{(i)}=\mL$ for all $i$, so from
%
\expref{Lemmas}{lem:labelstruct} \expref{and}{lem:Dzeroform} %
we obtain that $\mL^{-1}$ has the desired form.
\end{proof}
With this corollary, we obtain that an accurate estimate of $\mL^{-1}$ implies an accurate estimate of $\mL^{-1}_{0,n}=(\I_N-c_q\I_{n+1}\otimes\Jw)^{-1}c^n\mW_0^n$, and by using \expref{Corollary}{cor:hbpCommutes} we will use this to bound the distance of the WPRG to $\mW_0^n$.

\subsection{Applying Richardson iteration}
From the previous subsection, we obtain a matrix $(\Lq)^{-1}$ which we will show is a good (enough) approximation of $\mL^{-1}$ to apply preconditioned Richardson iterations (\expref{Proposition}{lem:richardson}) to obtain a very accurate estimate of $\mL^{-1}$.

\begin{proposition}\label{prop:cycleEntrywise}
Given $n\in \N$ and $\delta,\eps \in (0,1/2)$, let $\ell=\lceil \log_{1/\delta}(1/\eps)\rceil$. Then for every length-$n$ permutation branching program $\bp$, let $\mL$ be defined as in \expref{Definition}{def:cycleDecomp} and $\Lq$ as in \expeqref{Equation}{eqn:Lq}. Define $\Err = \I_{2^qN}-(\Lq)^{-1}\mL$ and define
$$\Lt=\sum_{i=0}^\ell \Err^i\cdot (\Lq)^{-1}.$$
Then $\|\Lt - \mL^{-1}\|_{\max}\leq \eps \cdot \poly(n)$.
\end{proposition}


To prove this, we follow the argument of %
Ahmadinejad %
et al.~\cite{AhmadinejadKeMuPeSiVa19} which in turn is based on Cohen et al.~\cite{CKPPRSV17}. We wish to apply preconditioned Richardson iteration (\expref{Lemma}{lem:richardson}), but to do so, we must first define an appropriate submultiplicative
%
norm. 
\begin{itemize}
        \item For a matrix $\mA \in \C^{d \times d}$ we write $\mA^*$ to denote its conjugate transpose and write $\sym[\mA]=(\mA+\mA^*)/2$ to denote its \textbf{symmetrization}.
        \item We say a Hermitian matrix $\mA$ is \textbf{positive semidefinite} (PSD) or write $\mA \succeq 0$ if $x\mA x^*\geq 0$ for all $x \in \C^d$. For two Hermitian matrices $\mA, \mtA$, we use $\mA\succeq \mtA$ to denote $\mA-\mtA\succeq 0$ and define $\preceq$ analogously.
        \item For a PSD matrix $\mH$, we define the seminorm induced by $\mH$ as $\|x\|_\mH=\sqrt{x\mH x^*}$ and the corresponding matrix seminorm as $\|\mA\|_{\mH} = \max_{x\neq 0}\|x\mA \|_{\mH}/\|x\|_{\mH}$. Note that $\|\cdot\|_{\mH}$ is submultiplicative if $\mH$ is invertible. 
    \end{itemize}
    We give a basic proposition that allows us to pass from $\mH$-norm to $2$-norm.
    \begin{proposition}
        Let $\mH\in \C^{d\times d}$ be a Hermitian positive definite matrix with minimum eigenvalue $\alpha$ and maximum eigenvalue $\beta$. Then for every $x \in \C^d$, $\alpha \|x\|_2^2\leq \|x\|^2_\mH\leq \beta\|x\|_2^2$.
    \end{proposition}
    \begin{proof}
    Let $(\mu_i)_{i\in[d]}$ be an orthonormal eigenbasis for $\mH$ and $(\lambda_i)_{i\in [d]} \in \R$ the associated eigenvalues. Then for every $x\in \C^d$,  
    \[\|x\|^2_\mH = x\left(\sum_{i=1}^d\lambda_i\mu_i^*\mu_i\right)x^*  = \sum_{i=1}^d \lambda_i |x\mu_i^*|^2 \leq \beta \sum_{i=1}^d |x\mu_i^*|^2=\beta \|x\|_2^2\]
    and the lower bound is nearly identical.
    \end{proof}
    \begin{corollary}\label{prop:condNumber}
        Let $\mH\in \C^{d\times d}$ be a Hermitian positive definite matrix with minimum eigenvalue $\alpha$ and maximum eigenvalue $\beta$. Then for every $\mA \in \C^{d\times d}$, $\sqrt{\alpha/\beta}\|\mA\|_2\leq \|\mA\|_\mH\leq \sqrt{\beta/\alpha}\|\mA\|_2$.
    \end{corollary}
    \begin{proof}
    Applying the previous proposition twice we have
    \[\|\mA\|_{\mH} = \max_{x\neq 0}\frac{\|x\mA\|_{\mH}}{\|x\|_{\mH}}\leq \sqrt{\beta/\alpha}\max_{x\neq 0}\frac{\|x\mA\|_2}{\|x\|_2} = \sqrt{\beta/\alpha}\|\mA\|_2\]
    and the lower bound is nearly identical.
    \end{proof}
    
Using these tools, we can build a positive definite matrix $\mF$ such that $\Lq$ approximates $\mL^{-1}$ with respect to $\|\cdot\|_{\mF}$.
For all $i\in \{0,\ldots,q\}$, define 
\begin{equation}
    \mS^{(i)}=\begin{bmatrix}0 & 0\\ 0 & \I_{2^{q-i}N}-\mC_{q-i}\otimes c_{i}\mW_i \end{bmatrix} \in \R^{2^qN\times 2^qN}\label{eqn:Sdecomp}
\end{equation}
where the $0$ padding is added to make the dimensions of the matrices equal, and recall
\[\sym[\mS^{(i)}] = \frac{1}{2}\left(\mS^{(i)}+(\mS^{(i)})^T\right).\]

\begin{restatable}{lemma}{goodF}\label{lem:goodF}
Let $\mS^{(i)}$ and $\mL^{(i)}$ be defined as in \expeqref{Equation}{eqn:Sdecomp} and \expeqref{Equation}{eqn:Lq}, respectively. %
Then defining
$\mF=\frac{2}{q}\sum_{i=0}^q\sym[\mS^{(i)}]$, we have %
$$\lbar \I_{2^qN}-\Lt\mL^{-1}\rbar_\mF\leq \delta.$$
\end{restatable}
We follow the proof of~\cite{AhmadinejadKeMuPeSiVa19}, except that our choice of $c$ makes all of our matrices $\mL^{(i)}$ %
strictly diagonally dominant. The only aspect of the proof which differs is the observation that Schur complements of strictly diagonally dominant matrices are strictly diagonally dominant. As such, we defer the proof to \expref{Section}{sec:goodF}.

With our choice of $c$, we obtain a bound of $1/\poly(n)$ on the minimum eigenvalue of $\mF$.
\begin{lemma}\label{lem:condF}
Let $\mS^{(i)}$ and $\mL^{(i)}$ be defined as in \expeqref{Equation}{eqn:Sdecomp} and \expeqref{Equation}{eqn:Lq}, respectively, %
and define
$\mF=\frac{2}{q}\sum_{i=0}^q\sym[\mS^{(i)}]$. The eigenvalues of $\mF$ are contained in $[1/q(n+1),3]$, and $\|\mL^{-1}\|_2\leq n+1$.
\end{lemma}
\begin{proof}
For every $v \in \R^{2^qN}$ with $\|v\|_2=1$ we have  
\[\|v\mL\|_2=\|v(\I_{2^qN}- \mC_q\otimes c\mW_0)\|_2 \in [1-c,1+c]\subset[1/(n+1),2]\]
which suffices to bound $\|\mL^{-1}\|_2$ as desired. Furthermore, we have
\[\sym[\mS^{(0)}]=\I_{2^{q}N}-\sym[\mC_{q}\otimes c_{0}\mW_0]\]
and for all $i>0$ we have
\[\sym[\mS^{(i)}]=\begin{bmatrix}0 & 0\\ 0 & \I_{2^{q-i}N}-\sym[\mC_{q-i}\otimes c_{i}\mW_i] \end{bmatrix}\]
and since $\sym[\mC_{q-i}\otimes \mW_i]$ is a stochastic matrix, all eigenvalues of $\sym[\mS^{(i)}]$ are contained in $[0,2]$ for $i>0$ and $[1-c,2]$ for $i=0$. Therefore for every $v\in \R^{2^qN}$ we have
\[\|v\mF\|_2=\left\|\frac{2}{q}\sum_{i=0}^q v\sym[\mS^{(i)}]\right\|_2 \in [1/q(n+1),3]. \qedhere\] 
\end{proof}

Now that we have this $\mF$ norm, we can use Richardson iteration to prove \expref{Proposition}{prop:cycleEntrywise}.
\begin{proof}
Applying \expref{Lemma}{lem:richardson} with $\mA=\mL,\mP_0=(\Lq)^{-1}$, $m=\ell $, norm $\|\cdot\|=\|\cdot\|_\mF$ and $\alpha=\delta$, we obtain
\[\lbar\I_{2^qN}-\Lt \mL\rbar_\mF = \left\|\I_{2^qN}-\left(\sum_{i=0}^\ell \Err^i\cdot (\Lq)^{-1}\right)\mL\right\|_\mF
    = \|\I_{2^qN}-\Pm\mL\|_\mF
    \leq \delta^\ell
    \leq \eps.\]
Thus,
\begin{align*}
    \lbar\mL^{-1}-\Lt\rbar_{\max} &\leq \lbar\mL^{-1}-\Lt\rbar_{2}\\
    &= \lbar(\mL^{-1}-\Lt)\mL \mL^{-1}\rbar_2\\
    &\leq \lbar\I_{2^qN}-\Lt\mL\rbar_2\cdot\lbar\mL^{-1}\rbar_2\\
    &\leq \left(\lbar\I_{2^qN}-\Lt\mL\rbar_\mF\cdot \sqrt{3q(n+1)}\right) \lbar\mL^{-1}\rbar_2&& \text{(\expref{Proposition}{prop:condNumber})}\\
    &\leq \eps \cdot\poly(n) && \text{(\expref{Lemma}{lem:condF})}\qedhere
\end{align*}
\end{proof}

\subsection{Interpretation as a weighted PRG}
Now that we have a polynomial in the %
$\mW$  %
approximating $\mL^{-1}$ and thus $(\mL^{-1})_{0,n}=(\I_N-c_q\I_{n+1}\otimes\Jw)^{-1}c^n\mW_0^n$, to complete the proof of \expref{Theorem}{thm:cycleRichardson} we will interpret the polynomial as $(\I_N-c_q\I_{n+1}\otimes\Jw)^{-1}c^n\hbp[\gen_0]$ for a WPRG $\gen_0$.
\begin{lemma}\label{lem:baseG}
Given $n\in\N$ and $\delta \in (0,1/2)$, let $\INW_{j}$ be as defined in \expref{Definition}{def:basePRGfamily} with the same $n$ and $\delta$. Then for all $x,y \in \{0,\ldots,n\}$ there is a product of PRGs 
$$G_{x,y}=\INW_{i_1}\cdots\INW_{i_r}$$ 
where $r\leq 2q$ and given $x,y$ the index set $i_1,\dots,i_r$ is computable in space $O(\log n)$. Furthermore, let $m=y-x \mod n+1$. Then for every length-$n$ permutation branching program $\bp$ with lift transition matrix $\hbp$, let $\Lq$ be as defined in \expeqref{Equation}{eqn:Lq}. Then,
$$ (\I_N-c_q\I_{n+1}\otimes\Jw )^{-1}c^m\hbbp{m}{G_{x,y}}=(\Lq)^{-1}_{x,y}.$$
\end{lemma}
\begin{proof}
We apply the previous lemmas on the structure of $(\Lq)^{-1}$. %
\begin{align*}
(\Lq)^{-1}_{x,y} &= (\mD_0)_{x,y} && \text{(\expref{Lemma}{lem:Dzeroform})}\\
&= (\I_N-c_q\I_{n+1}\otimes\Jw )^{-1}c^m\mW_{i_1}\cdots\mW_{i_r} && \text{(\expref{Lemma}{lem:labelstruct})}\\
&=(\I_N-c_q\I_{n+1}\otimes\Jw )^{-1}c^m \hbbp{2^{i_1}}{\INW_{i_1}}\cdots\hbbp{2^{i_r}}{\INW_{i_r}} && \text{(\expref{Definition}{def:basePRGfamily})}\\
&=(\I_N-c_q\I_{n+1}\otimes\Jw )^{-1}c^m\hbbp{m}{\INW_{i_1}\cdots\INW_{i_r}}
\end{align*}

where $r\leq 2q$ and given $x,y$ the index set $i_1,\dots,i_r$ is computable in space $O(\log n)$. Furthermore, as the indices $i_1,\dots,i_r$ are independent of $\bp$, this holds for all such branching programs $\bp$ and we conclude.
\end{proof}

We can then describe the structure of $\Err$ in terms of combinations of these base PRGs.
\begin{lemma}\label{lem:errStruct}
For every $n\in\N$ and $\delta \in (0,1/2)$ and $i,j\in \{0,\ldots,n\}$, let $m=j-i\mod n+1$. Let $\{G_{x,y}:x,y\in\{0,\dots,n\}\}$ be defined as in \expref{Lemma}{lem:baseG} with the same $n$ and $\delta$ and recall $R$ is the trivial PRG on one bit. Then for every length-$n$ permutation branching program $\bp$, let $\Err\in \R^{(n+1)N\times (n+1)N}$ be the $(n+1)\times(n+1)$ block matrix as defined in \expref{Proposition}{prop:cycleEntrywise}. Then,
\[\Err_{i,j}=\begin{cases}
0 & i=j\\
c^{m}\hbbp{m}{G_{i,j-1}R-G_{i,j}} & i\neq j.
\end{cases}\]
\end{lemma}
\begin{proof}
The proof is nearly identical to that of \expref{Lemma}{lem:pathErrStruct}. We detail the case where $i<j$. %
\begin{align*}
    \Err_{i,j} &= -\sum_{k=0}^{n}(\Lq)^{-1}_{i,k}\cdot \mL_{k,j}\\
    &= -\left[(\Lq)^{-1}_{i,j}\cdot\mL_{j,j}+(\Lq)^{-1}_{i,j-1}\cdot\mL_{j-1,j}\right] && \text{(\expref{Definition}{def:cycleDecomp})}\\
    &= -\left[(\Lq)^{-1}_{i,j}\cdot\I_{N}+(\Lq)^{-1}_{i,j-1}\cdot-c\mW_0\right]\\
    &=  (\I_N-c_q\I_{n+1}\otimes\Jw)^{-1}\left(-c^{j-i}\hbbp{m}{G_{i,j}} + c^{j-i-1}\hbbp{m-1}{G_{i,j-1}}c\mW_0\right) && \text{(\expref{Lemma}{lem:baseG})}\\
    &= (\I_N-c_q\I_{n+1}\otimes\Jw)^{-1}c^{j-i}\left(\hbbp{m}{G_{i,j-1}R}-\hbbp{m}{G_{i,j}}\right)\\
    &= c^{j-i}\hbbp{m}{G_{i,j-1}R-G_{i,j}} && \text{(\expref{Corollary}{cor:hbpCommutes}).}
\end{align*}
and the case $i>j$ is nearly identical. Then for $i=j$,
\begin{align*}
    \Err_{i,i} &= \I_N-\sum_{k=0}^{n}(\Lq)^{-1}_{i,k}\cdot \mL_{k,j}\\
    &= \I_N -\left[(\Lq)^{-1}_{i,i}\cdot\I_{N}+(\Lq)^{-1}_{i,i-1}\cdot-c\mW_0\right] && \text{(\expref{Definition}{def:cycleDecomp})}\\ 
    &= \I_N -(\I_N-c_q\I_{n+1}\otimes\Jw)^{-1}\left[\I_N-c^{n+1}\hbbp{n+1}{G_{i,i-1}R}\right]&& \text{(\expref{Lemma}{lem:baseG})}\\
    &= \I_N-(\I_N-c_q\I_{n+1}\otimes\Jw)^{-1}\left[\I_N-c_q\I_{n+1}\otimes\Jw\right] && \text{(\expref{Corollary}{cor:longPRGtriv}).}\\
    &= 0 && \qedhere
\end{align*}
\end{proof}

We require one more lemma, which ensures that terms corresponding to WPRGs of length greater than $n$ fall out of the %
error-reduction %
procedure. 
\begin{lemma}\label{lem:noLongCycles}
  For every length-$n$ permutation branching program $\bp$, let $\Err\in \R^{(n+1)N\times (n+1)N}$ be the $(n+1)\times(n+1)$ block matrix as defined in \expref{Proposition}{prop:cycleEntrywise}. For every
%
sequence $0=i_0,\ldots,i_r$ of indices %
where there exists $j$ such that $i_j\geq i_{j+1}$,
\[\prod_{j=1}^r\Err_{i_{j-1},i_j} = 0.\]
\end{lemma}
\begin{proof}
If $i_j=i_{j+1}$ for any $j$ the claim is trivially satisfied, so we assume adjacent indices are distinct.
For $j\in[r]$ let $x_j=i_{j}-i_{j-1}\mod n+1$ and $x=\sum_{j=1}^rx_j$. By assumption, $x \geq n+1$. Then 
\begin{align*}
    \prod_{j=1}^r\Err_{i_{j-1},i_j} &= \prod_{j=1}^r c^{x_j}\hbbp{x_j}{G_{i_{r-1},i_r-1}R-G_{i_{r-1},i_r-1}}
    && \text{(\expref{Lemma}{lem:errStruct})}\\
    &= c^x\hbbp{x}{\prod_{j=1}^r \left(G_{i_{r-1},i_r-1}R-G_{i_{r-1},i_r-1}\right)}\\
    &= 0%
\end{align*}
%
where %
the final line follows from writing the %
preceding %
line as a sum of $2^r$ PRGs with positive and negative %
signs %
and applying \expref{Corollary}{cor:hbpCommutes}.
\end{proof}

We are now ready to describe the entries of $\Lt$. We define the index set in an analogous way to \expref{Definition}{def:pathIndex}.
\begin{definition}\label{def:cycleIndex}
For all $n,\ell \in \N$ define the index set $\vnl$ as
\[\vnl = \{0 =\sigma_0< \sigma_1<\dots<\sigma_r\leq n : \sigma_i\in \Z^+, \quad 0\leq r\leq \ell\}.\]
For $\sigma =(\sigma_0,\sigma_1,\dots,\sigma_r)\in \vnl$ we write $|\sigma|=r$. Note that we include the empty tuple $(0)$.
\end{definition}

We now index the nonzero summands of the Richardson polynomial, equivalent to \expref{Lemma}{lem:richardSummandsPath} with the products of INW PRGs $G$ taking the place of $\nis$.
\begin{lemma}\label{lem:richardSummandsCycle}
For all $n,\ell \in \N$ and $\delta \in (0,1/2)$, let $\{G_{x,y}:x,y\in\{0,\dots,n\}\}$ be defined as in \expref{Lemma}{lem:baseG} with the same $n$ and $\delta$ and $\vnl$ as in \expref{Definition}{def:cycleIndex} with the same $n,\ell$. Then for all $\sigma=(\sigma_0,\sigma_1,\dots,\sigma_r) \in \vnl$, define the WPRG (using the sum and product rules of 
%
\expref{Definitions}{def:addRule} \expref{and}{def:prodrule}) %
\[M_\sigma = \prod_{i=0}^{r-1}(G_{\sigma_i,\sigma_{i+1}-1}R-G_{\sigma_i,\sigma_{i+1}})G_{\sigma_r,n}\]
where $M_{(0)}=G_{0,n}$. For every length-$n$ permutation branching program $\bp$, let $\Err$ and $(\Lq)^{-1}$ be defined as in \expref{Proposition}{prop:cycleEntrywise}. Then,
\[\left(\sum_{r=0}^\ell\Err^{r}\cdot(\Lq)^{-1}\right)_{0,n} = (\I_N-c_q\I_{n+1}\otimes \Jw)^{-1}c^n \sum_{\sigma \in \vnl}\hbp^{(n)}[M_\sigma].\]
\end{lemma}
\begin{proof}
Fix any $r \in [\ell]$. %
Then we have  %
\begin{align*}
\MoveEqLeft{(\Err^{r}\cdot(\Lq)^{-1})_{0,n}} \\&= \sum_{(i_j) \in \{0..n\}^r} \Err_{0,i_1}\left(\prod_{j=1}^{r-1}\Err_{i_j,i_{j+1}}\right)(\Lq)^{-1}_{i_r,n}\\
&= \sum_{\sigma\in \vnl:|\sigma|=r}\left(\prod_{i=0}^{r-1}\Err_{\sigma_{i},\sigma_{i+1}}\right)(\Lq)^{-1}_{\sigma_r,n} && \text{(\expref{Lemma}{lem:noLongCycles})}\\
&= \sum_{\sigma\in \vnl:|\sigma|=r}\prod_{i=0}^{r-1}c^{\sigma_{i+1}-\sigma_i}\hbbp{\sigma_{i+1}-\sigma_i}{G_{\sigma_i,\sigma_{i+1}-1}R-G_{\sigma_i,\sigma_{i+1}}}(\Lq)^{-1}_{\sigma_r,n} 
&&\text{(\expref{Lemma}{lem:errStruct})}\\
&=(\I_N-c_q\I_{n+1}\otimes \Jw)^{-1}c^n \sum_{\sigma \in \vnl:|\sigma|=r}\hbbp{n}{M_\sigma} && \text{(\expref{Lemma}{lem:baseG})}.
\end{align*}
Then for $r=0$ we have
$$(\Lq)^{-1}_{0,n} = (\I_N-c_q\I_{n+1}\otimes \Jw)^{-1}c^n\hbbp{n}{G_{0,n}} = (\I_N-c_q\I_{n+1}\otimes \Jw)^{-1}c^n\hbbp{n}{M_{(0)}}$$
so we conclude.
\end{proof}

We then prove an analogue of \expref{Corollary}{cor:MstructPath}.
\begin{corollary}\label{cor:MstructCycle}
Given $n,\ell \in  \N$ and $\delta \in (0,1/2)$, let $\vnl$ be defined as in \expref{Definition}{def:cycleIndex} with the same $n,\ell$ and $\{M_\sigma:\sigma\in\vnl\}$ as in \expref{Lemma}{lem:richardSummandsCycle} with the same $n,\delta$. For all $\sigma\in \vnl$, let $r=|\sigma|$. Then
\[M_\sigma=\sum_{x\in \zo^r}\tau_{\sigma,x}\cdot P_{\sigma,x,1}\cdots P_{\sigma,x,k},\]
where $k\leq 3q(r+1)$ and for all $\sigma,x,i$ we have that $\tau_{\sigma,x}\in \{-1,1\}$ and $P_{\sigma,x,i}$ is an explicit PRG with seed length $s=O(\log n\cdot\log(\log(n)/\delta))$.
Furthermore given $\sigma\in \vnl,x\in \zo^r$ and $i\in [k]$, $\tau_{\sigma,x}$ can be computed and $P_{\sigma,x,i}$ can be evaluated in space $O(s+\log|\vnl|)$.
\end{corollary}
\begin{proof}
For every $\sigma\in \vnl$ and $x\in \zo^{r}$, let $\tau_{\sigma,x}=(-1)^{\sum_{i=1}^{r}x_i}$. 
For all $i\in[r]$, define %
\[D_{\sigma,x,i} = \begin{cases}
G_{\sigma_{i-1},\sigma_i-1}R & x_i=0\\
G_{\sigma_{i-1},\sigma_i} & x_i=1
\end{cases}\]
and define $D_{\sigma,x,r+1}=G_{\sigma_{r},n}$, where by \expref{Lemma}{lem:baseG} each $D_{\sigma,x}$ is a product of at most $3q$ explicit PRGs, each with seed length $s=O(\log n\cdot\log(\log(n)/\delta))$, and given $\sigma\in \vnl,x\in \zo^r$ and $i\in[r]$ the index set of the PRGs in $D_{\sigma,x,i}$ can be computed in space $O(\log n)$. Then define %
$$\prod_{j=1}^kP_{\sigma,x,j}=\prod_{i=1}^{r+1}D_{\sigma,x,i},$$ 
where given $\sigma\in \vnl$ and $x\in \zo^r$, every index $j\in[k]$ corresponds to a base PRG in the right hand product (and this PRG is computable in the desired space bound by \expref{Lemma}{lem:baseG}), and $k\leq (r+1)3q$. Finally, by definition
\[M_\sigma=\sum_{x\in \zo^r}\tau_{\sigma,x}\cdot P_{\sigma,x,1}\cdots P_{\sigma,x,k}.\qedhere\]
\end{proof}

We next show
that  %
the family of WPRGs jointly approximate $\hbbp{n}{U_n}$, which holds the distribution of random walks from layer $0$ to layer $n$ in the branching program.
\begin{lemma}\label{lemma:innerLift}
Given $n\in \N$ and $\eps,\delta \in (0,1/2)$, let $\ell=\lceil\log_{1/\delta}(1/\eps)\rceil$ and let $\{M_\sigma:\sigma \in \vnl\}$ be defined as in \expref{Lemma}{lem:richardSummandsCycle} with the same $n,\ell$ and $\delta$. Then for every length-$n$ permutation branching program $\bp$ with lift transition matrix $\hbp$,
$$\lbar\sum_{\sigma \in \vnl} \hbbp{n}{M_\sigma}-\hbbp{n}{U_n}\rbar_{\max}\leq \eps\cdot\poly(n).$$
\end{lemma}
\begin{proof}
Let $\Lt$ and $\mL$ be defined as in \expref{Proposition}{prop:cycleEntrywise}. We obtain
\begin{align*}
    \MoveEqLeft{\eps\cdot\poly(n)}\\ 
    &\geq \lbar(\Lt-\mL^{-1})_{0,n}\rbar_{\max} && \text{(\expref{Prop.}{prop:cycleEntrywise})}\\
    &= \lbar(\Lt)_{0,n}-(\I_N-c_q\I_{n+1}\otimes \Jw)^{-1}c^n\mW_0^n\rbar_{\max} && \text{(\expref{Cor.}{cor:trueCycleInverse})}\\
    &= \lbar (\I_N-c_q\I_{n+1}\otimes \Jw)^{-1}c^n\sum_{\sigma \in \vnl} \hbbp{n}{M_\sigma} - (\I_N-c_q\I_{n+1}\otimes \Jw)^{-1}c^n\mW_0^n\rbar_{\max} && \text{(\expref{Lemma}{lem:richardSummandsCycle})}\\
    &= \lbar c^n\sum_{\sigma \in \vnl} \hbbp{n}{M_\sigma} - c^n\hbbp{n}{U_n}\rbar_{\max} && \text{(\expref{Cor.}{cor:hbpCommutes})}\\
%
%
%
    &\geq \frac{1}{4}\lbar \sum_{\sigma \in \vnl}\hbbp{n}{M_\sigma}-\hbbp{n}{U_n}\rbar_{\max} %
\end{align*}
%
where the final inequality comes from our choice of $c=1-1/(n+1)$,
%
and the third equality follows from writing 
\[\sum_{\sigma\in \vnl}M_\sigma = \sum_{\sigma\in \vnl\setminus (0)}M_\sigma+M_{(0)}\]
and noting that $(\I_{n+1}\otimes \Jw)\hbbp{n}{\sum_{\sigma\in \vnl\setminus (0)}M_\sigma}=0$ by \expref{Corollary}{cor:hbpCommutes} (as all $M_\sigma$ except $M_{(0)}$ are a sum over $2^{|\sigma|}$ PRGs with positive and negative sign by \expref{Corollary}{cor:MstructCycle}) and $(\I_{n+1}\otimes \Jw)\hbbp{n}{M_{(0)}}=(\I_{n+1}\otimes \Jw)\mW_0^n$ by \expref{Corollary}{cor:hbpCommutes}. 
\end{proof}
We remark that the exact cancellation of all terms that do not correspond to WPRG outputs of length $n$ is the motivation for our placement of $\Jw$ in the lift transition matrix, rather than some arbitrary transition function as in~\cite{HozaPyVa21}. We are now prepared to prove \expref{Theorem}{thm:cycleRichardson}.
\cycleRichardson*
\begin{proof}
Let $\{M_\sigma:\sigma \in \vnl\}$ be defined as in \expref{Lemma}{lem:richardSummandsCycle} with the same $n,\ell$ and $\delta$, and let 
$$\{\tau_{i}\cdot P_{i,1}\cdots P_{i,k}:\sigma\in \vnl,x\in \zo^{|\sigma|}\}$$ 
be the family obtained from \expref{Corollary}{cor:MstructCycle} ranging over $\sigma$. Let $[V]$ be the set of terms $(\sigma,x)$ and define
\[\gen_0=\sum_{i\in[V]}\tau_{i}\cdot P_{i,1}\cdots P_{i,k}.\]
All explicitness and seed length conditions are satisfied from \expref{Corollary}{cor:MstructCycle}, and we have $V=n^{O(\ell)}$ as desired.
Finally, fixing an arbitrary length-$n$ permutation branching program $\bp$,
\begin{align*}
    \eps\cdot\poly(n) &\geq \lbar\sum_{\sigma \in \vnl} \hbbp{n}{M_\sigma}-\hbbp{n}{U_n}\rbar_{\max} && \text{(\expref{Lemma}{lemma:innerLift})}\\
    &= \lbar\sum_{i\in[V]} \hbbp{n}{\tau_{i}\cdot P_{i,1}\cdots P_{i,k}}-\hbbp{n}{U_n}\rbar_{\max} && \text{(\expref{Corollary}{cor:MstructCycle})}\\
    &\geq \lbar\sum_{i\in[V]}  \left(\hbbp{n}{\tau_{i}\cdot P_{i,1}\cdots P_{i,k}}\right)_{0,n}-\left(\hbbp{n}{U_n}\right)_{0,n}\rbar_{\max}\\
    &= \lbar\bbp\left[\sum_{i\in[V]}\tau_{i}\cdot P_{i,1}\cdots P_{i,k}\right]-\bbp[U_n]\rbar_{\max} && \text{(\expref{Proposition}{prop:HBPstruct})}\\
    &= \lbar \bbp[\gen_0] -\bbp[U_n]\rbar_{\max}.
\end{align*}
Finally, taking $\eps\la \eps/\poly(n)$ completes the proof.
\end{proof}

\subsection{Shorter seed length via derandomized PRG products}
We now have a set of explicit WPRGs whose sum provides a %
high-quality %
approximation of an arbitrary permutation branching program of length $n$. As in \expref{Section}{sec:path}, we wish to decrease the seed length of each summand. Applying \expref{Corollary}{cor:pathInner} would give a nearly-logarithmic dependence on width. To obtain a seed length independent of width, we use the main result of ~\cite{HozaPyVa21}. 

\begin{theorem}[\cite{HozaPyVa21} Theorem 1.4]\label{thm:HPVinner}
For every $k,d\in \N$ and $\delta\in(0,1/2)$, there is an explicit PRG $H:\zo^{\sINW}\ra [d]^k$ with seed length $\sINW=O(\log(d)+\log(k)(\log(1/\delta)+\log\log(k))$ such that for every permutation branching program $\bp$ of length $k$ and degree $d$,
$$\lbar \bbp[H]-\bbp[U_n]\rbar_{\max}\leq \delta.$$
\end{theorem}

We now state the inner derandomization lemma, the analogue of \expref{Corollary}{cor:pathInner}.
\begin{lemma}\label{lem:cycleInner}
Fix $\gamma\in(0,1/2)$ and a family of length-$n$ WPRGs, %
%
 $\{\tau_i\cdot  P_{i,1}\cdots P_{i,k}:i\in [V]\}$, satisfying the following 
%
conditions.  %
\begin{itemize}
    \item For all $i$, $\tau_i\in \{-1,1\}$.
    \item For all $i_j$, $P_{i,j}$ is an explicit PRG with seed length $s$.
    \item Given $i,j$, the coefficient $\tau_i$ can be computed and the generator $P_{i,j}$ can be evaluated in space $O(s+\log V)$.
\end{itemize}
Then there exists a $2V$-bounded explicit WPRG $\gen$ with seed length $O(s+ \log k\cdot \log(V\log(k)/\gamma))$ such that for every length-$n$ permutation branching program $\bp,$
$$\lbar\bbp[\gen]- \bbp\left[\sum_{i\in[V]}\tau_i \cdot P_{i,1}\cdots P_{i,k}\right]\rbar_{\max} \leq \gamma.$$
\end{lemma}
\begin{proof}
Fix an arbitrary permutation branching program $\bp$ of length $n$ and width $w$ (and degree $2$). Let $H:\zo^{\sINW}\ra (\zo^s)^k$ be the PRG obtained from applying \expref{Theorem}{thm:HPVinner} with $k=k$, $d=2^s$ and error $\delta=\gamma/V$.

Now fix arbitrary $i \in [V]$ and consider the product $P_{i,1}\cdots P_{i,k}$. For every $j\in[k]$ let $\{l_{j-1}+1,l_{j-1}+2,\ldots,l_j\}$ be the bits of the product output by $P_{i,j}$, where $l_{0}=0$, and define $\bp'_{j-1..j}[s]=\bp_{l_{j-1}..l_j}[P_{i,j}(s)]$. Note that this defines a length $1$, degree $2^s$ permutation branching program of the same width as $\bp$.
Then $\bp'$ is a degree $2^s$, length $k$ permutation branching program. Unrolling the definition,
\begin{align*}
    \gamma/V &\geq \lbar\bbp'[U_{[2^s]^k}]-\bbp'[H]\rbar_{\max}\\
    &= \lbar \prod_{j=1}^k\E[\bp_{l_{j-1}..l_j}[P_{i,j}(U_s)]] - \E_{x\la U_{\sINW}}\left[ \prod_{j=1}^k\bp_{l_{j-1}..l_j}[P_{i,j}(H(x)_j)]\right]\rbar_{\max}\\
    &= \lbar\bbp\left[\prod_{j=1}^k P_{i,j}\right] - \bbp\left[\left(\prod_{j=1}^k P_{i,j}\right)\circ H\right]\rbar_{\max}
\end{align*}
where $H(x)_j$ is the $j$th symbol output by $H$ on seed $x$. Then for all $i$, define 
\[\gen_i=\left(\prod_{j=1}^k P_{i,j}\right)\circ H\]
which is explicit by composition of space bounded algorithms and has seed length $\sINW=O(s+\log(k)\log(V\log(k)/\gamma))$.
Finally, we apply \expref{Proposition}{prop:multiAdd} and define the explicit WPRG
$$\gen=\sum_{i\in [V]} \tau_i \cdot \gen_i$$
which by the proposition is $2V$-bounded and has seed length \[\sINW+O(\log V)=O(s+\log k\cdot\log(V\log(k)/\gamma)).\]
Finally, 
\begin{align*}
    \lbar \bbp[\gen]-\bbp\left[\sum_{i\in[V]}\tau_i\cdot P_{i,1}\cdots P_{i,k}\right]\rbar_{\max} &=
    \lbar \bbp\left[\sum_{i\in[V]}\tau_i \cdot\gen_i\right]-\bbp\left[\sum_{i\in[V]}\tau_i\cdot P_{i,1}\cdots P_{i,k}\right]\rbar_{\max}\\ 
    &\leq \sum_{i\in [V]}\lbar \bbp[\gen_i] - \bbp[P_{i,1}\cdots P_{i,k}]\rbar_{\max}\\
    &\leq \frac{\gamma}{V}\cdot V. \qedhere 
\end{align*}
\end{proof}

\subsection{Putting it all together}
We are now prepared to prove our main theorem.
\cycleMatrix*
\begin{proof}
Applying \expref{Theorem}{thm:cycleRichardson} with $n=n$, $\eps=\eps/2$ and $\delta=\delta$ to be chosen later, we obtain 
$$\ell=\lceil \log_{1/\delta}(n/\eps)\rceil = O(\log(n/\eps)/\log(1/\delta))$$ 
and a generator
\[\gen_0 = \sum_{i\in [V]}\tau_i\cdot P_{i,1}P_{i,2}\cdots P_{i,k}\]
satisfying for every length-$n$ permutation branching program $\bp$,
$$\lbar\bbp[\gen_0] - \bbp[U_n]\rbar_{\max}\leq \eps/2.$$

Furthermore, the family $\{\tau_i\cdot P_{i,1}P_{i,2}\cdots P_{i,k}:i\in [V]\}$ satisfies the requirements of \expref{Lemma}{lem:cycleInner} with $V\leq |\vnl|2^\ell=n^{O(\ell)}$ and $k\leq 5\log(n)\ell$ and $s=O(\log n\cdot\log(\log(n)/\delta))$. Therefore, let $\gen$ be the WPRG obtained from applying \expref{Lemma}{lem:cycleInner} to this family with error $\gamma=\eps/2$.\\

Then unwinding definitions, we obtain
\begin{align*}
    \lbar\bbp[\gen]-\bbp[U_n]\rbar_{\max} &\leq \lbar\bbp[\gen]- \bbp[\gen_0]\rbar_{\max}+ \left\|\bbp[\gen_0] - \bbp[U_n]\right\|_{\max}\\
    &\leq \frac{\eps}{2}+\frac{\eps}{2} = \eps
\end{align*}
where the last line follows from our choice of error in \expref{Lemma}{lem:cycleInner} and \expref{Theorem}{thm:cycleRichardson}.

It remains to optimize parameters. By \expref{Lemma}{lem:cycleInner}, $\gen$ is explicit and has seed length
\[s=O\left(\log(n)\log(\log(n)/\delta) + \log(\log(n)\ell)(\ell+\log(|\vnl|/\eps))\right).\]

Finally, we choose $\delta=2^{-\sqrt{\log(n/\eps)\log\log(n/\eps)}}$. Thus we obtain 
$$\ell=O\left(\sqrt{\log(n/\eps)}/\sqrt{\log\log(n/\eps)}\right)$$
which implies $\log|\vnl|=O(\ell\log(n))=O\left(\log(n)\sqrt{\log(n/\eps)}/\sqrt{\log\log(n/\eps)}\right)$, which implies a final seed length of
\begin{align*}
    s &= O\left(\log(n)\log\log(n) +\log(n)\sqrt{\log(n/\eps)\log\log(n/\eps)} + \log\log(n/\eps)(\ell+\log(|\vnl|/\eps))\right)\\
    &= O\left(\log(n)\sqrt{\log(n/\eps)\log\log(n/\eps)} + \frac{\log\log(n/\eps)}{\sqrt{\log\log(n/\eps)}}\sqrt{\log(n/\eps)}\log(n)+\log\log(n/\eps)\log(1/\eps)\right)\\
    &= O\left(\log(n)\sqrt{\log(n/\eps)}\sqrt{\log\log(n/\eps)}+ \log(1/\eps)\log\log(n/\eps)\right). \qedhere
\end{align*}
\end{proof}



\section{Proof of \expref{Lemma}{lem:goodF}}\label{sec:goodF}
We require two prior results. We first recall notation.
\begin{itemize}
    \item For a matrix $\mA$, we use $\mA^+$ to denote the (Moore--Penrose) %
    pseudoinverse of $\mA$.  %
    \item For a PSD matrix $\mX$, we let $\mX^{1/2}$ to denote the square root of $\mX$, which is the unique PSD matrix such that $\mX^{1/2}\mX^{1/2}=\mX$. Furthermore, let $\mX^{+/2}$ denote the pseudoinverse  %
    of the square root of $\mX$.
\end{itemize}
In our case we exclusively work with invertible matrices, so $\mA^+=\mA^{-1}$, but we
state the results in their original %
form.  %
  
%
\begin{restatable}[{\cite[Lemma 6.7]{AhmadinejadKeMuPeSiVa19}}]{lemma}{AKMsmall}
\label{lem:cycleCore}
  Let $\mS^{(0)},\dots,\mS^{(q)}$ and $\mL^{(0)},\dots,\Lq$ be defined as in 
\expeqref{Equations}{eqn:Sdecomp} \expeqref{and}{eqn:Lq}, respectively. %
%

Then for 
$$\mF = \frac{2}{q}\sum_{i=0}^q\sym[\mS^{(i)}]$$
we have the following %
two properties. %
\begin{itemize}
    \item For each $0\leq i \leq q$,
    $$\lbar \mF^{+/2}(\mL-\mL^{(i)})\mF^{+/2}\rbar_2 \leq \delta/40q .$$ %
    
    \item The final matrix $\Lq$ satisfies 
    $${\Lq}^T \mF^{+}\Lq \succeq \frac{1}{40q^2}\mF.$$
\end{itemize}
\end{restatable}
The lemma is proved in~\cite{AhmadinejadKeMuPeSiVa19} for Laplacians of cycle lifts of Eulerian graphs without the scaling factor $c$. We reproduce the proof in \expref{Section}{sec:app}, with a small modification to account for the strictly diagonally dominant
%
$\mS^{(i)}$.  %
We now recall a further lemma required to bound the norm of $\Err$.
%
\begin{lemma}[{\cite[Lemma D.4]{AhmadinejadKeMuPeSiVa19}}]\label{lem:D4}  %
Suppose we are given matrices $\mL$, $\widetilde{\mL}$ and a positive
%
semidefinite  %
matrix $\mF$ such that $\ker(\mF) \subseteq \ker(\mL) = \ker(\mL^T) = \ker(\widetilde{\mL}) = \ker(\widetilde{\mL}^T)$ and
\begin{itemize}
    \item $\|\mF^{+/2}(\mL-\widetilde{\mL})\mF^{+/2}\|_2\leq \delta,$
    \item $\widetilde{\mL}\mF^+\widetilde{\mL}\succeq \gamma \mF,$
\end{itemize}
then $\|\I_{\text{im}(\mF)}-\widetilde{\mL}^+\mL\|_\mF \leq \delta \sqrt{\gamma^{-1}}.$
\end{lemma}
%
Now we restate and prove the lemma.  %
\goodF*
\begin{proof}
  We apply \expref{Lemma}{lem:D4} with $\delta=\delta/40q$, $\mL=\mL,\widetilde{\mL}=\Lq$ and $\mF=\mF$, all of which are invertible and thus immediately satisfy the kernel
%
conditions, %
and satisfy the other %
conditions  %
by \expref{Lemma}{lem:cycleCore}, which gives
\[\|\I_{2^qN}-(\Lq)^{-1}\mL\|_\mF\leq 40q\delta/40q= \delta. \]
\end{proof}

\section{Proof of \expref{Lemma}{lem:cycleCore}}\label{sec:app}
Here we reproduce the proof of
%
\cite[Lemma 6.7]{AhmadinejadKeMuPeSiVa19} %
to verify our claim. To maintain consistency with~\cite{AhmadinejadKeMuPeSiVa19}, we return to convention and index matrices starting from $1$.
%
First we  %
recall the formal definition of a Schur complement.

\begin{definition}[Schur complement] %
  For a matrix $\mA\in \C^{N\times N}$ and $F,C\subseteq[N]$, let $\mA_{FC}$ be the submatrix induced by the rows of $F$ and columns of $C$. If $F,C$ partition $[N]$ and $\mA_{FF}$ is invertible, then we
%
define  %
the Schur complement of $\mA$ onto the set $C$ by %
%
%
%
%
\[\Sc(\mA,C) = \begin{bmatrix}
0 & 0\\
0 & \mA_{CC}-\mA_{CF}\mA_{FF}^{-1}\mA_{FC}
\end{bmatrix}.\]
Note that the standard definition is $\Sc(\mA,C)=\mA_{CC}-\mA_{CF}\mA_{FF}^{-1}\mA_{FC}$, but we define the Schur complement as dimension-preserving by padding.
%
\end{definition}

\noindent
We then recall the lemma of Cohen et al. \cite{CKKPPRS18} %
that forms the core of the 
proof. %
%
%
%
\begin{lemma}[{\cite[Lemma 2.3]{CKKPPRS18}}]  %
\label{app:sc}
  Consider a sequence
%
%
$\mS^{(0)},\ldots,\mS^{(m)}$
of $m$-by-$m$ matrices %
such that  %
%
%
%
\begin{enumerate}
    \item $\mS^{(i)}$ has nonzero indices only on the indices $[i+1,m]$,
    \item the left-right kernels of $\mS^{(i)}$ are equal, and after restricting $\mS^{(i)}$ to the indices $[i+1,m]$, the kernel of the resulting matrix equals the coordinate restriction of the vectors in the kernel of $\mS^{(0)}$, and  %
    \item the symmetrization of each $\mS^{(i)}$, denoted $\sym[\mS^{(i)}]$, is positive
%
semidefinite.  %
\end{enumerate}
Let $\mM=\mM^{(0)}=\mS^{(0)}$ and define $\mM^{(1)},\ldots, \mM^{(m)}$ iteratively by
$$\mM^{(i+1)} = \mM^{(i)} + (\mS^{(i+1)} - \Sc(\mM^{(i)},[i+1,m]))$$
If for a subsequence of indices $1 =i_0<i_1<\ldots<  i_{p_{\max}}$ associated scaling parameters $0< \theta_0,\ldots,\theta_{p_{\max}-1}<1/2$ such that $\sum_{p=0}^{p_{\max}-1}\theta_p=1$, and some global error $0< \delta <1/2$, we have for 
every $0\leq p < p_{\max}$ %
$$\|\sym[\mS^{(i_p)}]^{+/2}(\mM^{(i_p)}-\mM^{(i_{p+1})})\sym[\mS^{(i_p)}]^{+/2}\| \leq \theta_p\delta$$
then for a matrix-norm defined from the symmetrization of the $\mS^{(i_p)}$ matrices and the scaling parameters
$$\mF = \sum_{0\leq p<p_{\max}}\theta_p \sym[\mS^{(i_p)}]$$
we have the following.  %
\begin{enumerate}
    \item For each $0\leq i\leq p_{\max}$,
    $$\|\mF^{+/2}(\mM-\mM^{(i)})\mF^{+/2}\|_2 \leq \delta\,.$$
        %
    \item The final matrix $\mM^{(p_{\max})}$ satisfies
    $${\mM^{(p_{\max})}}^T\mF^+\mM^{(p_{\max})}\succeq \frac{1}{10p^2_{\max}}\mF.$$
\end{enumerate}
\end{lemma}

We require one more basic derivation and two statements on unit circle equivalence.
%
\begin{lemma}[{\cite[Lemma D.2]{AhmadinejadKeMuPeSiVa19}}] %
\label{app:diff}
Let $\mL^{(i)}$ be the $2^qN\times 2^qN$ matrices defined in  \expeqref{Equation}{eqn:Lq}. Then
$$\mL^{(i+1)}-\mL^{(i)} = \begin{bmatrix} 0 & 0\\ 0 & -\mC_{q-i-1}\otimes c_{i+1}\mW_{i+1}+\mC_{q-i-1}\otimes c_{i+1}\mW_{i}^2\end{bmatrix}.$$
\end{lemma}

This is proven without the scaling factor $c_{i+1}$ but the construction is identical.

\begin{lemma}[{\cite[Corollary 4.6]{AhmadinejadKeMuPeSiVa19}}]  %
\label{app:ucCycle}
Let $\mtW, \mW \in \C^{N\times N}$ be
%
(not necessarily symmetric)  %
matrices such that $\mtW \uc_\delta \mW$. For all $k \in \N$ let $\mC_{(k)}$ be the transition matrix for the directed cycle on $k$ vertices. Then $\mC_{(k)}\otimes\mtW \uc_\delta \mC_{(k)} \otimes \mW$.
\end{lemma}
\begin{lemma}[{\cite[Lemma 3.8]{AhmadinejadKeMuPeSiVa19}}] %
\label{app:ucLap}
Let $\mtW, \mW \in \C^{N\times N}$ be possibly asymmetric matrices. Then if $\mtW \uc_\delta \mW$ we have $\lbar\sym[\I_N-\mW]^{+/2}(\mtW-\mW)\sym[\I_N-\mW]^{+/2}\rbar_2\leq \delta$.
\end{lemma}

We are now prepared to prove the result.
\AKMsmall*
\begin{proof}
From the %
%
$\mS^{(i)}$ and the $\mL^{(i)}$  %
we build a sequence of    
%
%
%
matrices $\hS^{(j)}$ and $\hM^{(j)}$ %
%
that satisfy the conditions of \expref{Lemma}{app:sc}, and using that we derive
the statement of the lemma. For $0\leq i < q$, and $0\leq j <2^{q-i-1}N$, define
$a_i= (2^q-2^{q-i})N$, $\hS^{(a_q)}=\mS^{(q)}$, 
and
\[ 
\hS^{(a_i+j)} = \begin{cases}
\mS^{(i)} & \text{if } j=0\\
\Sc(\hM^{a_i+j-1},[a_i+j,2^qN]) & \text{otherwise}\end{cases}\]
and
$$\hM^{(h+1)} = \hM^{(h)}+(\hS^{(h+1)}-\Sc(\hM^{(h)},[h+1,2^qN])) \quad \forall 0\leq h \leq (2^q-1)N$$

Note that the %
$\hS^{(i)}$  %
satisfy all the three premises in \expref{Lemma}{app:sc}. First $\hS^{(i)}$ has non-zero entries only on the indices $[i+ 1,2^q N]$. Furthermore, as $\hS^{(i)}$ restricted to the indices $[i+ 1,2^q N]$ is a Schur complement of a diagonally dominant matrix
%
%
and diagonal dominance is preserved under Schur complements, %
%
%
%
all the restricted %
$\hS^{(i)}$   %
are invertible so the kernel %
conditions  %
are trivially satisfied and $\sym[\hS^{(i)}]$ is PSD, so all premises of the lemma hold. 
Next we show that for all %
$i$,  %
the matrix %
%
$\mL^{(i+1)}$ approximates $\mL^{(i)}$ in the norm defined by $\sym[\mS^{(i)}]$.  By
\expref{Lemma}{app:diff},
$$\mL^{(i+1)}-\mL^{(i)} = \begin{bmatrix} 0 & 0\\ 0 & -\mC_{q-i-1}\otimes c_{i+1}\mW_{i+1}+\mC_{q-i-1}\otimes c_{i+1}\mW_{i}^2\end{bmatrix}\,.$$  %
Now, given $c_{i+1}\mW_{i+1} \uc_{\delta/40q^2}c_{i+1}\mW_{i}^2$ by \expref{Lemma}{lem:WsUC} and
\expref{Corollary}{cor:UCscale}, from
%
\expref{Lemmas}{app:ucCycle} \expref{and}{app:ucLap} %
%
we obtain
\[\|\sym[\Sc(\mS^{(i)},H_i)]^{+/2}(\mL^{(i+1)}-\mL^{(i)})\sym[{\Sc(\mS^{(i)},H_i)}]^{+/2}\|_2 \leq \delta/40q^2\]
where $H_i$
%
is the array of %
indices used for the $i+1$st Schur complement. %
Then,  %
since $\sym[\Sc(\mS^{(i)},H_i)] \preceq 2\sym[\mS^{(i)}]$,
\[\|\sym[{\mS^{(i)}}]^{+/2}(\mL^{(i+1)}-\mL^{(i)})\sym[{\mS^{(i)}}]^{+/2}\|_2 \leq 2\delta/40q^2\,.\]
%
By construction, we have $\hS^{(a_i)}=\mS^{(i)}$ and $\hM^{(a_i)}=\mL^{(i)}$ for all $0\leq i\leq q$. Therefore, we have
$$\|\sym[\hS^{(a_i)}]^{+/2}(\hM^{(a_{i})}-\hM^{(a_{i+1})})\sym[\hS^{(a_i)}]^{+/2}\|_2 \leq 2\delta/40q^2 \,.$$ %
Thus by \expref{Lemma}{app:sc} for $\mF=\frac{2}{q}\sum_{i=0}^q \sym[\mS^{(i)}]$ we obtain
\[\|\mF^{+/2}(\mL-\mL^{(i)})\mF^{+/2}\|_2\leq \delta/40q \quad \forall 0\leq i\leq q\]
and 
\[{\Lq}^T\mF^+\Lq \succeq \frac{1}{40q^2}\mF\,. \qedhere\] 
%
\end{proof}

\section{Proofs of claims in the Introduction}\label{app:intro} %
We first show that weighted PRGs that consist entirely of positive coefficients are not substantially different %
from
%
unweighted PRGs.
\begin{proposition}\label{app:prop:unweight}
    Given an explicit WPRG $(G,\rho):\zo^s\ra\zo^n\times \R$
    such that $|\E_{x\la U_s}[\rho(x)]-1|\leq  \eps$, there is an explicit PRG $F$ with seed length $s'=s+O(\log(1/\eps))$ such that for any function $B:\zo^n\ra\zo$, \[|\E[B(F(U_{s'}))] - \E[\rho(U_s)\cdot B(G(U_s))] |\leq 2\eps.\]
\end{proposition}
Note that if $(G,\rho)$ is an $\eps$-WPRG with all positive coefficients for a class of functions that contains the constant function $B(x)=1$, we have $|\E_{x\la U_s}[\rho(x)\cdot 1]-1|\leq\eps$, so we can apply the above result and obtain that there exists an explicit $3\eps$-PRG for the class.

Finally, we give the proof of \expref{Proposition}{app:prop:unweight}.
\begin{proof}
    Let $\mu=\E_{x\la U_s}[\rho(x)]$ and define $\rho_T:\zo^s\ra\R$ such that $\rho_T(x)$ is equal to $\rho(x)/\mu$ rounded up or down to a multiple of $2^{-d}$ while ensuring that $\E_{x\la U_s}[\rho_T(x)]=1$. Choosing $d=\lceil\log(1/\eps)\rceil$, we obtain that for arbitrary $B:\zo^n\ra\zo$,
    \begin{align*}
        \left|\E_{x\la U_s}[\rho(x)B(G(x))] -\E_{x\la U_s}[\rho_T(x)B(G(x))]\right|  &= \left|\E_{x\la U_s}[\rho(x) - \rho_T(x)]\right|\\
        &= \left|\E_{x\la U_s}[\rho(x) -\rho(x)/\mu+\rho(x)/\mu- \rho_T(x)] \right|\\
        &\leq \left|\E_{x\la U_s}[\rho(x) -\rho(x)/\mu]\right|+\left|\E_{x\la U_s}[\rho(x)/\mu- \rho_T(x)] \right|\\
        &\leq \eps+\eps.
    \end{align*}
    So the weighted expectation of $(G,\rho_T)$ is within $2\eps$ of that of $(G,\rho)$ on every boolean function. Then let $F:\zo^{s+d}\ra\zo^n$ be an explicit PRG where for each seed $x\in \zo^s$, the output $G(x)$ appears with multiplicity $\rho_T(x)\cdot 2^d$ among the outputs $\{F(y)\}_{y\in \zo^{s+d}}$. Again fixing an arbitrary function $B$, we have
%
%
\begin{align*}
 \E_{y\la U_{s+d}}[B(F(y))] &= \frac{1}{2^s\cdot 2^d}\sum_{y\in \zo^s\times\zo^d}B(F(y))\\
 &=\frac{1}{2^s}\sum_{x\in \zo^s}\frac{\rho_T(x)\cdot 2^d}{2^d}B(G(x))
        =\E_{x\la U_s}[\rho_T(x)\cdot B(G(x))].
\end{align*}
    So we obtain the desired result. The seed length is $s'=s+d=s+O(\log(1/\eps))$.
\end{proof}

We next prove that an arbitrary explicit WPRG can be decomposed into a linear combination of two unweighted PRGs. %
\begin{proposition}
Given an explicit WPRG $(G,\rho) : \zo^s \ra \zo^n\times \R$ and $\eps>0$, there are 
explicit generators $G_+ : \zo^{s'}\rightarrow \zo^n$ and $G_- : \zo^{s'}\rightarrow \zo^n$ with seed length
$s'=O(s+\log(1/\eps))$ and coefficients $\rho_+,\rho_-\in \R^{\geq 0}$ such that for every function $B : \zo^n\rightarrow \zo$, we have %
\[\left|\E_x[\rho(x)\cdot B(G(x))] - \left( 
\rho_+\cdot \E_x\left[B(G_+(x))\right] - \rho_-\cdot \E_x\left[B(G_-(x))\right]\right)\right|\leq \eps.\]
\end{proposition}
\begin{proof}
Let $\rho_+=\E_{x\in U_s}[\rho(x)\cdot \mathbb{I}[\rho(x)\geq 0]]$ and $\rho_-=-\E_{x\in U_s}[\rho(x)\cdot \mathbb{I}[\rho(x)< 0]]$ be the average magnitude of the positive and negative coefficients,
%
respectively (over the entire set of seeds). Then $R^+=(G,\frac{\rho}{\rho_+}\mathbb{I}[\rho \geq 0])$ and $R^-=(G,-\frac{\rho}{\rho_-}\mathbb{I}[\rho< 0])$ are explicit WPRGs with all non-negative coefficients and expected weight exactly 1. We then apply \expref{Proposition}{app:prop:unweight} to $R^+$ with error $\eps=\eps/2\rho_+$ and obtain an unweighted WPRG $G_+:\zo^{s'}\ra\zo^n$ with seed length $s'=O(s+\log(1/\eps))$ such that for every $B:\zo^n\ra\zo$,
\[\left|\E_{x\la U_{s'}}[B(G_+(x))] - \E_{x\la U_{s}}\left[\frac{\rho(x)}{\rho_+}\mathbb{I}[\rho(x)\geq 0]B(G(x))\right]
\right|\leq \frac{\eps}{2\rho_+}\,.\] %
Applying an identical transformation to $R^-$ and applying the triangle inequality, we obtain for every $B:\zo^n\ra\zo$,
\begin{align*}
%
    \MoveEqLeft{\left|\E_x[\rho(x)\cdot B(G(x))] - \left(\rho_+ \cdot \E_x[B(G_+(x))]-\rho_-\cdot \E_x[B(G_-(x))]\right)\right|}\\ 
    &= \left|\left(\E_x[\rho(x)(\mathbb{I}[\rho(x)\geq 0]+\mathbb{I}[\rho(x)< 0])\cdot B(G(x))] \right)-\left( \rho_+ \cdot \E_x[B(G_+(x))]-\rho_-\cdot \E_x[B(G_-(x))] \right)\right|\\
    &\leq \left|\E_x[\rho(x)\mathbb{I}[\rho(x)\geq 0]\cdot B(G(x))]- \rho_+ \cdot \E_x[B(G_+(x))]\right| \\
%
%
    &\qquad +\left|\rho_-\cdot \E_x[B(G_-(x))]-\E_x[\rho(x)\mathbb{I}[\rho(x)< 0]\cdot B(G(x))]\right|\\
    &\leq \rho_+\left(\frac{\eps}{2\rho_+}\right)+\rho_-\left(\frac{\eps}{2\rho_-}\right)\,.
\end{align*}
%
Furthermore, note that if $(G,\rho)$ $\eps$-fools the function that accepts on all inputs, it must be the case that $\rho_+-\rho_-\in [1- \eps,1+\eps]$ and so we can take $\rho_-=1-\rho_+$ at the cost of an additive $O(\eps)$ loss in approximation.
\end{proof}
We next prove that a WPRG for a class that includes indicator test functions must have bounded coefficients. We remark that permutation branching programs %
of  %
unbounded width have this property, whereas permutation branching programs of width less
than $n$ do not. 
\begin{proposition}
  Given an explicit $\eps$-WPRG $(G,\rho):\zo^s\ra\zo^n\times \R$ that fools a
%
class $\cF$ of functions   %
that includes the indicator functions $\mathbb{I}[x=x_0]$ for every $x_0$, we have that $\rho$ is $(\eps+2^{-n})2^{s}$-bounded.
\end{proposition}
\begin{proof}
    First, we assume that $G(x)\neq G(x')$ for every $x\neq x'$, since otherwise we could combine $\rho(x),\rho(x')$ and decrease the seed length. Fix an arbitrary $x_0\in \zo^s$ and let $y=G(x_0)$, and consider the function $f(z)=\mathbb{I}[z=y]$. As $f\in \cF$, we have
    \begin{align*}
        \eps &\geq \left|\E[f(U_n)] - \E_{x\la U_s}[\rho(x)f(G(x))] \right|\\
        &= \left|2^{-n}-\frac{\rho(x_0)}{2^s}\cdot %
        f(G(x_0))\right|
    \end{align*}
    and thus $(\eps+2^{-n})2^s\geq |\rho(x)|$.
\end{proof}

Finally, we give a contrived example of a class for which WPRGs obtain exponentially shorter seed length than any PRG. We do this by requiring that the class satisfy a functional equation that WPRGs can exploit but PRGs cannot.
\begin{proposition}
Let $\cB$ be the set of all functions $B:\zo^n\ra\zo$ such that 
\begin{equation}\label{eqn:functional}
    \E[B(U_n)] = \E[B(U_S)]-\E[B(U_T)]\,,  %
\end{equation}
where $S=\{0^{n-\log n}z:z\in \zo^{\log n}\}$ and $T=\{1^{n-\log n}z:z\in \zo^{\log n}\}$. Then there exists an explicit 0-WPRG for $\cB$ with seed length $O(\log n)$. Furthermore, any $(1-1/n)$-PRG for $\cB$ has seed length $\Omega(n)$.
\end{proposition}
\begin{proof}
Let $(G,\rho):\zo^{1+\log n}\ra\zo^n\times \R$ be the WPRG where $\rho(x)=2$ if $x_1=1$ and $\rho(x)=-2$ if $x_1=0$ and $G(x)=0^{n-\log n}x_{2..\log n+1}$ if $x_1=1$ and $G(x)=1^{n-\log n}x_{2..\log n+1}$ if $x_1=0$. Then for every $B\in \cB$,
\begin{align*}
    \left|\E_{y\la U_{1+\log n}}[\rho(y)\cdot B(G(y))] - \E_{y\la U_n}[B(y)]\right| &= \left|\left(\E_{y\la U_{S}}[B(y)]-\E_{y\la U_T}[B(y)]\right) - \E_{y\la U_n}[B(y)]\right|\\
    &= 0\,.     %
\end{align*}
Now consider an arbitrary PRG $F:\zo^s\ra\zo^n$ with seed length $s\leq n-2\log(n)$. Let $\im(F)$ be the image of $F$. Now choose some $x^*\in T$ where $\Pr[F(U_s)=x^*]\leq 1/|T|$, which is always possible since $\min_{x\in T}\Pr[F(U_s)=x]\cdot |T|\leq \sum_{x\in T}\Pr[F(U_s)=x]\leq 1$. We now define the function $B$ such that $B(x)=1$ for all $x\in M=(S\cup T \cup \im(F))\setminus \{x^*\}$. Note that $|\zo^n\setminus M| \geq (1-3/n)2^n$. %
We have that $\E[B(U_S)]-\E[B(U_T)]=1-(1-1/n)=1/n$, so set $B(x)=1$ at a sufficient number of points in $\zo^n\setminus M$ so that $B$ satisfies \expeqref{Equation}{eqn:functional}, and otherwise set $B(x)=0$. By construction, $B$ satisfies \expeqref{Equation}{eqn:functional} with $\E[B(U_n)]=1/n$. However, $\E[B(F(U_s))]=1$, so $F$ is not a $(1-1/n)$-PRG for the class.
\end{proof}

We now prove facts about samplers for functions with bounded variance. Such functions form a natural class where general nonadaptive samplers obtain smaller sampler complexity than averaging samplers.

\begin{definition}
An $(\eps,\delta)$-nonadaptive sampler $G$ for a 
%
class $\cF$ of functions  %
is a randomized function that makes nonadaptive queries to $f\in \cF$ and returns as estimate $\rho$ such that, over the randomness of the algorithm,
\[\Pr[|\rho-\E[f]|\leq \eps]\geq 1- \delta.\]
We say $G$ is an \textbf{averaging sampler} if $G$ generates (possibly correlated) points $x_1,\ldots,x_t$ in the domain of $f$ and returns $\rho=\frac{1}{t}\sum_{i=1}^tf(x_i)$.
\end{definition}
We then define the model to study.
\begin{definition}
Let $\cF=\{f:\zo^m\ra\R:\var(f)\leq 1\}$ be the set of unbounded functions with variance at most $1$.
\end{definition}
We first prove that samplers exist with the claimed parameters. 
\begin{proposition}
There is an averaging sampler for $\cF$ with sample complexity $\min\{2^m,O(1/\eps^2\delta)\}$, and a nonadaptive sampler with sample complexity $\min\{2^m,O(\log(1/\delta)/\eps^2)\}$.
\end{proposition}
\begin{proof}
Fixing $\eps,\delta>0$, first note that if either minimum is $2^m$ the relevant statement is immediate by querying on all points of $\zo^m$ and returning the average, which is exactly the expectation of the function.

For the averaging sampler, let $t=1/\eps^2\delta$. Now fix an arbitrary $f\in \cF$. The sampler generates $t$ independent points $X_1,\ldots,X_t$ from $U_m$. Let $Y_1,\ldots,Y_t$ be the random variables where $Y_i=f(X_i)$, and define $Y=\frac{1}{t}\sum_{i=1}^tY_i$. Then $\E[Y]=\E[f(U_m)]$ via linearity of expectation and $\var(Y)=\frac{1}{t^2}\sum_{i=1}^t\var(Y_i)=1/t$ by independence. Applying Chebyshev's Inequality to $Y$, we %
obtain, %
for any $k>0$,
\[\Pr\left[|Y-\E[f(U_m)]| \geq \frac{k}{\sqrt{t}}\right] \leq \frac{1}{k^2}\]
and choosing $k=1/\sqrt{\delta}$ we obtain
\[\Pr\left[|Y-\E[f(U_m)]| \geq \frac{1}{\sqrt{\delta}}\eps\sqrt{\delta}\right] \leq \delta\,,%
\]
%
exactly the required condition for the averaging sampler.\\

For the nonadaptive sampler, choose $t=1/10\eps^2$ and $T=10\log(1/\delta)$. The sampler generates $X_{i,j}\la U_m$ for $i\in[t]$ and $j\in[T]$, and we %
define the random variables $Y_{i,j}=f(X_{i,j})$. Let $Y_j=\frac{1}{t}\sum_{i=1}^tY_{i,j}$ for all $j$. Then by Chebyshev, for all $j$,
\[\Pr[|Y_j-\E[f(U_m)]| \geq \eps] \leq 1/10.\]
Note that the $Y_j$ are independent and have mean $\E[f(U_m)]$ for all $j$. Then define the (independent) indicator variables $B_j=\mathbb{I}[|Y_j-\E[f(U_m)]|\geq \eps]$ and note that $\E[B_j]=\Pr[|Y_j-\E[f(U_m)]|\geq \eps]\leq 1/10$ for all $j$. If $\sum_{j=1}^TB_j<T/3$, we have that the median of the %
$Y_j$  %
is within $\eps$ of $\E[f(U_m)]$, so the sampler succeeds. We bound the probability of this
failing 
to occur by 
%
Bernstein--Hoeffding
%
%
%
%
%
where for $\rho>0$ we have %
\[\Pr\left[\sum_{j=1}^TB_j\geq T/3\right]\leq \Pr\left[\sum_{j=1}^TB_j - T/10 \geq (T/10)\rho\right] \leq \exp(-\rho^2T/10(2+\rho))\leq \delta\,,\] %
%
where  %
the final step follows from choosing $\rho=2$ and recalling $T=10\log(1/\delta)$.
\end{proof}

Finally, we prove a lower bound on sampler length for averaging samplers. %
\begin{proposition}
Let $A$ be an $(\eps,\delta)$ averaging sampler for $\cF$. Then $A$ makes $t=\min\{\widetilde{\Omega}(1/\eps^2\delta),2^{\Omega(m)}\}$ queries.
\end{proposition}
\begin{proof}
We first define our set of test functions. For $1\geq \rho\geq 4/2^{m}$ to be chosen later, let $\cF_\rho\subset \cF$ be the class of functions obtained by selecting sets $S^-\subseteq \zo^m$ and $S^+\subseteq \zo^m\setminus S^-$ both of size $\lfloor (\rho/2)\cdot 2^m\rfloor$ uniformly at random without replacement, and setting
\[f_{S^-,S^+}(x) = \begin{cases}
1/\sqrt{\rho} & x\in S^+\\
-1/\sqrt{\rho} & x \in S^-\\
0 & \text{otherwise.}\end{cases}\]
For all $f\in \cF_\rho$ we obtain $\E[f(U_m)]=0$ and $\operatorname{Var}(f) \leq  (1/\sqrt{\rho})^2\rho=1$.


If $t\geq 2^{m/3-1}$ the statement is immediately true, so we assume this is not the case. Furthermore, if $\delta \leq 2^{-m/3}$ we set $\delta=2^{-m/3}$, which does not affect the asymptotic lower bound, and likewise if $t\geq 2^{1/2\delta}$ we set $\delta=1/2\log t$ (as decreasing $\delta$ cannot affect the bound).

For every random seed $\sigma$ for the sampler, there is $k(\sigma)\in \{1,\ldots,\log t\}$ such that at least $t/\log t$ of the (not necessarily distinct) points queried by $A$ are queried with multiplicity $\{2^{k(\sigma)-1},\ldots,2^{k(\sigma)}\}$ by the pigeonhole principle. Furthermore, there is $k\in \{1,\ldots,\log t\}$ such that for at least a $1/\log t$ fraction of the seeds, $k(\sigma)=k$, again via pigeonhole. 

Let $C(\sigma)$ be the event that there are at least $t/2^k\log t$ \textit{distinct} points queried by $A$ with multiplicity at least $2^{k-1}$. We have $\Pr_\sigma[C(\sigma)]\geq 1/\log t$ as with probability $1/\log t$ over the random seed at least $t/\log t$ of the total points are queried with multiplicities in the range $\{2^{k-1},\ldots,2^k\}$, so there are at least $t/2^{k}\log t$ distinct points with the desired property.

We now choose $\rho=\delta\log^2(t)2^k/t$ (which is valid as $\rho\geq \delta/t\geq 2^{-2m/3}$ and $\rho\leq \delta\log^2 t\leq 1$) and consider the behavior of the sampler on $f$ drawn uniformly at random from $\cF_\rho$. Fix $\sigma$ for which $C(\sigma)$ occurs and let $x_1,\ldots,x_t$ be the multiset of points queried by $A$. Let $B(\sigma,f)$ be the event that $f$ takes a nonzero value on at least one point queried with multiplicity at least $2^{k-1}$. We have
\begin{align*}
    \Pr_{\sigma,f}[B(\sigma,f)] &\geq \Pr_\sigma[C(\sigma)]\cdot \left[1-(1-\rho)^{t/2^k\log^2 t}\right]\\
    &\geq \Pr_\sigma[C(\sigma)]\cdot\left[(\rho)(t/2^k\log t)-\rho^2(t/2^k\log t)^2\right]\\
    &\geq (1/\log t)(\delta \log(t)/2)\\
    &= \delta/2%
\end{align*}
%
where %
the first line follows because this event becomes strictly less likely if the nonzero points of $f$ are sampled independently with probability $\rho$, the second follows from two rounds of inclusion-exclusion, and the third from the assumption that $\rho(t/2^k\log t)\leq \delta\log t\leq 1/2$. 

Conditioned on $B(\sigma,f)$ occurring, WLOG assume $x_1$ is queried with multiplicity at least $2^{k-1}$ and $f(x_1)\neq 0$. 
Letting $Y=\sum_{i=1:x_i\neq x_1}^tf(x_i)$ be the sum of the queried points excluding $x_1$, we claim
\[\Pr_{f}[f(x_1)\geq 0,Y\geq 0|B(\sigma,f)]\geq 1/8.\]
As $B(\sigma,f)$ is independent of the sign of $f$ on all nonzero points, $\Pr_f[f(x_1)\geq 0|B(\sigma,f)]=1/2$. Since $t \leq 2^{m/3-1}$ and there are at least $2^m(\delta/t)\geq 2^{m/3}$ nonzero points for all $f\in \cF_\rho$, for every $f$ there is some point $x'$ not queried by $A$ where $f(x')\neq 0$, and
\[\Pr_f[f(x')<0|B(\sigma,f),f(x_1)>0]\geq \frac{(\rho/2)2^m+1}{\rho 2^m}\geq 1/2.\]
Then, %
since we condition on one point of each sign, $\E_{f}[Y|f(x_1)>0,f(x')<0]=0$ and the distribution is symmetric, so $\Pr_{f}[f(x_1)\geq 0,Y\geq 0|B(f,\sigma)]\geq 1/4$. Thus with probability at least $\delta/16$,
\[\left|\frac{1}{t}\sum_{i=1}^tf(x_i)-\E[f]\right|\geq \frac{2^{k-1}}{t}f(x_1)=  \frac{2^{k-1}}{t}\sqrt{\frac{t}{2^k\delta\log^2t}}\,. \]  %
By taking $\delta\la 20\delta$ we can obtain that this event occurs with probability strictly greater than $\delta$, and so by the fact that $A$ is an $(\eps,\delta)$ sampler we obtain the constraint
\[2\eps\sqrt{20\delta} \geq \frac{1}{\sqrt{(t/2^k)\log^2t}} \geq \frac{1}{\sqrt{t\log^2 t}}\]
and thus derive $t\log^2 t\geq 1/80\eps^2\delta$, which is implied by $t=\Omega(1/\eps^2\delta\log^2(1/\eps\delta))$ (with a sufficiently small constant).
\end{proof}
\section{Acknowledgements}
We thank Jack Murtagh and Sumegha Garg for insightful discussions, and 
Oded Goldreich and the CCC'21 and \emph{Theory of Computing} reviewers
for feedback that improved our presentation.


